\documentclass[12pt,a4paper,oneside,english,leqno]{smfart}

\usepackage[T1]{fontenc}
\usepackage[utf8]{inputenc}
\usepackage[english]{babel}
\usepackage{newtxtext}
\makeatletter
\let\Bbbk\relax
\makeatother
\usepackage{newtxmath}
\usepackage{amscd}
\usepackage[margin=2.5cm]{geometry}
\ifdefined\pdfmapfile\pdfmapfile{+newtx.map}\fi
\usepackage[all]{xy}
\newcommand{\rank}{\operatorname{rank}}

%-------------------------------------------
%
%
%
%
%  Num\'{e}rotation des \'{e}quations
%
%
%
%
%
%---------------------------------------------

% VOIR PLUS LOIN
                                        %met le num\'{e}ro de section dans le num\'{e}ro d'\'{e}quation

%------------------------------------------------------------------
%
%
%
%
%          Definition des macros : thm, lemme, etc.
%
%
%
%
%-----------------------------------------------------------------

% --- Subsections: numbered X.Y, titled, and listed in the table of contents,
% --- exactly as in Compositio Math. 146 (2010), 1416-1506.
\newenvironment{paragr}[1][]{\refstepcounter{subsection}%
  \addcontentsline{toc}{subsection}{\protect\numberline{\thesubsection}#1}%
  \noindent \textbf{\thesubsection . \ #1}}{\medskip}
\newenvironment{sparagr}[1][]{\refstepcounter{subsubsection} \noindent \textbf{\thesubsubsection . \ #1}}{\medskip}

% --- Results are numbered <section>.<subsection>.<n>, as in the journal.
\newcounter{theo}
\numberwithin{theo}{subsection}
\newenvironment{theoreme}{ \medskip\refstepcounter{theo}  \noindent\textbf{Theorem \thetheo}. ---\em}{\em \medskip}
\newenvironment{proposition}{\medskip\refstepcounter{theo}   \noindent\textbf{Proposition \thetheo}. ---\em}{\em\medskip}
\newenvironment{corollaire}{\medskip\refstepcounter{theo}  \noindent\textbf{Corollary \thetheo}. ---\em}{\em\medskip}
\newenvironment{remarque}{\medskip \noindent \textbf{Remark}. --- }{}
\newenvironment{remarques}{\medskip \noindent \textbf{Remarks}. --- }{}
\newenvironment{lemme}{\medskip\refstepcounter{theo}   \noindent\textbf{Lemma \thetheo}. ---\em}{\em\medskip}
\newenvironment{conjecture}{ \medskip\refstepcounter{theo}  \noindent\textbf{Conjecture \thetheo}. ---\em}{\em \medskip}

%\newenvironment{conjecture}{\medskip\refstespcounter{theo}  \noindent\textbf{Conjecture \thetheo}. ---\em}{\em\medskip}
\newenvironment{definition}{\medskip\refstepcounter{theo}  \noindent\textbf{Definition \thetheo}. ---}{\medskip}
\newenvironment{hypothese}{\medskip\refstepcounter{theo}  \noindent\textbf{Hypothesis \thetheo}. ---}{\medskip}

\newenvironment{preuve}[1][]{\noindent \textbf{Proof.} #1 --- }{\hfill
  \ensuremath{\square} \medskip}

%---------------------------------------------------
%
%
%
%  METTRE LE NUMERO DE SECTION DANS LES NUMEROS D EQUATIONS
%
%
%
%
%
%
%------------------------------------------------------

%\numberwithin{equation}{section}

%---------------------------------------------
%
%
%
%     LES OPERATEURS MATHEMATIQUES
%
%
%
%
%---------------------------------------------

\DeclareMathOperator{\Lie}{Lie}
\DeclareMathOperator{\nd}{nd}
\DeclareMathOperator{\nr}{nr}
\DeclareMathOperator{\reg}{reg}
\DeclareMathOperator{\freg}{freg}
\DeclareMathOperator{\vol}{vol}
\DeclareMathOperator{\rang}{rang}
\DeclareMathOperator{\Ad}{Ad}
\DeclareMathOperator{\AD}{AD}
\DeclareMathOperator{\End}{End}
\DeclareMathOperator{\ad}{ad}
\DeclareMathOperator{\ab}{ab}
\DeclareMathOperator{\supp}{supp}
\DeclareMathOperator{\mes}{mes}
\DeclareMathOperator{\der}{der}
\DeclareMathOperator{\scnx}{sc}
\DeclareMathOperator{\SCNX}{SC}
\DeclareMathOperator{\Norm}{Norm}
\DeclareMathOperator{\Cent}{Cent}
\DeclareMathOperator{\el}{ell}
\DeclareMathOperator{\Out}{Out}
\DeclareMathOperator{\Aut}{Aut}
\DeclareMathOperator{\Gal}{Gal}
\DeclareMathOperator{\Hom}{Hom}
\DeclareMathOperator{\homc}{\mathcal{H}\mathit{om}}
\DeclareMathOperator{\inv}{inv}
\DeclareMathOperator{\Int}{Int}
\DeclareMathOperator{\Ind}{Ind}
\DeclareMathOperator{\Indp}{Ind-p}
\DeclareMathOperator{\Id}{Id}
\DeclareMathOperator{\stab}{stab}
\DeclareMathOperator{\aff}{aff}
\DeclareMathOperator{\loc}{loc}
\DeclareMathOperator{\geom}{geom}
\DeclareMathOperator{\Tran}{Tran}
\DeclareMathOperator{\Res}{Res}
\DeclareMathOperator{\cusp}{cusp}
\DeclareMathOperator{\disc}{disc}
\DeclareMathOperator{\Ker}{Ker}
\DeclareMathOperator{\Cok}{Cok}
\DeclareMathOperator{\Coker}{Coker}
\DeclareMathOperator{\Kil}{Kil}
\DeclareMathOperator{\rss}{rang_{ss}}
\DeclareMathOperator{\vect}{vect}
\DeclareMathOperator{\val}{val}
\DeclareMathOperator{\SL}{SL}
\DeclareMathOperator{\Spec}{Spec}
\DeclareMathOperator{\Spf}{Spf}
\DeclareMathOperator{\Div}{Div}
\DeclareMathOperator{\Tor}{Tor}
\DeclareMathOperator{\Ext}{Ext}
\DeclareMathOperator{\trace}{trace}
\DeclareMathOperator{\Sym}{Sym}
\DeclareMathOperator{\car}{\mathfrak{car}}
\DeclareMathOperator{\lng}{long}
\DeclareMathOperator{\Def}{Def}
\DeclareMathOperator{\card}{card}

%----------------------------------
%
%
%
%    MACROS MATHS
%
%
%
%
%
%----------------------------------
%-----------------------------Ensembles ---------------------
\newcommand{\ZZ}{\mathbb{Z}}
\newcommand{\Gm}{\mathbb{G}_m}
\newcommand{\Gmk}{\mathbb{G}_{m,k}}
\newcommand{\SG}{\mathfrak{S}}
\newcommand{\NN}{\mathbb{N}}
\newcommand{\RR}{\mathbb{R}}
\newcommand{\AAA}{\mathbb{A}}
\newcommand{\CC}{\mathbb{C}}
\newcommand{\BBB}{\mathbb{B}}
\newcommand{\QQ}{\mathbb{Q}}
\newcommand{\FF}{\mathbb{F}}
\newcommand{\Fq}{\mathbb{F}_q}
\newcommand{\Qp}{\mathbb{Q}_p}
\newcommand{\Ql}{\mathbb{Q}_{\ell}}
\newcommand{\Qlb}{\bar{\mathbb{Q}}_{\ell}}
\newcommand{\GA}{G(\mathbb{A})}
\newcommand{\CGA}{C^{\infty}_{c}(G(\mathbb{A}))}
\newcommand{\GQ}{G(\mathbb{Q})}
\newcommand{\ga}{\gamma}

%---------------Lettres calligraphi\'{e}es ------------------

\newcommand{\oc}{\mathcal{O}}
\newcommand{\rc}{\mathcal{R}}
\newcommand{\Sc}{\mathcal{S}}
\newcommand{\tc}{\mathcal{T}}
\newcommand{\uc}{\mathcal{U}}
\newcommand{\vc}{\mathcal{V}}
\newcommand{\wc}{\mathcal{W}}
\newcommand{\ec}{\mathcal{E}}
\newcommand{\gc}{\mathcal{G}}
\newcommand{\mc}{\mathcal{M}}
\newcommand{\yc}{\mathcal{Y}}
\newcommand{\lc}{\mathcal{L}}
\newcommand{\dc}{\mathcal{D}}
\newcommand{\fc}{\mathcal{F}}
\newcommand{\pc}{\mathcal{P}}
\newcommand{\hc}{\mathcal{H}}
\newcommand{\nc}{\mathcal{N}}
\newcommand{\Ac}{\mathcal{A}}
\newcommand{\bc}{\mathcal{B}}
\newcommand{\kc}{\mathcal{K}}
\newcommand{\Ic}{\mathcal{I}}
\newcommand{\Qc}{\mathcal{Q}}
\newcommand{\xc}{\mathcal{X}}

%--------------------Gros chapeaux
\newcommand{\Gc}{\widehat{G}}
\newcommand{\Tc}{\widehat{T}}
\newcommand{\Pc}{\widehat{P}}
\newcommand{\Bc}{\widehat{B}}
\newcommand{\Mc}{\widehat{M}}
\newcommand{\Lc}{\widehat{L}}
\newcommand{\Rc}{\widehat{R}}
\newcommand{\Sch}{\widehat{S}}
\newcommand{\Uc}{\widehat{U}}
\newcommand{\Hc}{\widehat{H}}
\newcommand{\Zc}{\widehat{Z}}

%--------------------Lettres grasses math\'{e}matiques ---------
\newcommand{\GG}{\mathbf{G}}
\newcommand{\HH}{\mathbf{H}}
\newcommand{\PP}{\mathbf{P}}
\newcommand{\QQQ}{\mathbf{Q}}
\newcommand{\MM}{\mathbf{M}}
\newcommand{\XX}{\mathbf{X}}
\newcommand{\LL}{\mathbf{L}}
\newcommand{\CCC}{\mathbf{C}}
\newcommand{\RRR}{\mathbf{R}}
\newcommand{\TT}{\mathbf{T}}
\newcommand{\BB}{\mathbf{B}}

%----------------Lettres gothiques ------------------------
\newcommand{\ggo}{\mathfrak{g}}
\newcommand{\kgo}{\mathfrak{k}}
\newcommand{\gA}{\mathfrak{g}(\mathbb{A})}
\newcommand{\gQ}{\mathfrak{g}(\mathbb{Q})}
\newcommand{\SgA}{\mathcal{S}(\mathfrak{g}(\mathbb{A}))}
\newcommand{\of}{\mathfrak{o}}
\newcommand{\Ogo}{\mathfrak{O}}
\newcommand{\mgo}{\mathfrak{m}}
\newcommand{\ngo}{\mathfrak{n}}
\newcommand{\ago}{\mathfrak{a}}
\newcommand{\cgo}{\mathfrak{c}}
\newcommand{\pgo}{\mathfrak{p}}
\newcommand{\qgo}{\mathfrak{q}}
\newcommand{\sgo}{\mathfrak{s}}
\newcommand{\rgo}{\mathfrak{r}}
\newcommand{\hgo}{\mathfrak{h}}
\newcommand{\tgo}{\mathfrak{t}}
\newcommand{\lgo}{\mathfrak{l}}
\newcommand{\ugo}{\mathfrak{u}}
\newcommand{\zgo}{\mathfrak{z}}
\newcommand{\Zgo}{\mathfrak{Z}}
\newcommand{\bgo}{\mathfrak{b}}
\newcommand{\Xgo}{\mathfrak{X}}
\newcommand{\Fgo}{\mathfrak{F}}

%--------------lettres avec une barre ---------------

\newcommand{\gb}{\bar{g}}
\newcommand{\qb}{\bar{q}}
\newcommand{\mb}{\bar{m}}
\newcommand{\nb}{\bar{n}}
\newcommand{\hb}{\bar{h}}
\newcommand{\xb}{\bar{x}}

%----------------------------------------------

\newcommand{\Gg}{G_{\gamma}}
\newcommand{\GAun}{G(\mathbb{A})^1}

\newcommand{\npa}{\ngo_P(\AAA)}
\newcommand{\sig}{\sigma_1^2}
\newcommand{\quo}{\GQ \backslash \GAun}
\newcommand{\quop}{P_1(\QQ) \backslash \GAun}

\newcommand{\rgoe}{\mathfrak{r}^\ast}
\newcommand{\mti}{\tilde{\mathfrak{\mgo}}}
\newcommand{\nbar}{\overline{\mathfrak{n}}}
\newcommand{\Oreg}{\mathcal{O}_{\reg}}

\newcommand{\hgoe}{\mathfrak{h}^\ast}
\newcommand{\Yb}{\overline{Y}}
\newcommand{\CT}{\mathcal{C}_0^2 (T_1)}

%-------------------Lettres grecques ---------------------------
\newcommand{\al}{\alpha}
\newcommand{\be}{\beta}
\newcommand{\Si}{\Sigma}
\newcommand{\zreg}{\zgo_{M,\reg}}
\newcommand{\am}{a_M}
\newcommand{\ams}{a_M^\ast}
\newcommand{\wm}{\wc(\am)}
\newcommand{\om}{\omega}
\newcommand{\Om}{\Omega}
\newcommand{\Lo}{\Lambda_\om}

\newcommand{\la}{\lambda}

\newcommand{\Ga}{\Gamma}
\newcommand{\back}{\backslash}
\newcommand{\ie}{i.e.\ }
\newcommand{\Cc}{C_c^\infty}

%-----------------------lettres à petit chapeau -------------
\newcommand{\ic}{\hat{i}}
\newcommand{\jc}{\hat{j}}
\newcommand{\ac}{\hat{a}}
\newcommand{\rf}{\hat{r}}

\newcommand{\geg}{\Ga^{\ec}(\ggo)}
\newcommand{\bg}{\langle}
\newcommand{\bd}{\rangle}

%--------------------lettres etoilees -----------------------
\newcommand{\Le}{L^\ast}
\newcommand{\Ge}{G^\ast}
\newcommand{\Me}{M^\ast}
\newcommand{\Te}{T^\ast}
\newcommand{\Se}{S^\ast}
\newcommand{\gem}{\Ga^{\ec}_{\Ge}(\mgo)}
\newcommand{\mgoe}{\mgo^\ast}
\newcommand{\ggoe}{\ggo^\ast}
\newcommand{\lgoe}{\lgo^\ast}
\newcommand{\sgoe}{\sgo^\ast}

\newcommand{\lce}{\lc^\ast}
\newcommand{\geme}{\Ga^{\ec}_{\Ge}(\mgoe)}
\newcommand{\gege}{\Ga^{\ec}(\ggoe)}
\newcommand{\gegell}{\Ga^{\ec}_{\el}(\ggo)}
\newcommand{\gag}{\Ga(\ggo)}
\newcommand{\gge}{\Ga(\ggoe)}
\newcommand{\fe}{f^\ast}
\newcommand{\dge}{\Delta_{\Ge}}
\newcommand{\dme}{\Delta_{\Me}}
\newcommand{\Fcl}{\overline{F}}
\newcommand{\kcl}{\overline{l}}

%------------------- L-groupes ---------------
\newcommand{\LG}{{}^LG}
\newcommand{\LM}{{}^LM}
\newcommand{\LT}{{}^LT}
\newcommand{\LLL}{{}^LL}
%---------------------Lettres pointees ---------------
\newcommand{\Gp}{\dot{G}}
\newcommand{\Mp}{\dot{M}}
\newcommand{\eps}{\varepsilon}
\newcommand{\epse}{\epsilon^\ast}
\newcommand{\Qe}{Q^\ast}
\newcommand{\Pe}{P^\ast}
\newcommand{\ome}{\omega^\ast}
\newcommand{\Ome}{\Omega^\ast}
\newcommand{\Gps}{G_{\eps}}
\newcommand{\Geps}{\Ge_{\epse}}

%--------------------------------
%
%
%   nouveaux <= et >=
%
%
%
%----------------------------
\renewcommand{\leq}{\leqslant}
\renewcommand{\geq}{\geqslant}

\usepackage[colorlinks=true,linkcolor=blue,citecolor=blue,urlcolor=blue,bookmarks=false,
            pdftitle={The Weighted Fundamental Lemma. I. Geometric Constructions},
            pdfauthor={Pierre-Henri Chaudouard, Gerard Laumon}]{hyperref}
\newcommand{\doi}[1]{\href{https://doi.org/#1}{\texttt{doi:#1}}}
\newcommand{\urlref}[1]{\href{#1}{\nolinkurl{#1}}}
\setcounter{tocdepth}{2}
\makeatletter
\def\l@subsection{\@tocline{2}{0pt}{1pc}{3.2em}{}}
\makeatother

\title[The Weighted Fundamental Lemma. I]{The Weighted Fundamental Lemma. I. Geometric Constructions}
\author{Pierre-Henri Chaudouard and G\'{e}rard Laumon}
\date{}
\subjclass{14D20 (primary); 11F70, 11F72, 11R39, 22E55 (secondary)}
\keywords{Hitchin fibration, Higgs bundles, weighted orbital integral, Arthur--Selberg trace formula, Langlands--Shelstad fundamental lemma, weighted fundamental lemma}

\begin{document}

\numberwithin{equation}{subsection}

\pagestyle{plain}
\raggedbottom
\hbadness=10000
\setlength{\emergencystretch}{2em}

\begin{abstract}
This work is the geometric part of our proof of the weighted fundamental lemma, which is an extension of Ng\^{o} Bao Ch\^{a}u's proof of the Langlands--Shelstad fundamental lemma. Ng\^{o}'s approach is based on a study of the elliptic part of the Hitchin fibration. The total space of this fibration is the algebraic stack of Hitchin bundles and its base space is the affine space of ``characteristic polynomials''. Over the elliptic set, the Hitchin fibration is proper and the number of points of its fibers over a finite field can be expressed in terms of orbital integrals. In this paper, we study the Hitchin fibration over an open set larger than the elliptic set, namely the ``generically regular semi-simple set''. The fibers are in general neither of finite type nor separated. By analogy with Arthur's truncation, we introduce the substack of $\xi$-stable Hitchin bundles. We show that it is a Deligne--Mumford stack, smooth over the base field and proper over the base space of ``characteristic polynomials''. Moreover, the number of points of the $\xi$-stable fibers over a finite field can be expressed as a sum of weighted orbital integrals, which appear in the Arthur--Selberg trace formula.
\end{abstract}

\maketitle

\begin{center}
\small English translation of \emph{Le lemme fondamental pond\'er\'e. I. Constructions g\'eom\'etriques}, Compositio Math. \textbf{146} (2010), 1416--1506,\\
\doi{10.1112/S0010437X10004756}\\[2pt]
The section, subsection, theorem and equation numbering of the published version has been preserved throughout.
\end{center}

\tableofcontents

\section{Introduction}

\begin{paragr}[Objectives.] ---
This work is the geometric part of our proof of the weighted fundamental lemma, whose strategy follows that developed by Ng\^{o} Bao Ch\^{a}u in his proof of the Langlands-Shelstad fundamental lemma.

This geometric part pursues two objectives. On the one hand, we truncate the Hitchin fibration so that the truncated Hitchin fibers are proper. On the other hand, we express the number of points over a finite field of such a fiber in terms of global weighted orbital integrals of Arthur that appear in the Arthur-Selberg trace formula; more precisely, these are variants of these integrals for Lie algebras.

Our constructions apply to any semisimple algebraic group and our results hold in this generality. However, for the purposes of this introduction, we shall restrict ourselves to the case of the group $\mathrm{SL}(n)$.
\end{paragr}

\begin{paragr}[The Hitchin Fibration for $\mathrm{SL}(n)$.] ---
Let $C$ be a projective, smooth, connected curve of genus $g$ over an algebraically closed field $k$, $n$ a positive integer, and $D$ an effective divisor on $C$ of degree $d>2g$. We assume that the characteristic of $k$ is either zero or greater than $n$.

A Hitchin bundle for the group $\mathrm{SL}(n)$ is a pair $(\mathcal{E},\theta)$ where $\mathcal{E}$ is a vector bundle of rank $n$ on $C$, equipped with a trivialization of its determinant, and $\theta:\mathcal{E}\rightarrow \mathcal{E}(D)$ is a twisted endomorphism of $\mathcal{E}$ with trace zero.

Let $\mathbb{M}$ denote the algebraic stack over $k$ of Hitchin bundles. The Hitchin fibration is the morphism
$$
f:\mathbb{M}\rightarrow\mathbb{A}
$$
with base the affine space
$$
\mathbb{A}=\bigoplus_{i=2}^{n}H^{0}(C,\mathcal{O}_{C}(iD)),
$$
which sends $(\mathcal{E},\theta)$ to the characteristic polynomial of $\theta$, denoted
$$
P_{a}(u):=u^{n}+a_{2}u^{n-2}+\cdots +a_{n},
$$
with $a_i \in H^{0}(C,\mathcal{O}_{C}(iD))$ for $2\leq i\leq n$.

For every $a=(a_2,\ldots,a_n)\in \mathbb{A}$, one has the spectral curve $Y_{a}$ defined by the equation
$$P_{a}(u)=0$$
viewed on the ruled surface $\Sigma_{D}=\mathbb{V}(\mathcal{O}_{C}(-D))$. This is a projective curve which, in general, is neither smooth nor irreducible, nor even reduced. The canonical projection $\pi_{a}:Y_{a}\rightarrow C$ is a finite covering of degree $n$.

Let $\mathbb{A}^{\mathrm{reg}}$ be the open subset where $Y_{a}$ is reduced, that is, where the equation $P_{a}(u)=0$ has only simple roots at the generic point of $C$. The fiber over $a\in \mathbb{A}^{\mathrm{reg}}$ of the Hitchin fibration can be interpreted as the stack of torsion-free $\mathcal{O}_{Y_a}$-modules $\mathcal{F}$ that are of rank $1$ at every generic point of $Y_a$ and that are equipped with a trivialization of the determinant of $\pi_{a,\ast}\mathcal{F}$ (cf. \cite{BNR89}). The open subset $\mathbb{A}^{\mathrm{reg}}$ contains the elliptic open subset $\mathbb{A}^{\mathrm{ell}}$ of those $a$ such that $Y_{a}$ is irreducible and hence integral. The Hitchin fiber at any point of the elliptic open is a proper Deligne-Mumford stack. In contrast, the Hitchin fiber at a non-elliptic point $a$ of $\mathbb{A}^{\mathrm{reg}}$ is an Artin stack which is neither separated nor of finite type.
\end{paragr}

\begin{paragr}[Truncation.] ---
In this work, we truncate the non-elliptic Hitchin fibers by a suitable notion of stability. Our construction depends on auxiliary data that we now introduce.

Let $\infty$ be a closed point of $C$ which is not in the support of $D$. We restrict to the open subset $\mathbb{A}^{\infty\textrm{-reg}}$ of $\mathbb{A}^{\mathrm{reg}}$ where $P_{a}(u)$ has only simple roots at the point $\infty$. This is not very restrictive since the opens $\mathbb{A}^{\infty\mathrm{-reg}}$ cover $\mathbb{A}^{\textrm{reg}}$ as the point $\infty$ varies. Let
$$\Ac\to  \mathbb{A}^{\mathrm{reg}}$$
be the finite \'{e}tale Galois cover with Galois group the symmetric group $\mathfrak{S}_{n}$, which consists, for each point $a$ of $\mathbb{A}^{\mathrm{reg}}$, in giving a total ordering on the roots of $P_{a}(u)$ at the point $\infty$. Let
$$
f:\mathcal{M}\rightarrow \mathcal{A}
$$
be the fibration deduced from the Hitchin fibration by the base change $\mathcal{A}\rightarrow \mathbb{A}$. A point of $\mathcal{M}$ is thus a triplet $(\mathcal{E},\theta,t)$ consisting of
\begin{itemize}
\item a Hitchin bundle $(\mathcal{E},\theta)$;
\item a point $t=(t_{1},\ldots ,t_{n})$ in $k^{n}$ whose coordinates are pairwise distinct and whose sum is zero;
\end{itemize}
which satisfy the following condition: the $n$-tuple $(t_{1},\ldots ,t_{n})$ is the ordered collection of eigenvalues of $\theta_{\infty}\in \mathrm{End}(\mathcal{E}_{\infty})$.

Let $\xi\in \mathbb{R}^{n}$ be such that $\xi_{1}+\cdots +\xi_{n}=0$. The fiber at $\infty$ of a vector subbundle $\mathcal{F}$ of $\mathcal{E}$ which is preserved by $\theta$ is a direct sum of eigenspaces for $\theta_\infty$, that is, one has
$$
\mathcal{F}_{\infty}=\bigoplus_{i\in I}\mathrm{Ker}(\theta_{\infty}-t_{i})
$$
for some subset $I\subset\{1,\ldots ,n\}$. We then set
$$
\mathrm{deg}_{\xi}(\mathcal{F})=\mathrm{deg}(\mathcal{F})+\sum_{i\in I}\xi_{i}.
$$

\begin{definition}
We say that $(\mathcal{E},\theta,t_{\infty})\in \mathcal{M}$ is $\xi$-stable if for every nonzero proper subbundle $\mathcal{F}\subsetneq \mathcal{E}$ preserved by $\theta$, one has
$$
\frac{\mathrm{deg}_{\xi}(\mathcal{F})}{\mathrm{rank}(\mathcal{F})} < 0.
$$
\end{definition}

For $\xi=0$, we recover the usual notion of stability. We show that the condition of $\xi$-stability is open and thus defines an open substack $\mathcal{M}^{\xi}$ of $\mathcal{M}$.

\begin{theoreme}\label{thm:sln}
Assume that $\sum_{i\in I}\xi_{i}\notin \mathbb{Z}$ for every nonempty proper subset $I\subsetneq\{1,\ldots ,n\}$. Then the stack $\mathcal{M}^{\xi}$ is a Deligne-Mumford stack proper over $\mathcal{A}$ which is smooth over $k$.
\end{theoreme}

Our definition of $\xi$-stability is reminiscent of a weighted stability definition due to Esteves \cite{Est01} in the context of compactified Jacobians. For a general reductive group, our definition is inspired by Arthur's truncations as well as by the work of Behrend (cf. \cite{Beh95}) on the stability and canonical reduction of group schemes. Our definition of stability is close to that for ordinary or Hitchin bundles with parabolic structure studied in \cite{HS} and \cite{BY96}.

The above theorem is proved using Langton's method \cite{Lan75}, as adapted by Nitsure \cite{Nit91} in the context of Hitchin bundles. In characteristic zero and for a general group, one may use Faltings' method (cf. \cite{Fal93}). Our approach is in fact an adaptation of an adelic argument due to Heinloth (cf. \cite{Hei08}), which has the advantage of still working in positive characteristic.
\end{paragr}

\begin{paragr}[Counting and Weighted Orbital Integrals.] ---
From now on, the base field $k$ is an algebraic closure of a finite field $\mathbb{F}_{q}$. We assume that all our data are defined over $\mathbb{F}_{q}$. Let $F$ be the function field of $C$, $\mathbb{A}$ its ring of adeles, and $\mathcal{O}$ the maximal compact subring of $\mathbb{A}$. Let $V$ be the set of closed points of $C$ and let $(d_v)_{v\in V}$ be the almost-zero family of positive integers such that $D=\sum_{v\in V}d_v v$. For every $v\in V$, let $\varpi_v$ be a uniformizer of the completion of the local ring at $v$. Set $\varpi^{-D}=(\varpi_v^{-d_v})_{v\in V}$. We have
$$
H^{0}(C,\mathcal{O}_{C}(iD))=F\cap \varpi^{-iD}\mathcal{O}.
$$

Let $(a,t)$ be an $\mathbb{F}_{q}$-rational point of $\mathcal{A}$. Let $P_{1},\ldots ,P_{s}$ be the irreducible factors of $P_{a}(u)$ viewed as an element of $F[u]$. These factors can be written at the point $\infty$ as
$$
P_{j}(\infty )(u)=\prod_{i\in I_{j}}(u-t_{i})
$$
for a unique partition $\{1,\ldots ,n\}=I_{1}\amalg\cdots\amalg I_{s}$. Let $(e_i)_{1\leq i\leq n}$ be the canonical basis of $k^n$ and, for every $I\subset\{1,\ldots,n\}$, let $V_I$ be the subspace of $k^n$ spanned by $e_i$ for $i\in I$. Let $M\subset \mathrm{SL}(n)$ be the stabilizer of the subspaces $V_{I_j}$ for $1\leq j\leq s$. This is a Levi subgroup. Let $\mathfrak{m}$ and $\mathfrak{sl}(n)$ be the Lie algebras of $M$ and $\mathrm{SL}(n)$, respectively. Let $X\in \mathfrak{m}(F)$ be a semisimple element. We assume that $X$ is regular in the sense that its centralizer $T_X$ in $\mathrm{SL}(n)$ is a maximal torus. Note that $T_X\subset M$. Let $\varphi$ be a locally constant function with compact support on $\mathfrak{sl}(n,\mathbb{A})$.

The weighted orbital integral $J_{M}(X,\varphi)$ is defined by the formula
$$
J_{M}(X,\varphi)=\int_{T_{X}(\mathbb{A})\backslash \mathrm{SL}(n,\mathbb{A})}
\varphi (g^{-1}X g)
\mathrm{v}_{M}(g)\frac{\mathrm{d}g}{\mathrm{d}t},
$$
where
\begin{itemize}
\item $\mathrm{d}g$ is the Haar measure on $\mathrm{SL}(n,\mathbb{A})$ normalized so that the maximal compact subgroup $\mathrm{SL}(n,\mathcal{O})$ has volume $1$;
\item $\mathrm{d}t$ is a Haar measure on the torus $T_{X}(\mathbb{A})$;
\item $\mathrm{v}_{M}(g)$ is Arthur's weight function defined in terms of volume and is left-invariant under $M(\mathbb{A})$.
\end{itemize}
The volume
$$\mathrm{vol}(T_{X}(F)\backslash T_{X}(\mathbb{A})^{1},\mathrm{d}t)$$
is finite, where $T_{X}(\mathbb{A})^{1}\subset T_{X}(\mathbb{A})$ is the intersection of the kernels of all homomorphisms $T_{X}(\mathbb{A})\rightarrow \mathbb{Z}$ obtained by composing an $F$-rational character $T_{X}\rightarrow \mathbb{G}_{m,F}$ with the degree map $\mathbb{A}^{\times}\rightarrow \mathbb{Z}$.

\begin{theoreme}\label{thm:slncomptage}
Assume that $\sum_{i\in I}\xi_{i}\notin \mathbb{Z}$ for every nonempty proper subset $I\subsetneq\{1,\ldots ,n\}$. Then the number of $\mathbb{F}_{q}$-rational points of the fiber over $(a,t)$ of the truncated Hitchin fibration $\mathcal{M}^{\xi}\rightarrow \mathcal{A}$ does not depend on $\xi$. Up to a coefficient which depends only on the normalization of Arthur's weight function, it is equal to the following sum of weighted orbital integrals of Arthur
$$
\sum_{X}\mathrm{vol}(T_{X}(F)\backslash T_{X}(\mathbb{A})^{1},\mathrm{d}t)
J_{M}(X,\mathbf{1}_D),
$$
where  $\mathbf{1}_D$ is the characteristic function of $\varpi^{-D}\mathfrak{sl}(n,\mathcal{O})$, and the sum is taken over the $M(F)$-conjugacy classes of elements $X\in \mathfrak{m}(F)$ such that the characteristic polynomial of the restriction of $X$ to $V_{I_j}$ is equal to $P_j$ for $1\leq j\leq s$.
\end{theoreme}
\end{paragr}

\begin{paragr}[Overview.] ---
Let us briefly describe the organization of this article. Some reminders and notations are given in Section \ref{sec:morcar}. Section \ref{sec:schema-car} introduces the base $\Ac$ of our Hitchin fibration and gives its essential properties. In Section \ref{sec:fib-Hitchin} we define the algebraic stack $\mc$ of Hitchin triples and the Hitchin morphism $f:\mc\to \Ac$. We also explain there the notion of reduction of a Hitchin triple to a parabolic subgroup and give a criterion for its existence and uniqueness. In Section \ref{sec:cvx}, we associate to each Hitchin triple a convex set which is constructed using the degrees of the reductions of that triple to the semi-standard parabolic subgroups. Using this convex set, we define in Section \ref{sec:xi-stab} the $\xi$-stability of a Hitchin triple. We also state there the main theorem of this article, which generalizes Theorem \ref{thm:sln} above (cf. Theorem \ref{thm:principal}). Section \ref{sec:description} provides an adelic description of Hitchin triples. Section \ref{sec:existence} is entirely devoted to the proof of the existence part of the valuative criterion of properness for the morphism $f$ restricted to the stack of $\xi$-semi-stable Hitchin triples. In Section \ref{sec:separation}, we prove that the restriction of the morphism $f$ to the stack of $\xi$-stable Hitchin triples is separated. In Section \ref{sec:DM}, we prove that the stack of $\xi$-stable Hitchin triples, for $\xi$ in ``general position\'', is a Deligne--Mumford stack. In Section \ref{sec:comptage}, we express the counting of the $\xi$-semi-stable Hitchin fibers in terms of Arthur\'s weighted orbital integrals (cf. Theorem \ref{thm:comptage}, which generalizes Theorem \ref{thm:slncomptage}). The article concludes in Section \ref{sec:rang1} with the study of an example of a non-elliptic Hitchin fiber for a semisimple group of rank $1$.
\end{paragr}

\section{Characteristic Morphism}\label{sec:morcar}

\begin{paragr}[Notations.] ---
Let $k$ be an algebraic closure of a finite field $\FF_q$. Let $G$ be a connected reductive algebraic group over $k$ and let $\ggo$ be its Lie algebra. In general, if an uppercase letter denotes a group, its Lie algebra is denoted by the corresponding lowercase Gothic letter. Let $\Int$ denote the action of $G$ on itself by inner automorphisms and $\Ad$ the adjoint representation. Let $T$ be a maximal subtorus of $G$ and $\tgo$ its Lie algebra. Let
$$W=W^G=W^G_T$$
be the Weyl group of $T$ in $G$, and
$$\Phi=\Phi^G=\Phi^G_T$$
the set of roots of $T$ in $G$. We denote by $|X|$ the cardinality of a finite set $X$. We assume that the order $|W|$ of the Weyl group is invertible in $k$.
\end{paragr}

\begin{paragr}[Characteristic Morphism.] \label{S:morcar}
Let $k[\ggo]$ be the $k$-algebra of regular functions on $\ggo$. The group $G$ acts on its Lie algebra by the adjoint action and, by duality, on $k[\ggo]$. Let $k[\ggo]^{G}$ denote the subalgebra of $G$-invariant functions and let
$$\car=\Spec(k[\ggo]^{G})$$
be the categorical quotient of $\ggo$ by $G$. Hence, we have a canonical $G$-invariant morphism
\begin{equation}
  \label{eq:chi}
  \chi \ :\ \ggo \to \car,
\end{equation}
called the characteristic morphism.

Let $k[\tgo]$ be the $k$-algebra of regular functions on $\tgo$ and let $k[\tgo]^W$ be the subalgebra of functions invariant under the Weyl group $W$. The restriction morphism $k[\ggo] \to k[\tgo]$ induces a morphism
$$k[\ggo]^G \to k[\tgo]^W,$$
which is in fact an isomorphism by Chevalley's theorem. It follows that, on the one hand, the restriction of the morphism $\chi$ to $\tgo$, still denoted $\chi$, induces an isomorphism
$$\tgo//W=\Spec(k[\tgo]^W)\simeq \car,$$
and on the other hand, there exist $n$ elements in $k[\ggo]^G$ (with $n=\rank(G)$), which are homogeneous and algebraically independent, that generate the $k$-algebra $k[\ggo]^G$ (see \cite{Bou68} Sections~5 and~6). The degrees of these elements depend only on the group $G$. In particular, the scheme $\car$ is isomorphic to the standard affine space of dimension $n=\rank(G)$.
\end{paragr}

\begin{paragr}[Discriminant.]
For every root $\al\in \Phi$, let $d\al\in k[\tgo]$ be its derivative. The discriminant $D^G$ defined by
$$D^G=\prod_{\al\in \Phi} d\al$$
belongs to $k[\tgo]^{W}$. Let $\car^{\reg}$ be the regular open subset of $\car$, i.e. the open subset where $D^G$ does not vanish.
\end{paragr}

\begin{paragr}[Some Properties of the Characteristic Morphism.]
By definition, an element of $\ggo$ is \emph{regular} if its $G$-orbit has maximal dimension, i.e. equal to $\dim(G)-\rank(G)$. Let $\ggo^{\reg}$ denote the open subset of $\ggo$ consisting of regular elements and
$$\tgo^{\reg}=\tgo\cap \ggo^{\reg}.$$

The morphism
$$\chi \ : \ \tgo \to \car$$
is a finite Galois covering with Galois group $W$. It induces on the regular opens an \'{e}tale covering $\tgo^{\reg}\to \car^{\reg}$.
\end{paragr}

\begin{paragr}[Kostant Section.] ---  \label{S:Kostant}
For every $\al\in \Phi$, the root space
$$\ggo_\al=\{X\in \ggo \mid \forall t\in T,\ \Ad(t)X=\al(t)X\}$$
is one-dimensional over $k$.

Let $B$ be a Borel subgroup of $G$ containing $T$, and let $\Delta\subset \Phi$ be the set of simple roots of $T$ in $B$. For every $\al\in \Delta$, let $X_\al$ be a nonzero element of $\ggo_\al$. Let $\al^\vee$ be the coroot corresponding to $\al$. Let $X_{-\al}$ be the unique element of $\ggo_{-\al}$ such that $[X_{\al},X_{-\al}]$ equals the coroot $\al^\vee$, viewed as an element of $\tgo$. Define
$$X_\pm=\sum_{\al\in \Delta}X_{\pm\al},$$
and let
$$\ggo_{X_+}=\{Y\in \ggo \mid [Y,X_+]=0\}.$$

According to a theorem of Kostant (cf. \cite{Kos63} or \cite[\S 2.4]{Kot99}), the affine space $X_- +\ggo_{X_+}$ is contained in the regular open $\ggo^{\reg}$, and the characteristic morphism $\chi \ :\ \ggo \to \car$ induces an isomorphism
$$X_- +\ggo_{X_+} \to \car.$$
The inverse morphism
$$\eps \ : \ \car \to  X_- +\ggo_{X_+}$$
is simply denoted $\eps$, or more precisely $\eps_{(B,\{X_\al\}_{\al\in \Delta})}$ if one wishes to indicate the dependence on the pinning $(B,\{X_\al\}_{\al\in \Delta})$.
\end{paragr}

\begin{paragr}[Parabolic Subgroups and Levi Subgroups.] --- \label{S:parab}
For any subgroups $M$ and $Q$ of $G$, denote by $\fc^Q(M)$ the set of parabolic subgroups $P$ of $G$ satisfying $M\subset P\subset Q$. Let $\fc=\fc^G(T)$ be the set of \emph{semi-standard} parabolic subgroups of $G$ (in the sense that they contain $T$).

A \emph{Levi subgroup} of $G$ is a Levi factor of a parabolic subgroup of $G$. Let $\lc=\lc^G(T)$ be the set of \emph{semi-standard} Levi subgroups of $G$ (in the sense that they contain $T$).

For every $P\in\fc$, let $N_P$ denote the unipotent radical of $P$, and let $M_P\in \lc$ be the unique Levi subgroup of $P$ that contains $T$. For every Levi subgroup $M\in \lc$, let $\pc(M)$, respectively $\lc(M)$, denote the subset of those $P\in \fc$ such that $M_P=M$, respectively those $L\in \lc$ such that $M\subset L$.
\end{paragr}

\begin{paragr}[The Schemes $\car_M$.] ---
Let $M\in \lc$. The notations of the preceding paragraphs, when subscripted or superscripted by $M$, apply to the reductive group $M$. Take care not to confuse the opens $\car_M^{G\textrm{-}\reg}$ and $\car_M^{M\textrm{-}\reg}$ of $\car_M$: the first is the locus where $D^G$ does not vanish, while the second is defined by $D^M\neq 0$. Since the latter does not appear in what follows, we set
$$\car_M^{\reg}=\car_M^{G\textrm{-}\reg}.$$

Similarly, we distinguish the $G$-regular open subset $\mgo^{G\textrm{-}\reg}$ and the $M$-regular open subset $\mgo^{M\textrm{-}\reg}$ of $\mgo$. We will only be interested in the former; hence we set
$$\mgo^{\reg}=\mgo^{G\textrm{-}\reg}.$$

Let $P\in \fc$ be a semi-standard parabolic subgroup and let $M=M_P$. Let $k[\pgo]$ be the algebra of regular functions on $\pgo$, and $k[\pgo]^P$ the subalgebra of $P$-invariant functions. Let $\car_P=\Spec(k[\pgo]^P)$ be the affine $k$-scheme, and consider the morphism
\begin{equation}
  \label{eq:chiP}
  \chi_P \ :\ \pgo \to \car_P,
\end{equation}
given by the inclusion $k[\pgo]^P\subset k[\pgo]$.

\begin{lemme}
The restriction morphism $k[\pgo]\to k[\mgo]$ induces an algebra isomorphism $k[\pgo]^P\simeq k[\mgo]^M$ and an isomorphism of schemes $\car_P\simeq \car_M$.
\end{lemme}

\begin{preuve}
Let $N=N_P$ and let $\mgo'\subset \mgo$ be the open subset consisting of $G$-regular semisimple elements. The open subset $\mgo'\oplus \ngo$ of $\pgo$ is the union of the $N$-conjugates of $\mgo'$. It follows that the morphism $k[\pgo]^P\to k[\mgo]^M$ is injective. It admits a section given by $\phi\in k[\mgo]^M \mapsto \phi\circ p\in k[\pgo]^P$, where $p$ is the projection $\mgo\oplus\ngo\to \mgo$.
\end{preuve}
\end{paragr}

\section{Characteristic Schemes}\label{sec:schema-car}

\begin{paragr}[Conventions.] --- \label{S:Conventions}
Let $S$ be a $k$-scheme. For any $k$-scheme $U$, let
$$U_S = U \times_k S$$
be the fiber product of $U$ and $S$ over $k$. For all $k$-schemes $U$ and $V$ and every $k$-morphism
$$\phi \ : \ U \to V,$$
let $\phi_S \ : \ U_S \to V_S$ be the $S$-morphism obtained by base change. When $S=\Spec(K)$ is the spectrum of a field $K$ extension of $k$, we simply denote by $U_K$ the $K$-scheme and by $\phi_K$ the corresponding $K$-morphism. If $s \in S$, we denote by $k(s)$ the residue field at $s$ and we set $U_s = U_{k(s)}$, etc.

Let $S$ be a $k$-scheme and $H$ an affine algebraic group over $k$. For any $H$-torsor $\ec$ over $S$ and any affine $S$-scheme $U$ equipped with a left action of $H$, the group $H$ acts on the fiber product $\ec \times_S U$ on the right by
$$ (e,u).h = (eh, h^{-1}u) $$
for all $h \in H$ and $(e,u) \in \ec \times_S U$. Let $\ec \times_S^H U$ denote the contracted product by $H$, i.e. the quotient of $\ec \times_S U$ by $H$. When $S=\Spec(k)$, we omit the index $S$.
\end{paragr}

\begin{paragr}[Notations.] --- \label{S:lacourbe}
Let $C$ be a projective, smooth, and connected curve over $k$ of genus $g$. Let $D$ be an effective divisor of degree greater than $2g$, and let $\infty$ be a point in $C(k)$ which is disjoint from the support of $D$.
Let $\lc_D$ be the $\Gmk$-torsor on $C$ associated to $D$, equipped with its canonical trivialization on the open complement of the support of $D$.
\end{paragr}

\begin{paragr}[The Bundles $\tgo_D$ and $\car_{M,D}$.] --- \label{S:tgoD}
For any $k$-scheme $V$ equipped with an action of $\Gmk$, let
$$V_D = \lc_D \times_k^{\Gmk} V$$
be the contracted product. If $V$ is a $k$-vector space and if the action of $\Gmk$ is linear, then $V_D$ is a vector bundle over $C$.

Let $M \in \lc$. The scaling action of the multiplicative group $\Gmk$ on $\tgo$ induces an action of $\Gmk$ on $\car_M$ for which the characteristic morphism $\chi_M$ is $\Gmk$-equivariant. Moreover, the respective actions of $\Gmk$ on $\tgo$ and $\car_M$ preserve the open subsets $\tgo^{\reg}$ and $\car^{\reg}_M$. By the above construction, one obtains a vector bundle $\tgo_D$ and a bundle $\car_{M,D}$, as well as the open subsets
$$\tgo^{\reg}_D \subset \tgo_D$$
and
$$\car^{\reg}_{M,D} \subset \car_{M,D}.$$
One also obtains a morphism, again abusively denoted $\chi_M$,
$$\chi_M \ : \ \tgo_D \to \car_{M,D},$$
which is a finite covering, Galois with group $W^M$, and \'{e}tale over the open subset $\car^{\reg}_D$.

When $M=G$, we omit the index $M$ in the notations.
\end{paragr}

\begin{paragr}[Characteristic Scheme $\Ac_M$.] --- \label{S:A}
Let $M \in \lc$. Define $\Ac_M^{G\textrm{-}\reg}$ to be the $k$-scheme such that for every $k$-scheme $S$, the set $\Ac_M^{G\textrm{-}\reg}(S)$ of $S$-points is the set of pairs
$$(a,t)$$
consisting of a section $a$ of the bundle $\car_{M,D,S}$ over $C_S$ and a point $t \in \tgo^{G\textrm{-}\reg}_D(S)$ satisfying
$$a(\infty_S) = \chi_M(t).$$
In what follows, we call $\Ac_M^{G\textrm{-}\reg}$ the characteristic scheme of $M$.

\begin{remarque}
The superscript $G\textrm{-}\reg$ indicates that the point $t \in \tgo^{G\textrm{-}\reg}_D(S)$ is $G$-regular in the sense that it belongs to $\ggo^{\reg}(S)$. Henceforth, we omit this superscript and, abusively, write
$$\Ac_M = \Ac_M^{G\textrm{-}\reg}.$$
When $M=G$, we omit the index $M$ and write
$$\Ac = \Ac_G.$$
\end{remarque}

\begin{proposition} \label{prop:irredA_M}
The scheme $\Ac_M$ is smooth and irreducible of dimension
$$\dim(\Ac_M) = \deg(D) \Bigl(\frac{|\Phi^M|}{2} + \rank(M)\Bigr) + \rank(M)(1-g).$$
\end{proposition}

\begin{preuve}
The forgetful morphism which omits the data $t$ makes $\Ac_M$ into a finite, \'{e}tale Galois cover with Galois group $W^M$ over the scheme of global sections of $\car_{M,D}$ that are $G$-regular at the point $\infty$. But this latter scheme is an open subset of the scheme of global sections of $\car_{M,D}$, which is non-canonically isomorphic to a vector space that one sees, by the Riemann--Roch formula, to have the announced dimension (once $\deg(D) > 2g-2$, cf. \cite[Lemma 4.4.1]{Ngo08}).

The forgetful morphism of the data $a$ is a surjective and open morphism from $\Ac_M$ to $\tgo^{\reg}$ whenever $\deg(D) > 2g-1$ (cf. \cite{Ngo08}, proof of Lemma 5.5.2). Its fibers are isomorphic to affine spaces. It follows that $\Ac_M$ is irreducible.
\end{preuve}
\end{paragr}

\begin{paragr}[Morphisms Between Characteristic Schemes.] --- Let $M\subset L$ be two semi-standard Levi subgroups. Let
$$\chi^M_L  \ :\ \car_M\to\car_L$$
be the morphism defined by the inclusion $k[\tgo]^{W^L} \subset k[\tgo]^{W^M}$. This is a finite morphism that is \'{e}tale over $\car_L^{\reg}$. As this morphism is $\Gmk$-equivariant, it induces a finite morphism, \'{e}tale over $\car_{L,D}^{\reg}$,
$$\chi^M_{L,D}  \ :\ \car_{M,D}\to\car_{L,D}.$$
Abusing notation, we still denote by $\chi^M_L$ the morphism
$$\chi^M_L \ : \ \Ac_M \to \Ac_L,$$
which, for every affine $k$-scheme $S$, sends $(a,t) \in \Ac_M(S)$ to the pair $(\chi^M_{L,D}(a), t) \in \Ac_L$.

\begin{proposition}\label{prop:immersionfermee}
The morphism
$$\chi^M_L \ : \ \Ac_M \to \Ac_L$$
is a closed immersion.
\end{proposition}

The proof of this proposition is postponed to §\ref{S:immersionfermee}.
\end{paragr}

\begin{paragr}[Elliptic Characteristics.] ---  \label{S:car-elliptique}
Let $M\in \lc(T)$. Define
$$\Ac_{M,\el} = \Ac_M - \bigcup_{L\in \lc,\ L\subsetneq M} \Ac_L.$$
We call $\Ac_{M,\el}$ the elliptic characteristic scheme of $M$.

\begin{proposition}
The scheme $\Ac_{M,\el}$ is a nonempty open subscheme of $\Ac_M$ and a locally closed subscheme of $\Ac$. Its closure in $\Ac$ is the closed subscheme $\Ac_M$.
\end{proposition}

\begin{preuve}
First, $\Ac_{M,\el}$ is nonempty for dimension reasons (cf. the dimension formula in Proposition \ref{prop:irredA_M}). The remainder of the first assertion follows from Proposition \ref{prop:immersionfermee}. The second assertion follows from the irreducibility of $\Ac_M$ (cf. Proposition \ref{prop:irredA_M}).
\end{preuve}

\begin{proposition} \label{prop:reunionAM}
The scheme $\Ac$ is the disjoint union of the locally closed subschemes $\Ac_{M,\el}$ for $M\in \lc(T)$, i.e.,
$$\Ac = \bigcup_{M\in \lc(T)} \Ac_{M,\el}.$$
\end{proposition}

The proof of this proposition is postponed to §\ref{S:reunionAM}.
\end{paragr}

\begin{paragr}[Proof of Proposition \ref{prop:immersionfermee}.] --- \label{S:immersionfermee}
This is a variant of the proofs of Lemma 7.3 in \cite{Ngo06} and Proposition 6.3.5 in \cite{Ngo08}. For the reader's convenience, we recall here the main arguments based in part on the following lemma (cf. Lemma 7.3 of \cite{Ngo06}).

\begin{lemme}\label{lem:relevtNgo}
Let $S$ be a normal, integral $k$-scheme and let $U\subset S$ be a nonempty open subscheme. Let $V'$ and $V$ be two $S$-schemes. For any commutative diagram
$$\xymatrix{
U \ar[d]_{i} \ar[r]^{h'} & V' \ar[d]^{\pi} \\
S  \ar[r]^{h} & V,
}$$
where $i$ is the canonical inclusion, $h'$ and $h$ are sections, and $\pi$ is a finite $S$-morphism, the section $h'$ extends uniquely to a section $S \to V'$ lifting $h$.
\end{lemme}

First, we show that the morphism $\chi^M_L$ is radiciel. That is, we must show that for every field $K$ extension of $k$, the morphism $\chi^M_L$ induces an injection on the corresponding sets of $K$-points. For $i=1,2$, let $(a_i,t_i) \in \Ac_M(K)$ be two elements having the same image $(a,t)$ in $\Ac_L$. Thus $t = t_1 = t_2$ and we have a commutative diagram
$$\xymatrix{
C_K \ar@<1ex>[r]^{a_1} \ar[r]_{a_2} \ar[rd]_{a} & \car_{M,D,K} \ar[d]^{\chi_{L}^M} \\
& \car_{L,D,K}.
}$$

Recall that the morphism $\chi^M_{L,D}$ is finite. By Lemma \ref{lem:relevtNgo} above, for $a_1=a_2$ it suffices that $a_1$ and $a_2$ coincide on an open subset of $C_K$. But $a_1$ and $a_2$ take the same value at the point $\infty_K$, namely $\chi_M(t)$. Since $t$ is $G$-regular and the morphism $\chi_L^M$ is \'{e}tale over $\car_{L,D,K}^{\reg}$, the sections $a_1$ and $a_2$ coincide on the spectrum of the completion of the local ring of $C_K$ at the point $\infty_K$, hence they coincide on an open subset of $C_K$.

\medskip

Next, we show that the morphism $\chi^M_{L,D}$ is proper. For this, we apply the valuative criterion of properness. Let $A$ be a discrete valuation ring with field of fractions $K$. Let $S = \Spec(A)$. Let $(a_M,t_M) \in \Ac_M(K)$ and $(a,t) \in \Ac_L(S)$ be such that $\chi_L^M((a_M,t_M))=(a,t)$. Thus $t_M=t$. By Lemma \ref{lem:relevtNgo} above, the section $a_1$ extends uniquely to a morphism from $C_S$ to $\car_{M,D}$ which fits (dotted) into the following commutative diagram:
$$\xymatrix{
C_K   \ar[r]^{a_1}  \ar[d] & \car_{M,D,S} \ar[d]^{\chi_{L}^M} \\
C_S \ar[r]^{a} \ar@{.>}[ur] & \car_{L,D,S}.
}$$
Once the extension of $a_1$ is obtained, it remains to verify that $a_1 \circ \infty_S = \chi^M_{L}(t)$. This equality follows from the properness of the morphism $\chi_{L}^M$ since the two sections $a_1 \circ \infty_S$ and $\chi^M_{L,D}(t)$ fit (dotted) into the commutative diagram
$$\xymatrix{
\Spec(K)   \ar[r]^{a_1\circ \infty_K}  \ar[d] & \car_{M,D,S} \ar[d]^{\chi_{L}^M} \\
S \ar[r]^{a\circ\infty_S} \ar@{.>}[ur] & \car_{L,D,S}.
}$$

Next, we show that the morphism $\chi^M_{L,D}$ is unramified. Since $\Ac_M$ and $\Ac_L$ are smooth, it suffices to show that the tangent map is injective. That is, one must show that $\chi_L^M$ induces an injection $\Ac_M(k[\eps])\to \Ac_L(k[\eps])$ (where $\eps$ is an indeterminate with $\eps^2=0$). So let $(a_1,t)$ and $(a_2,t)$ be two elements of $\Ac_M(k[\eps])$ having the same image $(a,t)$ in $\Ac_L(k[\eps])$. In what follows, we denote by a subscript $\eps$ the base change from $k$ to $k[\eps]$. Let $a_0$ be the section of $\car_{L,D}$ defined by restricting $a$ to $C$, and let $U \subset C$ be the open subset defined as the inverse image by $a_0$ of the regular locus $\cgo_{L,D}^{\reg}$. Note that $\infty \in U$. The section $a$ then induces a morphism $U_\eps \to \cgo_{L,D,\eps}^{\reg}$. Since the morphism $\chi_{L,\eps}^M$ is \'{e}tale, hence unramified, over $\cgo_{L,D,\eps}^{\reg}$, it follows that the sections $a_1$ and $a_2$ coincide on $U_\eps$, and hence on $C_\eps$ by flatness of $k[\eps]$ over $k$.

\medskip In conclusion, the morphism $\chi^M_{L,D}$, which is radiciel, proper, and unramified, is a closed immersion.
\end{paragr}

\begin{paragr}[Conjugation and the Characteristic Morphism.] ---
In this and the next paragraph, we gather some useful results for what follows and for the proof of Proposition \ref{prop:reunionAM}.

\begin{lemme} \label{lem:cjsep}
Let $P$ be a parabolic subgroup or a Levi subgroup of $G$. Let $F$ be an extension of $k$ and $F_s$ a separable closure of $F$. For all $X$ and $Y$ in $\pgo(F_s)$ and any $a\in \car_P^{\reg}(F_s)$ such that
$$\chi_P(X)=\chi_P(Y)=a,$$
there exists $p\in P(F_s)$ such that $X=\Ad(p)Y$.

If, moreover, we assume that $F$ is an extension of transcendence degree $1$ over an algebraically closed field and that the elements $X$ and $Y$ belong to $\pgo(F)$, then one may even take $p\in P(F)$.
\end{lemme}

\begin{preuve}
The $F_s$-scheme defined by
$$\{p\in P \mid \Ad(p)X=Y\}$$
is nonempty since it has points over an algebraic closure of $F$. It is clearly a torsor under the centralizer $T_X$ of $X$ in $G\times_k F_s$. Since $X$ is semisimple and regular, this centralizer is a torus over $F_s$. Hence the above scheme is smooth over $F_s$, so it has points in $F_s$.

Now suppose additionally that $X$ and $Y$ belong to $\pgo(F)$. Let $p\in P(F_s)$ be such that $X=\Ad(p)Y$. Then for all $\sigma\in\Gal(F_s/F)$, the element $p\sigma(p)^{-1}$ belongs to $T_X(F_s)$. Thus we obtain a 1-cocycle of $\Gal(F_s/F)$ with values in $T_X$, which is a coboundary since $T_X$ is connected; in particular, if $F$ is of dimension $\leq 1$, then $H^1(F_s/F, T_X)$ is trivial (see \cite{Ser94}, Chapter III, \S2 Theorem 1' and the following remarks). This allows one to conclude since an extension of transcendence degree $1$ over an algebraically closed field is of dimension $\leq 1$.
\end{preuve}

We now use the following definition of an elliptic element.

\begin{definition}
Let $F$ be an extension of $k$ and $X\in \ggo^{\reg}(F)$ a semisimple, $G$-regular element. Let $T_X$ be its centralizer in $G_F$ (a maximal $F$-torus of $G_F$) and let $A_X \subset T_X$ be the maximal split $F$-subtorus of $T_X$. We say that $X$ is an \emph{elliptic} element of $\ggo(F)$ if
$$A_X = A_G,$$
where $A_G$ is the connected component of the center of $G_F$.
\end{definition}
\end{paragr}

\begin{paragr}[Characteristics and Regular Semisimple Elements.] --- \label{S:lemmesAM}
Let $K$ be an extension of $k$. Let $F$ be the function field of the curve $C_K$ over $K$. Let $\oc_\infty$ be the completion of the local ring of $C_K$ at $\infty_K$, and let $F_\infty$ be the field of fractions of $\oc_\infty$. With these notations, we can state the following proposition.

\begin{proposition}\label{prop:caractdesAM}
For every $(a,t)\in \Ac_G(K)$, let $t_a\in \tgo^{\reg}(\oc_\infty)$ be the unique lift of $t$ such that the restriction of $a$ to $\Spec(\oc_\infty)$ is equal to $\chi(t_a)$. Let $a_\eta \in \car^{\reg}(F)$ be the restriction of $a$ to the generic point of $C_K$.

For every $M \in \lc(T)$ and $(a,t)\in \Ac_G(K)$, the following two assertions are equivalent:
  \begin{enumerate}
  \item $(a,t) \in \chi^M_G(\Ac_{M}(K))$;
  \item there exists $X$, a semisimple, $G$-regular element in $\mgo(F)$, and $m\in M(F_\infty)$ such that
    \begin{enumerate}
    \item $\chi_G(X) = a_\eta$;
    \item $\Ad(m)X = t_a$.
    \end{enumerate}
  \end{enumerate}
\end{proposition}

\begin{preuve}
Let $(a,t)\in \Ac_G(K)$. A point $t_a$ with the above properties exists and is unique since $\chi$ is \'{e}tale over $\car^{\reg}$.

First, we prove that (1) implies (2). Suppose $(a,t)\in \chi^M_G(\Ac_M(K))$, i.e. there exists $(a_M,t)\in \Ac_M(K)$ such that
\begin{equation}
  \label{eq:chi(aM)=a}
  \chi_G^M(a_M)=a,
\end{equation}
and
\begin{equation}
  \label{eq:am=chi(t)}
  a_M(\infty_K) = \chi_M(t).
\end{equation}
Let $a_{\infty}$ and $a_{M,\infty}$ be the restrictions of $a$ and $a_M$ to $\Spec(\oc_\infty)$, respectively. The morphisms given by $a_{M,\infty}$ and $\chi_M(t_a)$ coincide on the special fiber by (\ref{eq:am=chi(t)}) and both fit into the commutative diagram
$$\xymatrix{
\Spec(\oc_\infty) \ar@<1ex>[rr]^{\chi_M(t_a)} \ar[rr]_{a_{M,\infty}} \ar[rrd]_{a_\infty} & & \car_{M}^{\reg} \ar[d]^{\chi_{G}^M} \\
& & \car^{\reg}.
}$$
Since $\chi_{G}^M$ is \'{e}tale over $\car^{\reg}$, we have
$$a_{M,\infty}=\chi_M(t_a).$$
Let $a_{M,\eta}\in \car_{M}^{\reg}(F)$ be the restriction of $a_M$ to the generic point. Then there exists $X\in \mgo(F)$ such that
\begin{equation}
  \label{eq:chiM(gamma)=aM}
  \chi_M(X)=a_{M,\eta}.
\end{equation}
(to see this, one may use a Kostant section relative to $M$). Since $a_{M,\eta}$ is $G$-regular, $X$ is necessarily semisimple and $G$-regular. Thus, in $\car_M^{\reg}(F_\infty)$,
$$\chi_M(t_a) = \chi_M(X).$$

By Lemma \ref{lem:cjsep}, there exists a finite Galois extension $L$ of $F_\infty$ and $m\in M(L)$ such that
$$\Ad(m)X = t_a.$$
From this equality, for every $\sigma\in \Gal(L/F_\infty)$, the element $m\sigma(m)^{-1}$ normalizes $t_a$, so it belongs to $T(L)$. We thus obtain a $1$-cochain of the Galois group $\Gal(L/F_\infty)$ with values in $T(L)$ which is a coboundary, since $T$ being a split torus the group $H^1(L/F_\infty,T(L))$ is trivial; hence there exists $x\in T(L)$ such that $m\sigma(m)^{-1}=x\sigma(x)^{-1}$ for all $\sigma\in \Gal(L/F_\infty)$. Replacing $m$ by $x^{-1}m$, we may assume that $m\in M(F_\infty)$. This proves condition 2.(b). Condition 2.(a) follows from (\ref{eq:chi(aM)=a}) and (\ref{eq:chiM(gamma)=aM}).

\medskip

Now, we prove the converse. Suppose there exists $X\in \mgo(F)$, semisimple, $G$-regular, and $m\in M(F_\infty)$ satisfying conditions 2.(a) and 2.(b). Let
\begin{equation}
    a_{M,\eta}=\chi_M(X) \in \car_{M}^{\reg}(F).\label{eq:aMeta}
\end{equation}
This $F$-point fits into the commutative diagram
$$\xymatrix{
\Spec(F) \ar[r]^{a_{M,\eta}} \ar[d] & \car_{M,D} \ar[d]^{\chi_{G,D}^M} \\
C_K \ar[r]^{a} & \car_{G,D}.
}$$
The properness of the morphism $\chi_{G}^M$ implies that $a_{M,\eta}$ extends to a morphism
$$a_M \ : \ C_K \to \car_{M,D}$$
lifting $a$. We have the following equality in $\car_M^{\reg}(K)$:
\begin{equation}
  \label{eq:chi(t)=aM}
  \chi_M(t) = a_M(\infty_K).
\end{equation}
Indeed, by (\ref{eq:aMeta}), the restriction of $a_M$ to $\Spec(F_\infty)$ equals $\chi_M(X)$, hence equals $\chi_M(t_a)$ by 2.(b). It follows that the restriction of $a_M$ to $\Spec(\oc_\infty)$ is also equal to $\chi_M(t_a)$. Evaluating at $\infty_K$ yields (\ref{eq:chi(t)=aM}).

Thus, the pair $(a_M,t)$ defines an element of $\Ac_M(K)$. By the construction of $a_M$, we have
$$\chi_G^M(a_M,t)=(a,t),$$
which is assertion 1.
\end{preuve}

\begin{corollaire}\label{cor:caractdesAM}
With the notations of Proposition \ref{prop:caractdesAM}, the following two assertions are equivalent:
  \begin{enumerate}
  \item $(a,t) \in \chi_G^M(\Ac_{M,\el}(K))$;
  \item there exists $X$, a semisimple, $G$-regular, and elliptic element in $\mgo(F)$, and $m\in M(F_\infty)$ such that
    \begin{enumerate}
    \item $\chi_G(X)=a_\eta$;
    \item $\Ad(m)X = t_a$.
    \end{enumerate}
  \end{enumerate}
\end{corollaire}

\begin{preuve}
Let $(a,t)\in \Ac_G(K)$.

First, assume that $(a,t)\in \chi_G^M(\Ac_{M,\el}(K))$. By Proposition \ref{prop:caractdesAM}, there exist $X\in \mgo(F)$, semisimple and $G$-regular, and $m\in M(F_\infty)$ satisfying conditions 2.(a) and 2.(b) of Proposition \ref{prop:caractdesAM}. We will show that $X$ is necessarily elliptic in $\mgo(F)$. Suppose, for contradiction, that it is not. Let $A_X$ be the maximal split subtorus of $T_X$, the centralizer of $X$ in $G_F$. Up to conjugating $X$ by an element of $M(F)$ and translating $m$ by the same element, we may assume that $A_X$ is contained in the standard torus $T_F$. Then there exists $L\in \lc^G(T)$ such that the centralizer of $A_X$ in $G_F$ is $L_F$. Note that $L\subset M$ and that this inclusion is strict since $X$ is not elliptic in $M$. However, condition 2.(b) of Proposition \ref{prop:caractdesAM} implies
$$m T_{X,F_\infty}m^{-1} = T_{F_\infty}.$$
Since the tori $T_{X,F_\infty}$ and $T_{F_\infty}$ are two maximal $F_\infty$-split tori in $L_{F_\infty}$, they are conjugate by an element of $L(F_\infty)$. Hence,
$$m \in \Norm_{M(F_\infty)}(T) \, L(F_\infty).$$

Since $T$ is a split torus over $k$, we have $\Norm_{M(F_\infty)}(T)= \Norm_{M(k)}(T) \cdot T(F_\infty)$. Therefore, there exists $w\in  \Norm_{M(k)}(T)$ and $l\in L(F_\infty)$ such that
$$m = w l.$$
Set $L_1 = w L w^{-1}$, $l_1 = wlw^{-1}$, and $X_1 = \Ad(w)X$. By construction, $L_1\in \lc(T)$ is a proper Levi subgroup of $M$; moreover, $l_1\in L_1(F_\infty)$ and $X_1\in \lgo_1(F)$ satisfy $\chi_G(X_1) = a_\eta$ and $\Ad(l_1)X_1 = t_a$. It then follows from Proposition \ref{prop:caractdesAM} that $(a,t)\in \chi^{L_1}_G(\Ac_{L_1}(K))$. This contradicts our assumption, and hence $X$ is elliptic in $\mgo(F)$.

\medskip

Now, assume that assertion 2 holds. That is, there exist $X\in \mgo(F)$, semisimple, $G$-regular, and elliptic, and $m\in M(F_\infty)$ such that $\chi_G(X)=a_\eta$ and $\Ad(m)X=t_a$. Then, by Proposition \ref{prop:caractdesAM}, we have $(a,t)\in \chi_G^M(\Ac_M(K))$. Let $L\in \lc(T)$ be a semi-standard Levi subgroup contained in $M$ such that $(a,t)\in \chi^L_G(\Ac_{L,\el}(K))$. Then there exists $Y\in \lgo(F)$, semisimple, $G$-regular, and elliptic, and $l\in L(F_\infty)$ such that $\chi_G(Y)=a_\eta$ and $\Ad(l)Y=t_a$. By Lemma \ref{lem:cjsep}, the equality $\chi_G(X)=\chi_G(Y)$ implies the existence of $g\in G(F_s)$, for some separable closure $F_s$ of $F$, such that $\Ad(g)X = Y$. The centralizer of $Y$ in $G_F$ is a maximal $F$-torus. For every $\sigma \in \Gal(F_s/F)$, we have $\sigma(g)g^{-1} \in T_Y$, and the inner automorphism $\Int(g)$ of $G_{F_s}$ induces an $F$-isomorphism from $T_X$ to $T_Y$, hence an $F$-isomorphism between the maximal $F$-split subtori $A_X$ and $A_Y$. But these are nothing other than the connected centers $A_{M,F}$ and $A_{L,F}$ of the Levi subgroups $M_F$ and $L_F$. Thus $\dim(A_M) = \dim(A_L)$. Since we also have $A_M\subset A_L$, it follows that $A_M = A_L$ and hence $M = L$. Therefore, $(a,t)\in \chi_G^M(\Ac_{M,\el}(K))$, as desired.
\end{preuve}
\end{paragr}
\begin{paragr}[Proof of Proposition \ref{prop:reunionAM}.] --- \label{S:reunionAM}
In this paragraph, the only non-obvious point is that the union is disjoint. We resume the notations from §\ref{S:lemmesAM}. Let $(a,t)\in \Ac_G(K)$ and let $M_1$ and $M_2$ be two Levi subgroups in $\lc(T)$ such that $a$ lies in the intersection of the images of $\Ac_{M_1,\el}(K)$ and $\Ac_{M_2,\el}(K)$. By Corollary \ref{cor:caractdesAM}, for $i=1,2$, there exist $X_i\in \mgo_i(F)$, elliptic and $G$-regular, and $m_i\in M_i(F_\infty)$ such that $\Ad(m_i)X_i=t_a$ and $\chi_G(X_i)=a_\eta$. By Lemma \ref{lem:cjsep}, the equality $\chi_G(X_1)=\chi_G(X_2)$ implies the existence of $g\in G(F_s)$, for a separable closure $F_s$ of $F$, such that $\Ad(g)X_1=X_2$. Arguing as in the proof of Corollary \ref{cor:caractdesAM}, one deduces that
$$ g A_{M_1} g^{-1} = A_{M_2}.$$
But since $g$ conjugates $X_1$ and $X_2$, we have that $m_2 g m_1^{-1}$ centralizes $t_a$, which is $G$-regular. Hence, $m_2 g m_1^{-1} \in T(F_s \otimes_F F_\infty)$, so that $A_{M_1} = A_{M_2}$ and consequently $M_1 = M_2$, as required.
\end{paragr}
\bigskip
\section{The Hitchin Fibration}\label{sec:fib-Hitchin}

We continue with the notations from the previous sections.
\medskip

\begin{paragr}[Notations.] --- \label{S:VD} We complete the notations of Sections \ref{S:Conventions} through \ref{S:tgoD}. Let $S$ be a $k$-scheme and $\ec$ a $G$-torsor over $S$. Define
  $$\Aut_G(\ec)= \ec\times_k^{G,\Int} G$$
to be the group scheme over $S$ of $G$-equivariant automorphisms of $\ec$.

Let $V$ be a finite-dimensional $k$-vector space. For any linear algebraic representation
$$\rho \ :\ G \to GL(V),$$
let
$$\rho(\ec)=G\times_k^{G,\rho} V$$
be the vector bundle over $S$ obtained by pushing forward $\ec$ via the representation $\rho$.

When $\ec$ is a $G$-torsor over $C_S=C\times_k S$ (where $C$ is the curve of \S\ref{S:lacourbe}), we set
$$\rho_D(\ec)=\ec \times_{C_S}^{G,\rho} V_{D,S},$$
where $V_{D,S}=V_D\times_k S$ is the vector bundle over $C_S$ deduced from the vector bundle $V_D$ on $C$ (see \S\ref{S:tgoD}) by base change.
\end{paragr}

\begin{paragr}[Generic trivialization of $G$-torsors.] --- Let $K$ be an extension of $k$ and $\ec$ a $G$-torsor over $C_K$. In what follows, we will make repeated use of the following lemma.

  \begin{lemme}\label{lem:Sen}
    There exists a finite extension $K'$ of $K$ such that the $G$-torsor $\ec\times_{C_K} C_{K'}$ is trivial at the generic point of $C_{K'}$.
  \end{lemme}

  \begin{preuve} Clearly, it suffices to prove the lemma when $K$ is algebraically closed. In this case, the field of functions $F$ of $C_K$ has dimension $\leq 1$. Let $F_s$ be a separable closure of $F$. The torsor $\ec$, which is smooth over $C_K$, admits sections over $F_s$. Since $G$ is connected and $F$ has dimension $\leq 1$, we know that the set $H^1(F_s/F,G)$ is trivial (see \cite{Ser94} chap. III \S2 theorem 1' and the remarks that follow).
 \end{preuve}

\end{paragr}

\begin{paragr}[Reductions of torsors to a subgroup.] \label{S:reduc} --- We continue with the notations of the previous paragraph. Let $P$ be an algebraic subgroup of $G$. A reduction of the $G$-torsor $\ec$ over $S$ to $P$ is the data of a $P$-torsor $\ec_P$ over $S$ together with an isomorphism of the $G$-torsor
$$\ec_P\times^{P} G,$$
where $G$ is equipped with the left translation action by $P$, with $\ec$. The obvious $P$-equivariant morphism $\ec_P\to  \ec_P\times^{P} G\simeq \ec$ gives, upon passing to the quotient by $P$, a section
$$\sigma_P \ :\ S \to \ec/P.$$
Conversely, $\ec$ is a $P$-torsor over $\ec/P$ which, when pulled back by a section of $\ec/P$, yields a reduction of $\ec$ to $P$. In this way, we obtain a natural bijection between the sections of $\ec/P$ and the isomorphism classes of reductions of $\ec$ to $P$.
\end{paragr}

\begin{paragr}[Characteristic morphism.] --- Recall that in \S\ref{S:lacourbe} a curve $C$ over $k$ and a point $\infty\in C(k)$ were fixed. Let $S$ be a $k$-scheme and denote by $C_S$ and $\infty_S$ the objects obtained by base change. Let $\ec$ be a $G$-torsor over $C_S$. The construction of paragraph \ref{S:VD} gives a vector bundle $\Ad_D(\ec)$ over $C_S$. The characteristic morphism defined in \S\ref{S:morcar} (\ref{eq:chi})
$$\chi \ :\ \ggo \to \car$$
is $G$-invariant and $\Gmk$-equivariant. It therefore induces a morphism, by abuse of notation still denoted by $\chi$,
$$\chi  \ :\ \Ad_D(\ec) \to \car_{D,S},$$
and hence a morphism, still denoted by the same symbol,
$$\chi \ :\ H^0(C_S,\Ad_D(\ec)) \to H^0(C_S,\car_{D,S}).$$

The group scheme $\Aut_G(\ec)$ acts on the bundle $\Ad_D(\ec)$ by the adjoint action. For any section $\theta\in H^0(C_S,\Ad_D(\ec))$, denote by
$$\Aut_G(\ec,\theta)$$
the subgroup scheme of $\Aut_G(\ec)$ that centralizes $\theta$.
\end{paragr}

\begin{paragr}[The Hitchin Algebraic Stack.] --- \label{S:Cham-Hitchin}
We now introduce the following definition.

  \begin{definition}
The Hitchin stack $\mc_G$ is the $k$-algebraic stack whose fiber category over a $k$-scheme $S$ is the groupoid whose objects are triples $(\ec,\theta,t)$ satisfying the following conditions:
  \begin{enumerate}
  \item $\ec$ is a $G$-torsor over $C\times_k S$;
  \item $\theta$ is a section of the vector bundle $\Ad_D(\ec)$;
  \item $t$ belongs to $\tgo^{G\textrm{-}\reg}(S)$ and satisfies
$$\chi(\theta)(\infty_S)=\chi(t).$$
  \end{enumerate}
and the set of morphisms from an object $(\ec,\theta,t)$ to another $(\ec',\theta',t')$ is given by
\begin{itemize}
\item the empty set if $t\not= t'$;
\item the set of $G$-equivariant isomorphisms from $\ec$ to $\ec'$ such that the induced isomorphism from $\Ad_D(\ec)$ to $\Ad_D(\ec')$ sends $\theta$ to $\theta'$.
\end{itemize}
\end{definition}

\begin{remarque} This is not the usual definition of the Hitchin stack which only considers pairs $(\ec,\theta)$. Here the additional datum of an element $t$ that is $G$-regular forces the section $\theta$ to be regular semisimple at $\infty_S$, and hence generically as well. The morphism which forgets $t$ makes our Hitchin stack an \'{e}tale Galois covering with group $W$ of an open substack of the usual Hitchin stack.
\end{remarque}

\end{paragr}

\begin{paragr}[The Hitchin Fibration.] --- This is the morphism
$$f_G \ : \ \mc_G \to \Ac_G,$$
defined for every $k$-scheme $S$ and every triple $m=(\ec,\theta,t)$ in $\mc(S)$ by
$$f(m)=(\chi(\theta),t).$$
In what follows, we will often omit the index $G$ and denote the above morphism by
$$f\ :\ \mc\to \Ac.$$

\end{paragr}

\begin{paragr}[The Morphism $\chi_P$.] --- Let $K$ be an extension of $k$ and $P\in \fc$ a semi-standard parabolic subgroup. Let $\ec$ be a $P$-torsor over $C_K$. Let $\Ad$ denote the adjoint action of $P$ on its Lie algebra $\pgo$. The construction of paragraph \ref{S:VD} applied to the group $P$ provides a vector bundle $\Ad_D(\ec)$ over $C_K$. The characteristic morphism defined in \S\ref{S:parab} (\ref{eq:chiP})
$$\chi_P  \ :\ \pgo \to \car_P$$
is $P$-invariant and $\Gmk$-equivariant. It therefore induces, by abuse of notation, morphisms still denoted by $\chi_P$,
$$\chi_{P}  \ :\ \Ad_D(\ec) \to \car_{P,D,S},$$
and
$$\chi_P \ :\ H^0(C_S,\Ad_D(\ec)) \to H^0(C_S,\car_{P,D,S}).$$

\end{paragr}

\begin{paragr}[Reduction to a Parabolic Subgroup.] --- Let $S$ be an affine $k$-scheme and $m=(\ec,\theta,t)\in \mc(S)$. We introduce the following definition.

  \begin{definition} \label{def:reduc}
    A reduction of $m$ to $P$ is a triple $m_P=(\ec_P,\theta_P,t_P)$ consisting of

\begin{itemize}
    \item a reduction $\ec_P$ of $\ec$ to $P$, that is, a $P$-torsor $\ec_P$ over $C\times_k S$ together with an isomorphism
\begin{equation}
        \label{eq:onpousse}
        \ec_P\times_k^P G \simeq \ec \ ;
      \end{equation}
    \item a section $\theta_P$ of $\Ad_D(\ec_P)$ over $C\times_k S$;
    \item a point $t_P\in \tgo^{G\textrm{-}\reg}(S)$ satisfying
$$\chi_P(\theta_P)(\infty_S)=\chi_P(t_P);$$
\end{itemize}
which furthermore satisfies the following two conditions:
  \begin{enumerate}
  \item The isomorphism, deduced from (\ref{eq:onpousse}),
$$\Ad_D(\ec_P\times_k^P G)\simeq \Ad_D(\ec),$$
composed with the morphism
$$\Ad_D(\ec_P)\hookrightarrow \ec_P\times_k^{P,\Ad}\ggo \simeq  \Ad_D(\ec_P\times_k^P G),$$
sends the section $\theta_P$ to $\theta$;
\item The points $t_P$ and $t$ are equal.
\end{enumerate}
  \end{definition}

In what follows, we say that two reductions $(\ec_P,\theta_P,t)$ and $(\ec_P',\theta_P',t)$ of $m$ are isomorphic if there exists an isomorphism of $P$-torsors from $\ec_P$ to $\ec_P'$ such that the induced isomorphism $\Ad_D( \ec_P) \to \Ad_D( \ec_P')$ sends $\theta_P$ to $\theta_P'$.
\end{paragr}

\begin{paragr}[Existence of Reductions on an Open Set.] --- \label{S:reductouvert} Let $S$ be a test $k$-scheme and $(\ec,\theta,t)\in \mc(S)$. Let $M\in \lc$ and $(a_M,t)\in \Ac_M(S)$ such that
$$f(m)=\chi_G^M((a_M,t)).$$
  Let $U$ be the open subset of $C_S$ defined as the inverse image via $a_M$ of the open set $\car_M^{\reg}$. Let $X$ be a section of $\mgo\times_k U$ such that $\chi_M(X)$ coincides with $a_M$ on the open set $U$. The existence of a Kostant section guarantees that such a section exists. Note that $X$ is everywhere $G$-regular semisimple. Let $T_X$ be the torus subscheme of $M_U$ centralizing $X$ in $G_U$. We then define $\ec'$ over $U$ as follows: for any $U$-scheme $\Omega$, the sections of $\ec'$ over $\Omega$ are the sections of $\ec$ over $\Omega$ such that the induced isomorphism $\Ad_D(\ec_{\Omega}) \to \ggo \times_k \Omega$ sends $\theta_\Omega$ to $X_\Omega$. Since locally for the \'{e}tale topology the torus subschemes of $G\times_k U$ are conjugate to the constant torus $T\times_k U$, and since we have the equality
$$\chi_G(\theta)=\chi_G^M(a_M)=\chi_G(X)$$
on $U$, we see that $\ec'$ has sections locally for the \'{e}tale topology. It is then clear that $\ec'$ is a $T_X$-torsor over $U$ which, when pushed forward via the morphism $T_X \to G_U$, becomes isomorphic to $\ec_U$. In particular, for every parabolic subgroup $P\in \fc(M)$, let $\ec_P$ be the $P$-torsor over $U$ obtained from $\ec'$ by pushing forward via the morphism $T_X\to P_U$. One then verifies that $(\ec_P,\theta_U,t_U)$ is a (natural) reduction of the triple obtained from $(\ec,\theta,t)$ by base change to $U$.
\end{paragr}

\begin{paragr}[Existence and Uniqueness of Reductions.] --- We now state the main result concerning the existence and uniqueness of reductions to a parabolic subgroup of a Hitchin triple over the spectrum of a field $K$ extending $k$.

  \begin{proposition}\label{prop:Preduc}
    Let $m\in \mc(K)$ and $M\in \lc$ such that $f(m)\in \chi^M_G(\Ac_{M,\el}(K))$ (cf. Proposition \ref{prop:reunionAM}). Let $P\in \fc$. We have the following alternative:
\begin{itemize}
\item either $P$ does not contain $M$: in this case, there is no reduction of $m$ to $P$;
\item or $P$ contains $M$: in this case, there exists a unique reduction of $m$ to $P$ (up to isomorphism).
\end{itemize}
\end{proposition}

\begin{preuve} Let $P\in \fc$. We first show that the condition $M\subset P$ is necessary for the existence of a reduction to $P$. Suppose $m_P=(\ec_P,\theta_P,t)$ is a reduction of $m$ to $P$. The pair $(\chi_P(\theta_P),t)$ defines an element of $\Ac_{M_P}(K)$ whose image in $\Ac_G(K)$ equals $f(m)$. The hypothesis $f(m)\in \chi^M_G(\Ac_{M,\el}(K))$ and Proposition \ref{prop:reunionAM} imply that $M\subset M_P$, hence $M\subset P$.
\medskip

Now assume that $M\subset P$. We first prove the existence of a reduction to $P$. The construction in paragraph \S\ref{S:reductouvert} provides a reduction of $m$ to $P$, but only over an open subset $U$ of $C_K$. It remains to show that this reduction extends, i.e., that the $P$-torsor $\ec_P$ extends to all of $C_K$. By the bijection described in \S\ref{S:reduc}, this amounts to showing that the section $U\to \ec/P$ deduced from $\ec_P$ extends to all of $C_K$. But such a section extends by the properness of the quotient $\ec/P$ over $C$, hence existence is proved.

\medskip

Finally, we prove the uniqueness of a reduction of $m$ to $P$. Let $m_P=(\ec_P,\theta_P,t)$ and $m'_P=(\ec_P',\theta_P',t)$ be two such reductions. Let $\sigma_P$ and $\sigma_P'$ be the sections of $\ec/P$ deduced from $\ec_P$ and $\ec_P'$ by the construction described in \S \ref{S:reduc}. It suffices to prove that these sections are equal, and by the properness of $\ec/P$ over $C_K$, it is enough to check at the generic point $\Spec(F)$ of $C_K$, and even at the point $\Spec(F_s)$ where $F_s$ is a separable closure of the function field $F$ of $C_K$.

Now, the section $\sigma_P$ gives a section $\Spec(F_s)\to \ec/P$, which is the $P(F_s)$-orbit of the image of a point $e\in \ec_P(F_s)$ via the morphism $\ec_P\to \ec$. Such a point $e\in \ec_P(F_s)$ gives a $P$-equivariant isomorphism
$$\Phi_P \ : \ \ec_{P,F_s}\to P\times_k F_s,$$
an isomorphism of group schemes
$$\iota_P \ : \ \Aut_P(\ec_{F_s}) \to P\times_k F_s,$$
and a derived isomorphism
$$d\iota_P \ : \ \Ad_P(\ec_{F_s}) \to \pgo \times_k F_s.$$
Let $Y=d\iota_P(\theta_{F_s})\in \pgo(F_s)$. By the definition of the reduction, $\chi_P(Y)$ equals $a_{M,F_s}$, the restriction of the section $a_M  \ : C \to \cgo_M$ to $\Spec(F_s)$. Using Proposition \ref{prop:caractdesAM} for the group $M$ (instead of $G$), we obtain the existence of $X\in \mgo(F)$ semisimple such that
\begin{enumerate}
\item $\chi_M(X)$ is equal to the restriction $a_{M,F}$ of $a_M$ at the generic point;
\item there exists $m\in M(F_\infty)$ such that $\Ad(m)X=t_a$, where $t_a\in \tgo^{\reg}(\oc_\infty)$ lifts $t$ and satisfies $\chi_M(t_a)=t$.
\end{enumerate}
Note that since $a_{M,F}$ belongs to $\car_M^{G\textrm{-}\reg}$, the elements $X$ and $t_a$ are automatically $G$-regular. Since we have the equality $\chi_P(Y)=\chi_P(X)$ in $\car_M^{G\textrm{-}\reg}(F_s)$, by Lemma \ref{lem:cjsep} the elements $X$ and $Y$ are conjugate under $P(F_s)$. After possibly translating $e$ by an element of $P(F_s)$, we may assume that $Y=X$. The image of $e$ in $\ec(F_s)$ gives a $G$-equivariant isomorphism
 \begin{equation}
   \label{eq:Phi}
   \Phi \ : \ \ec_{F_s}\to G\times_k F_s,
 \end{equation}
 an isomorphism of group schemes
 \begin{equation}
   \label{eq:iota}
   \iota \ : \ \Aut_G(\ec_{F_s}) \to G\times_k F_s,
 \end{equation}
 and a derived isomorphism
 \begin{equation}
   \label{eq:diota}
   d\iota \ : \ \Ad_G(\ec_{F_s}) \to \ggo \times_k F_s.
 \end{equation}
 This latter isomorphism is such that $X=d\iota(\theta)$ is an element of $\mgo(F)$ which is semisimple and $G$-regular and satisfies the assertions 1 and 2 above. Denote by a prime the objects attached to the reduction $m_P'$. Let $x\in G(F_s)$ be such that $e'=e\cdot x$. The automorphism of $\ggo\times_k F_s$ given by $d\iota'\circ d\iota^{-1}$ is then equal to $\Ad(x)$. Consequently, $\Ad(x)X=X'$. Hence, we have the equalities
$$t_a=\Ad(m')X'=\Ad(m'x)X=\Ad(m'xm^{-1})t_a,$$
which implies that $m'xm^{-1}\in T(F_\infty\otimes_F F_s)$ since $t_a$ is $G$-regular. It follows that $x\in M(F_s)$. Thus, the sections $\sigma_P$ and $\sigma'_P$ coincide at the generic point. Therefore, there exists an isomorphism from $\ec_P$ to $\ec_P'$. With the previous notation, we have seen that $X$ and $X'$ are conjugate under $M(F_s)$. Hence, we can choose this isomorphism so that $\theta_P$ is mapped to $\theta'_P$ (at least at the point $\Spec(F_s)$, and hence everywhere). This completes the proof of uniqueness and the proposition.
\end{preuve}
\end{paragr}
\section{Convex Set Associated to a Hitchin Triple}\label{sec:cvx}
\begin{paragr}[The Spaces $\ago_M$.] --- \label{S:espacesa}
  For any group $M$ defined over $k$, let $X^*(M)$ be the group of rational characters of $M$ and
$$X_*(M)=\Hom_{\ZZ}(X^*(M),\ZZ).$$
When $M$ is a torus, we note that $X_*(M)$ identifies with the group of cocharacters of $M$ via the natural duality between the groups of characters and cocharacters.

Let $M\in \lc$ be a semi-standard Levi subgroup and let $A_M$ be the connected component of the center of $M$. The restriction morphism
$$X^*(M)\to X^*(A_M)$$
is injective and its cokernel is finite. Let $\ago_M^*$ be the $\RR$-vector space defined by
$$\ago_M^*=X^*(M)\otimes_\ZZ \RR=X^*(A_M)\otimes_\ZZ \RR.$$
Its dual, denoted $\ago_M$, identifies with
$$\ago_M=\Hom_\ZZ(X^*(M),\RR).$$
We will not confuse $\ago_M$ with the Lie algebra of $A_M$. For any $L\in \lc(M)$, the restriction morphisms
$$X^*(L)\to X^*(M)\to X^*(A_L)$$
induce morphisms
$$\ago_L^*\to \ago_M^*\to \ago_L^*$$
whose composition is the identity. In this way, $\ago_L^*$ identifies with a subspace of $\ago_M^*$. Let
$$(\ago_M^L)^*=\Ker(\ago_M^*\to \ago_L^*).$$
Thus, we have a decomposition
$$\ago_M^*=\ago_L^*\oplus (\ago_M^L)^*$$
and, dualizing, a decomposition
$$\ago_M=\ago_L\oplus \ago_M^L$$
where $\ago_M^L$ is the orthogonal complement of $\ago_L^*$ in $\ago_M$. Every $\xi\in \ago_M$ can be written
\begin{equation}
  \label{eq:decompoLevi}
  \xi=\xi_L+\xi^L
\end{equation}
according to this decomposition.

\medskip We now generalize the preceding definitions to a parabolic subgroup $P\in \fc$ as follows. Let $A_P=A_{M_P}$ and $\ago_P=\ago_{M_P}$. For any $Q\in \fc(P)$, set $\ago_P^Q=\ago_{M_P}^{M_Q}$; we denote by an exponent $*$ the dual of $\ago_P^Q$. Thus, we have a decomposition
$$\ago_P=\ago_P^Q\oplus \ago_Q$$
and every $\xi\in \ago_P$ can be written
\begin{equation}
  \label{eq:decompoParab}
  \xi=\xi^Q+\xi_Q
\end{equation}
according to this decomposition.
\end{paragr}

\begin{paragr}[Roots and Coroots.] --- \label{S:deltachapeau}
 Let $B\in \fc(T)$ be a Borel subgroup of $G$. Let $\Delta_B$ be the set of simple roots of $T$ in $B$. Let $\Delta_B^\vee$ be the set of simple coroots. The sets $\Delta_B$ and $\Delta_B^\vee$ form respective bases for the spaces $(\ago_T^G)^*$ and $\ago_T^G$.

Let $P\in \fc(B)$ and let $\Delta_{P,B}\subset \Delta_B$ be the subset consisting of the roots in $N_P$. Let $\Delta_{P,B}^\vee\subset \Delta_B^\vee$ be the corresponding subset of coroots. The restriction map $X^*(T)\to X^*(A_P)$ induces a bijection from $\Delta_{P,B}$ onto a set denoted $\Delta_P$. Similarly, the projection $\ago_T\to \ago_P$ induces a bijection from $\Delta_{P,B}^\vee$ onto a set denoted $\Delta_P^\vee$. Since any Borel subgroup contained in $P$ and containing $T$ belongs to the orbit of $B$ under the Weyl group $W^{M_P}(T)$, the sets $\Delta_P$ and $\Delta_P^\vee$ do not depend on the choice of $B$. Moreover, $\Delta_P$ and $\Delta_P^\vee$ form respective bases for $\ago_P^*$ and $\ago_P$. Let $\hat{\Delta}_P$ be the basis of $\ago_P^*$ dual to $\Delta_P^\vee$.

\begin{lemme} \label{lem:deltachapeau}
Let $Q\in \fc(P)$. Then we have the natural inclusion $\hat{\Delta}_Q\subset\hat{\Delta}_P$.
\end{lemme}

\begin{preuve}
Let $B\subset P$ be a Borel subgroup containing $T$. We have the inclusion $\Delta_{Q,B}^\vee\subset \Delta_{P,B}^\vee$ and the complement $\Delta_{P,B}^\vee-\Delta_{Q,B}^\vee$ is contained in $\ago^Q_T$. Let $\Delta_{Q,P}^\vee$ be the image of  $\Delta_{Q,B}^\vee$ under the projection $\ago_T\to \ago_P$. Then the complement $\Delta_{P}^\vee-\Delta_{Q,P}^\vee$ is contained in $\ago^Q_P$. The projection $\ago_P\to \ago_Q$ induces a bijection from $\Delta_{Q,P}^\vee$ onto $\Delta_Q^\vee$.

Let $\varpi\in \hat{\Delta}_Q\subset \ago_Q^*$. Let $\al^\vee\in \Delta_Q^\vee$ be such that $\varpi(\al^\vee)=1$ and $\varpi$ vanishes on $\Delta_Q^\vee-\{\al^\vee\}$. Let $\beta^\vee\in \Delta_P^\vee$. Then we have the following alternative. Either $\beta^\vee\notin \Delta_{Q,P}^\vee$ and then $\varpi(\beta^\vee)=0$, or $\beta^\vee\in \Delta_{Q,P}^\vee$ and then $\varpi(\beta^\vee)$ equals $1$ if $\beta^\vee$ is the unique element of $\Delta_{Q,P}^\vee$ and $0$ otherwise. This shows that $\varpi\in \hat{\Delta}_P$.

\end{preuve}

\end{paragr}

\begin{paragr}[Cones in $\ago_P$.] --- Let $P\in \fc$. Let $^+\ago_P$ be the \emph{obtuse} open cone in $\ago_P$ defined by
$$^+\ago_P=\{H\in \ago_P \ |\ \varpi(H)>0 \ \forall \varpi\in \hat{\Delta}_P \}.$$
Let $\overline{^+\ago_P}$ be the closure of $^+\ago_P$ in $\ago_T$, i.e.,
$$\overline{^+\ago_P}=\{H\in \ago_P \ |\ \varpi(H)\geq 0 \ \forall \varpi\in \hat{\Delta}_P \}.$$

\begin{lemme} \label{lem:cones} Let $P\subsetneq G$ be a semi-standard parabolic subgroup. We have the following assertions:
  \begin{enumerate}
  \item for all $Q\in \fc(P)$, we have $^+\ago_P \subset \, ^+\ago_Q + \ago_P^Q$;
  \item the cone $^+\ago_P$ is the intersection of the cones $^+\ago_Q + \ago_P^Q$ as $Q$ runs over the maximal proper parabolic subgroups of $G$ containing $P$.
  \end{enumerate}
These assertions also hold for the closed cones.
\end{lemme}

\begin{preuve}
    Let $Q\in \fc(P)$. It is clear that we have
$$^+\ago_Q + \ago_P^Q=\{H\in \ago_P \ |\ \varpi(H)>0 \ \forall \varpi\in \hat{\Delta}_Q \}.$$
The first assertion then follows from the inclusion $\hat{\Delta}_Q\subset\hat{\Delta}_P$ given by Lemma \ref{lem:deltachapeau}. To show the second assertion, it suffices, by the first assertion, to prove that the intersection of the cones $^+\ago_Q + \ago_P^Q$ as $Q$ runs over the maximal elements of $\fc(P)$ is contained in $^+\ago_P$. Let $H$ be in this intersection. We need to show that $\varpi(H)>0$ for all $\varpi\in \hat{\Delta}_P$. But this is clear since for every $\varpi\in \hat{\Delta}_P$ there exists a unique proper maximal parabolic subgroup $Q$ containing $P$ such that $\hat{\Delta}_Q=\{\varpi\}$.
  \end{preuve}

\end{paragr}

\begin{paragr}[Degree of a Reduction.] --- Let $K$ be an extension of $k$, $m\in \mc(K)$ a Hitchin triple, $P\in \fc$ a semi-standard parabolic subgroup and $m_P=(\ec,\theta,t)$ a reduction of $m$ to $P$ (cf. Definition \ref{def:reduc}). The \emph{degree} of the reduction $m_P$ is the unique element $\deg(m_P)\in \ago_P$ that satisfies, for every character $\mu\in X^*(P)$, the following equality
$$ \mu(\deg(m_P))=\deg(\mu(\ec))$$
where $\mu(\ec)$ is the line bundle on $C_K$ obtained by pushing forward the torsor $\ec$ via the representation
$$\mu \ : \ P \to \mathrm{GL}(1)$$
and $\deg(\mu(\ec))$ is its degree.

\begin{lemme}\label{lem:projdeg}
  Let $P\in \fc$ and let $m_P$ be a reduction of $m$ to $P$. Let $Q\in \fc(P)$ and $m_Q$ a reduction of $m$ to $Q$. Then the degree $\deg(m_Q)$ is equal to the projection of $\deg(m_P)$ onto $\ago_Q$ with respect to the decomposition $\ago_P=\ago_Q\oplus \ago_P^Q$.
\end{lemme}

\begin{preuve}
Recall that the reduction $m_P$ is a certain triple $(\ec_P,\theta_P,t)$ (cf. Definition \ref{def:reduc}). Let $\ec_Q$ be defined by $\ec_P\times_k^P Q$ and let $\theta_Q$ be the section of $\Ad_D(\ec_Q)$ obtained by composing $\theta_P$ with the natural inclusion $\Ad_D(\ec_P)\hookrightarrow \Ad_D(\ec_Q)$. One verifies that the triple $(\ec_Q,\theta_Q,t)$ is a reduction of $m$ to $Q$, so it is isomorphic to $m_Q$ (cf. Proposition \ref{prop:Preduc}). Let $\mu\in X^*(Q)$. Clearly, $\mu(\ec_Q)$ is isomorphic to $\mu(\ec_P)$. Hence,
$$\mu(\deg(m_P))=\deg(\mu(\ec_P))=\deg(\mu(\ec_Q))=\mu(\deg(m_Q)).$$
It follows that $\deg(m_P)-\deg(m_Q)$ belongs to $\ago_P^Q$, which proves the lemma.

\end{preuve}

\end{paragr}

\begin{paragr}[Convex Sets Associated to a Hitchin Triple.] --- \label{S:lescvxes} We continue with the notations of the previous paragraph. Let $M$ be the unique semi-standard Levi subgroup such that $f(m)$ belongs to $\Ac_{M,\el}(K)$ (cf. Proposition \ref{prop:reunionAM}). For every $P$ in the set $\fc(M)$ of parabolic subgroups of $G$ containing $M$, let $m_P$ be the unique reduction of $m$ to $P$ (cf. Proposition \ref{prop:Preduc}) and set
$$ \,^+\mathcal{C}_{P,m}=-\deg(m_{P}) - \, ^+\ago_P +\ago^P_T$$
to be the open cylinder in $\ago_T$ with base the obtuse chamber $-\deg(m_{P}) - \, ^+\ago_P$, and
$$ \overline{\,^+\mathcal{C}_{P,m}}=-\deg(m_{P}) - \, \overline{^+\ago_P} +\ago^P_T$$
its closure in $\ago_T$. Note that, for $P=G$, we have $\,^+\mathcal{C}_{G,m}=\ago_T$.

Let  $\mathcal{C}_m$ be the open polyhedron contained in $\ago_T$ defined by
$$\mathcal{C}_m = \bigcap_{P\in \fc(M)}  \,^+\mathcal{C}_{P,m} $$
and let
$$ \overline{\mathcal{C}}_m=\bigcap_{P\in \fc(M)} \overline{\,^+\mathcal{C}_{P,m}}$$
be its closure.
\medskip

If $M=G$, then $\fc(G)=\{G\}$ and $\mathcal{C}_m= \overline{\mathcal{C}}_m=\ago_T$.

\begin{remarques}
We will see an equivalent definition of the preceding polyhedra in Proposition \ref{prop:envpecvx}. These definitions are directly suggested by Arthur's definition of the weights that appear in the definition of weighted orbital integrals (cf. \cite{Art76} \S\S 2,3). The link with Arthur's constructions will become more evident once we have an adelic description of the Hitchin triples (cf. \S\ref{S:description-adelic}). When $M=T$, one recovers the ``complementary polyhedron'' that Behrend associates to the group scheme $\Aut_G(\ec)$ and the torus $T$ (cf. \cite{Beh95}\S6).
\end{remarques}

\begin{proposition}\label{prop:descripCm} With the above notations, in the definition of $\mathcal{C}_m$ or $\overline{\mathcal{C}}_m$, one may replace $\fc(M)$ by $\pc(M)$. If $M\subsetneq G$, one may replace $\fc(M)$ by the maximal elements of $\fc(M)-\{G\}$.
\end{proposition}

\begin{preuve} This is an immediate consequence of the following lemma.
\end{preuve}

\begin{lemme} Let $P\in \fc(M)$ be a semi-standard parabolic subgroup. Then we have the following assertions:
  \begin{enumerate}
  \item For all $Q\in \fc(P)$, we have $ \,^+\mathcal{C}_{P,m} \subset \,^+\mathcal{C}_{Q,m}$;
  \item The cone $ \,^+\mathcal{C}_{P,m}$ is the intersection of the cones $\,^+\mathcal{C}_{Q,m}$ as $Q$ runs over the maximal proper parabolic subgroups of $G$ containing $P$.
  \end{enumerate}
These assertions also hold for the closed cones.

\end{lemme}

\begin{preuve}
Let $Q\in \fc(P)$. Lemma \ref{lem:projdeg} implies that we have the equality
 $$\,^+\mathcal{C}_{Q,m}=-\deg(m_{P})- \, ^+\ago_Q +\ago_T^Q.$$
The lemma then follows immediately from Lemma \ref{lem:cones}.
\end{preuve}
\end{paragr}

\begin{paragr}[Semi-continuity of the Convex Sets.] ---
Let $S$ be an affine $k$-scheme and $m\in \mc(S)$. For every $s\in S$, let $m_s\in \mc(k(s))$ be the Hitchin triple over $k(s)$, the residue field at $s$, deduced from $m$ by base change. Let $\mathcal{C}_m$, respectively $\overline{\mathcal{C}}_m$, be the function on $S$ which associates to each $s\in S$ the open convex $\mathcal{C}_{m_s}$, respectively the closed convex $\overline{\mathcal{C}}_{m_s}$. As one of the referees pointed out to us, the following proposition and its proof are similar to Proposition 7.1.3 of Behrend's thesis (cf. \cite{Beh}).

\begin{proposition} \label{prop:scontinue} Suppose $S$ is noetherian. The function $\mathcal{C}_m$ from $\mc(k(s))$ into the set of subsets of $\ago_T$, ordered by inclusion, is lower semicontinuous, that is, for any subset $\Xi$ of $\ago_T$, the set
$$\{s\in S \ | \Xi \subset \mathcal{C}_{m_s} \}$$
is open. The same holds for the function $\overline{\mathcal{C}}_m$.
\end{proposition}

\begin{preuve}
Let $m=(\ec,\theta,t)\in \mc(S)$. We treat only the function $\mathcal{C}_m$ as the same proof applies to $\overline{\mathcal{C}}_m$.

First, we show that $\mathcal{C}_m$ is constructible and takes only a finite number of values on $S$. For this, after possibly replacing $S$ by an irreducible component, we may assume that $S$ is irreducible. Let $\eta$ be the generic point of $S$ and $C_\eta=C\times_k k(\eta)$. Let $m_\eta=(\ec_{\eta},\theta_{\eta},t_\eta)$ be the Hitchin triple on $C_{\eta}$ deduced from $m$ by the base change $\Spec(k(\eta))\to S$. Let $M\in \lc$ be such that $f(m_\eta)\in \chi_G^M(\Ac_{M,\el}(k(\eta)))$ (cf. Proposition \ref{prop:reunionAM}). Let $P\in \fc(M)$ and let $m_{P,\eta}=(\ec_{P,\eta},\theta_{\eta},t_\eta)$ be ``the'' reduction of $m_{\eta}$ to $P$ (cf. Proposition \ref{prop:Preduc}). To the $P$-reduction $\ec_{P,\eta}$ described in \S\ref{S:reduc}, one associates a section $C_\eta \to \ec_{\eta}/P$. This section extends to all of $C_S$ minus a certain closed subset. The image of this closed subset by the proper morphism $C_S\to S$ is a closed subset which does not contain the generic point $\eta$. Let $U$ be the complement of this closed subset in $S$. By construction, the section $C_\eta \to \ec_{\eta}/P$ extends to a section $C_U\to \ec_{\eta}/P$. Using the bijection in \S\ref{S:reduc}, one deduces from this section a $P$-torsor $\ec_{P,U}$ which is an extension of $\ec_{P,\eta}$ and is a reduction of $\ec_U$ to $P$. It follows that the triple $(\ec_{P,U},\theta_{U},t_{|U})$ is a reduction of $m$ to $P$ over $U$. Hence, the function that associates to $u\in U$ the degree of $m_{P,u}$ is constant. After possibly shrinking $U$, the above function is constant for every $P\in \fc(M)$. After further reducing $U$, we may assume that $f(m_u)\in \chi_G^M(\Ac_{M,\el}(k(u)))$ for all $u\in U$. But then the function $\mathcal{C}_m$ is constant on $U$. By noetherian induction, it takes only finitely many values and there exists a partition of $S$ into constructible sets on which $\mathcal{C}_m$ is constant.
\medskip

Next, we show that for every $\Xi\subset \ago_T$, the set
  $$\{s\in S \ | \Xi \subset \mathcal{C}_{m_s} \}$$
is open. By the previous discussion, it is constructible. It suffices to show that it is stable under generization. For this question, we may assume that $S$ is the spectrum of a discrete valuation ring. Let $s$ be the special point and $\eta$ the generic point of $S$. It suffices to show the inclusion
 \begin{equation}
   \label{eq:incluCm}
   \mathcal{C}_{m_s}\subset \mathcal{C}_{m_\eta}.
 \end{equation}
 Let $(a,t)=f(m)\in \Ac_G(S)$. Let $M\in \lc$ be such that $(a_\eta,t_\eta)$ belongs to $\chi_G^M(\Ac_{M,\el}(k(\eta)))$ (cf. Proposition \ref{prop:reunionAM}). Since the morphism $\chi_G^M$ is a closed immersion (cf. Proposition \ref{prop:immersionfermee}), there exists a pair $(a_M,t)\in \Ac_M(S)$ whose image in $\Ac_G(S)$ is $(a,t)$.
If $M=G$, the inclusion (\ref{eq:incluCm}) is trivial since then $\mathcal{C}_{m_\eta}=\ago_T$. So assume $M\subsetneq G$. Using Proposition \ref{prop:descripCm}, we see that it suffices to show that for every maximal $P\in \fc(M)$, we have
 $$ \,^+\mathcal{C}_{P,m_s}\subset  \,^+\mathcal{C}_{P,m_\eta}$$
which amounts to showing that
$$\varpi(\deg(m_{P,s})-\deg(m_{P,\eta}))\geq 0$$
where $\varpi$ is the unique element of $\hat{\Delta}_P$. Let
$$\rho \ : \ G \to V$$
be the irreducible finite-dimensional representation with highest weight $\varpi$. Let $V_\varpi\subset V$ be the line generated by a highest weight vector of weight $\varpi$. In particular, this line is stable under the action of $P$.

Let $U$ be the open subset of $C_S$ defined as the inverse image via $a_M$ of the $G$-regular open set $\car^{\reg}_M$. In \S\ref{S:reductouvert} we constructed a triple $(\ec_{P,U},\theta_U,t_U)$ consisting of a reduction to $P$ of the triple deduced from $m$ by base change to $U$. This yields a section $U \to \ec/P$ which, by the projectivity of $\ec/P$, extends to an open subset $U_1\supset U$ containing all codimension $\leq 1$ points of $C_S$. Note that $U_1$ contains $C_\eta$ as well as the generic point of $C_s$. We then obtain a $P$-torsor $\ec_{P,U_1}$ which extends $\ec_{P,U}$ (cf. \S\ref{S:reduc}). One verifies that the triple $(\ec_{P,U_1},\theta_{U_1},t_{U_1})$ is a reduction of $m$ to $P$ over $U_1$.

The line bundle $\ec_{P,U_1}\times_k^{P,\rho} V_\varpi$ is a locally split subbundle of the restriction to $U_1$ of the vector bundle $\rho(\ec)=\ec\times_k^{P,\rho} V$. It extends uniquely to a line bundle $\vc_\varpi$ which is a subsheaf of $\rho(\ec)$. However, it may happen that outside $U_1$ this bundle is not locally split. Since $U_1$ contains $C_\eta$, we have the equality
\begin{equation}
  \label{eq:egVeta}
\deg( \vc_{\varpi,\eta})=\varpi(\deg(m_{P,\eta})).
\end{equation}

On the open set $U_1\cap C_s$ of $C_s$, by restricting the triple $(\ec_{P,U_1},\theta_{U_1},t_{U_1})$ we obtain a reduction of $m_s$ to $P$. By the same arguments, this reduction extends to $C_s$. We then obtain a triple $(\ec_{P,s},\theta_{s},t_s)$. The line bundle
$$  \vc_{\varpi,s}'=\ec_{P,s}\times_k^{P,\rho} V_\varpi$$
is a locally split subbundle of $\rho(\ec_s)$ and, as above, we have
\begin{equation}
  \label{eq:egVs}
\deg( \vc_{\varpi,s}')=\varpi(\deg(m_{P,s})).
\end{equation}
But this bundle $\vc_{\varpi,s}'$ coincides with $\vc_{\varpi,s}$ on the open set $U_1\cap C_s$. It follows that $\vc_{\varpi,s}'$ is the saturation of $\vc_{\varpi,s}$ in $\rho(\ec_s)$. Hence,
$$\deg(\vc_{\varpi,s}')\geq \deg(\vc_{\varpi,s})$$
with equality if and only if $\vc_{\varpi,s}'=\vc_{\varpi,s}$. Combining (\ref{eq:egVeta}) and (\ref{eq:egVs}) yields
$$\varpi(\deg(m_{P,s}))=\deg(\vc_{\varpi,s}')\geq \deg(\vc_{\varpi,s}) =\varpi(\deg(m_{P,\eta})),$$
as required.
\end{preuve}

\end{paragr}

\begin{paragr}[Reductions to Adjacent Parabolic Subgroups.] --- Let $K$ be an extension of $k$. Before stating the main result of this paragraph, we illustrate it in the case of $GL(2)$ endowed with its standard maximal torus $T$. Let $m=(\ec,\theta,t)$ be a Hitchin triple in $\mc_{GL(2)}(K)$. The data of $(\ec,\theta)$ is equivalent to that of a rank $2$ vector bundle $\vc$ over $C_K$, together with a twisted endomorphism
$$\theta \ : \ \vc\to\vc(D)$$
(denoted again by $\theta$). Assume that $f(m)$ is not elliptic. It follows that at the generic point $\eta$ of $C_K$, the endomorphism $\theta$ is diagonal in a basis denoted $(e_1,e_2)$ of the vector space $\vc_\eta$. Let $\vc_i\subset \vc$ be the line subbundle determined by $e_i$. The point $t\in \tgo(K)\simeq K^2$ writes as a pair $t=(t_1,t_2)$. There is an indeterminacy in numbering the bundles $\vc_i$ which is removed by requiring that at the point $\infty$ the endomorphism $\theta$ has eigenvalue $t_i$ on the space $\vc_{i,\infty}$.

Let $B\subset GL(2)$ be the standard Borel subgroup and $\overline{B}$ its opposite. It follows from Proposition \ref{prop:Preduc} that there exist reductions $m_B$ and $m_{\bar{B}}$ of $m$ to $B$ and $\overline{B}$ respectively. We identify $\ago_T$ with $\RR^2$ in the obvious way. One verifies the following formulas:
$$\deg(m_B)=(\deg(\vc_1), \deg(\vc/\vc_1))$$
and
$$\deg(m_{\bar{B}})=(\deg(\vc/\vc_2), \deg(\vc_2)).$$
It follows that we have
\begin{equation}
  \label{eq:GL2a}
  \deg(m_{\bar{B}})-\deg(m_B)= \mathrm{length}(\vc/(\vc_1+\vc_2))(1,-1)
\end{equation}
where the length, denoted $\mathrm{length}$, of the torsion module $\vc/(\vc_1+\vc_2)$ is given by
$$\deg(\vc)-\deg(\vc_1)-\deg(\vc_2)$$
since the morphism between line bundles
$$\vc_1 \to \vc/\vc_2$$
is injective.

At the generic point of $C_K$, we have a direct sum
$$\vc_\eta=\vc_{1,\eta}\oplus \vc_{2,\eta}.$$
Let $p_i$ denote the projection onto $\vc_{i,\eta}$. This projection commutes with $\theta$, which acts on the factor $\vc_{i,\eta}$ by the global section $\lambda_i$ of $\oc(D)$. Let $p_1(\vc)$ be the invertible sheaf image of $\vc$ under $p_1$. The projection $p_1$ induces an isomorphism from $\vc/\vc_2$ onto $p_1(\vc)$. The restriction of $p_1$ to $\vc_1$ is injective and induces an isomorphism from $\vc_1$ onto a subsheaf of $p_1(\vc)$. Similarly, one introduces the bundle $p_2(\vc)$ and obtains isomorphisms
$$p_1(\vc)/\vc_1\simeq \vc /(\vc_1+\vc_2) \simeq p_2(\vc)/\vc_2$$
compatible with the twisted endomorphisms induced by $\theta$. Now, $\theta$ acts by the scalars $\lambda_1$ on $p_1(\vc)$ and $\lambda_2$ on $p_2(\vc)$. It follows that the scalar $\lambda_1-\lambda_2$, which is zero since $\theta$ is regular, acts by $0$ on $p_1(\vc)/\vc_1$, hence
$$(\lambda_1-\lambda_2) p_1(\vc) \subset \vc_1(D),$$
from which one deduces the inequality
$$\mathrm{length}(p_1(\vc)/\vc_1) \leq 2\deg(D),$$
or equivalently,
\begin{equation}
  \label{eq:GL2b}
\mathrm{length}(\vc/(\vc_1+\vc_2))\leq 2\deg(D).
\end{equation}

\medskip

It is these formulas (\ref{eq:GL2a}) and (\ref{eq:GL2b}) that we shall generalize to an arbitrary group. For any Levi subgroup $M\in \lc$, two parabolic subgroups $P$ and $Q$ in $\pc(M)$ are said to be \emph{adjacent} if the intersection
$$\Delta_P^\vee \cap (-\Delta_Q^\vee)$$
consists of a single element.

\begin{lemme}\label{lem:pos} Let $m\in \mc(K)$ and $M\in \lc$ such that $f(m)\in \Ac_{M,\el}(K)$.

Let $P$ and $Q$ be two adjacent parabolic subgroups in $\pc(M)$. Let $\al^\vee$ be the unique coroot satisfying
$$\Delta_P^\vee \cap (-\Delta_Q^\vee)=\{\al^\vee\}.$$

Let $m_P$ and $m_Q$ be the respective reductions of $m$ to $P$ and $Q$ (cf. Proposition \ref{prop:Preduc}).

There exists a unique rational number $x_\al$ such that
    \begin{equation}
      \label{eq:positiv}
      -\deg(m_P)+\deg(m_Q) = x_\al \al^\vee.
    \end{equation}
Furthermore, $x_\al$ satisfies the inequality
\begin{equation}
  \label{eq:positiv2}
  0 \leq x_\al \leq  2\mu(\al^\vee)^{-1}\Big(\dim(\ngo_P/(\ngo_P\cap \ngo_Q))\Big)^2 \dim(\qgo/(\ngo_P\cap \ngo_Q)) \deg(D),
\end{equation}
where the character $\mu\in X^*(M)$ is given by $\det\big(\Ad(\cdot)_{|\ngo_P/(\ngo_P\cap \ngo_Q)}\big)$.

In particular, $x_\al$ is bounded on $\mc(K)$.
\end{lemme}

\begin{remarques}
  The left-hand inequality in (\ref{eq:positiv2}) is reminiscent of Lemma 3.6 in \cite{Art76}. In the context of reducing reductive group schemes to a Borel subgroup, it appears in \cite{Beh95} Proposition 6.6.
\end{remarques}

\medskip

\begin{preuve}
Let $m=(\ec,\theta,t)$ be the given point in $\mc(K)$. Let $m_P=(\ec_P,\theta,t)$ and $m_Q=(\ec_Q,\theta,t)$ be its reductions to $P$ and $Q$. The statement is trivial if $M=G$. Thus, we assume $M\subsetneq G$ henceforth.

Let $R \in \fc(M)$ be defined by $\Delta_R^\vee=\Delta_P^\vee-\{\al^\vee\}$. Thus, $\Delta_R^\vee \cup \{-\al^\vee\}=\Delta_Q^\vee$. In particular, $R$ contains both $P$ and $Q$. For every character $\mu\in X^*(R)$, we then have
$$\deg(\mu(m_R))=\deg(\mu(m_P))=\deg(\mu(m_Q)).$$
Consequently, $-\deg(m_P)+\deg(m_Q)$ belongs to the space $\ago_M^R$, which is nothing but the line spanned by $\al^\vee$. This establishes the existence of $x_\al$ satisfying (\ref{eq:positiv}).

It remains to verify inequality (\ref{eq:positiv2}). We begin with the case where $R=G$, that is, $M$ is maximal in $G$. In this case, $P\cap Q=M$. To simplify notation, let
$$\vc=\Ad(\ec),$$
$$\vc_1=\ec_P\times_k^{P,\Ad} \ngo,$$
where $\ngo$ is the Lie algebra of the unipotent radical of $P$, and
 $$\vc_2=\Ad(\ec_Q)=\ec_Q\times_k^{Q,\Ad}\qgo.$$
Then $\vc_1$ and $\vc_2$ are two vector subbundles of the vector bundle $\vc$ such that $\vc_1\cap \vc_2=0$. At the generic point of $C_K$, we have a direct sum
$$\vc_\eta=\vc_{1,\eta}\oplus \vc_{2,\eta}.$$
Let $p_i$ be the projection onto $\vc_{i,\eta}$ and let $p_i(\vc)$ be the vector bundle image of $\vc$ under $p_i$.
The section $\theta$ of $\vc(D)$ factors through $\vc_2(D)$. The adjoint action $\ad(\theta)$ of $\theta$ induces twisted endomorphisms $\vc_1\to \vc_1(D)$ and $\vc_2\to \vc_2(D)$. The vector bundles $\vc_1$ and $\vc/\vc_2$ have the same rank equal to $\dim(\ngo)$.
The obvious morphism
$$\vc_1 \to \vc /\vc_2$$
is injective. It follows that
$$\deg( \vc_1)\leq  \deg(\vc/\vc_2)$$
with equality if and only if the morphism is an isomorphism. Let $\mu$ be the character in $X^*(P)=X^*(M)$ given by $\det\big(\Ad_{|\ngo}\big)$. Hence, $\mu(\al^\vee)>0$. Note that $\mu$ is also equal to the character of $M$ given by $\det\big(\Ad_{|\ggo/\qgo}\big)$. It follows that
$$\mu(\deg(m_P))=\deg(\vc_1)\leq \deg(\vc/\vc_2)=\mu(\deg(m_Q))$$
and
$$\mu(\al^\vee) x_\al= \deg(\vc/\vc_2)-\deg(\vc_1)=\mathrm{length}(\tc),$$
where $\tc$ is the torsion sheaf defined by
$$ \tc= \vc/ (\vc_1+\vc_2).$$
Thus, we obtain the positivity in inequality (\ref{eq:positiv2}). It remains to bound the length of $\tc$. For this, we use the isomorphisms
\begin{equation}
  \label{eq:isoT}
  p_1(\vc)/\vc_1\simeq \tc \simeq p_2(\vc)/\vc_2,
\end{equation}
which are compatible with the twisted endomorphisms induced by $\ad(\theta)$. Let $\chi_i$ be the characteristic polynomial of the restriction of $\ad(\theta)$ to $\vc_{i,\eta}$ and let $\lambda$ be the resultant of these two polynomials. In fact, $\lambda$ is even a nonzero section of $\oc(ND)$ where
$$N=\dim(\ngo)\dim(\qgo).$$
By the Cayley--Hamilton theorem, $\chi_1(\ad(\theta))$ induces the trivial twisted endomorphism of $p_1(\vc)/\vc_1$, hence also of $p_2(\vc)/\vc_2$ (via the isomorphism (\ref{eq:isoT})), and likewise for $\chi_2(\ad(\theta))$. Therefore the twisted endomorphism $p_1(\vc)\to p_1(\vc)(ND)$ given by multiplication by $\lambda$ induces the trivial twisted endomorphism of $p_1(\vc)/\vc_1$. Hence,
$$\lambda \, p_1(\vc)\subset \vc_1(ND),$$
and we obtain the bound
$$\mathrm{length}(\tc) \leq 2\dim(\ngo)N \deg(D),$$
which shows the right-hand bound in inequality (\ref{eq:positiv2}) in the case where $M$ is maximal.

If $M$ is not maximal, we use the preceding results in the reductive group $M_R$, which is the unique element of $\lc$ that is a Levi factor of $R$. The Lie algebra of $R$ is $\pgo+\qgo$ and that of $M_R$ is isomorphic to $(\pgo+\qgo)/(\ngo_P\cap\ngo_Q)$.
\end{preuve}
\end{paragr}

\begin{paragr}[Convex Envelopes Associated to a Hitchin Triple.] --- We continue with the notations of the previous paragraph. Let $m\in \mc(K)$ and $M\in \lc$ such that $f(m)\in \Ac_{M,\el}(K)$. For every $P\in \pc(M)$, let $m_P$ be a reduction of $m$ to $P$. Define
$$\mathcal{CV}_m$$
to be the closed convex envelope of the points
$$-\deg(m_P)$$
for $P\in \pc(M)$.

The positivity in Lemma \ref{lem:pos} will allow us to give another description of the convex $\overline{\mathcal{C}}_m$ (for similar results in analogous contexts, cf. \cite{Art76} \S\S 2,3 and \cite{Beh95} \S2).

 \begin{proposition} \label{prop:envpecvx} For every $m\in \mc(K)$, we have the following equality
$$\overline{\mathcal{C}}_m=\mathcal{CV}_m+\ago_T^M+\ago_G.$$
  \end{proposition}

  \begin{preuve}
    It suffices to prove the equality between two sets that are clearly convex and closed subsets of $\ago_T$, invariant under translation by $\ago_T^M+\ago_G$. Let $X_P$ be the projection of $-\deg(m_P)$ onto $\ago_M^G$. It then suffices to prove the following equality between subsets of $\ago_M^G$:
$$\bigcap_{P\in \pc(M)} \Bigl(X_P -\Bigl(\overline{ ^+\ago_P}\cap \ago_M^G\Bigr) \Bigr)=\mathrm{cvx}(X_P),$$
where $\mathrm{cvx}(X_P)$ is the convex envelope of the $X_P$ for $P\in \pc(M)$.

We prove the inclusion $\subset$. Let $H$ be a point in $\ago_M^G$ that does not belong to $\mathrm{cvx}(X_P)$. Then there exists $\lambda\in (\ago_M^G)^*$ such that
$$\lambda(H) > \sup_{P\in \pc(M)} \lambda(X_P).$$

Let $P\in \pc(M)$ be such that $\lambda$ belongs to the cone positively generated by $\hat{\Delta}_P$. The inequality $\lambda(H) > \lambda(X_P)$ implies that there exists $\varpi\in \hat{\Delta}_P$ such that
 $$\varpi(H) > \varpi(X_P).$$
Consequently, $H$ does not belong to $X_P -\Bigl(\overline{ ^+\ago_P}\cap \ago_M^G\Bigr)$.

Conversely, we prove the inclusion $\supset$. It suffices to see that for any parabolic subgroups $P$ and $Q$ in $\pc(M)$, the vector $X_P-X_Q$ belongs to the cone $\overline{^+\ago_P}\cap \ago_M^G$. If $P$ and $Q$ are adjacent, this follows from Lemma \ref{lem:pos}. More generally, by taking a minimal-length sequence $P_i\in \pc(M)$ for $i=0,\ldots,n$ with $P_{i+1}$ and $P_i$ adjacent, $P_0=P$ and $P_n=Q$, one sees that $X_P-X_Q$
can be written as a positive linear combination of roots of $A_M$ in $N_P$, which yields the desired result.
  \end{preuve}

\end{paragr}
\section{The $\xi$-stability}\label{sec:xi-stab}

\begin{paragr}[Hitchin $\xi$-stable stacks.] --- Let $\xi\in \ago_T$.

\begin{definition} Let $S$ be an affine $k$-scheme and $m\in \mc(S)$. We say that $m$ is $\xi$-stable, resp. $\xi$-semi-stable, if for every $s\in S$ we have
$$\xi\in \mathcal{C}_{m_s},$$
resp.
$$\xi\in \overline{\mathcal{C}}_{m_s}.$$
\end{definition}

\begin{remarques}
 Since $ \mathcal{C}_{m_s}$ is invariant under translation by $\ago_G$, the $\xi$-(semi-)stability depends only on the projection of $\xi$ onto $\ago_T^G$.

In the case $G=GL(n)$, this definition of $\xi$-stability is reminiscent of a weighted stability notion due to Esteves \cite{Est01} in the setting of compactified Jacobians. When $\xi=0$, one recovers the usual notion of stability for Hitchin bundles. Note also that this definition is very close to that used for (ordinary or Hitchin) bundles with parabolic structure studied by Boden and Yokogawa in \cite{BY96} and Heinloth and Schmitt in \cite{HS}. It is also inspired by the truncations and weights that Arthur introduced in the context of the trace formula. Our presentation (in particular the use of the convex $\mathcal{C}_{m}$) is also inspired by Behrend's work on the stability of reductive group schemes in \cite{Beh95}.
\end{remarques}

\begin{definition}\label{def:mcxi}
   Let $\mc^{\xi}$, resp. $\overline{\mc^{\xi}}$, be the substack of $\mc$ such that for every affine $k$-scheme $S$, the category $\mc^{\xi}(S)$, resp. $\overline{\mc^{\xi}}(S)$, is the full subcategory of $\mc(S)$ whose objects are the $\xi$-stable, resp. $\xi$-semi-stable, points.
 \end{definition}

 \begin{definition}\label{def:posgenerale}
We say that $\xi\in \ago_T$ is \emph{in general position} if for every $P\in \fc$ with $P\subsetneq G$, the projection of $\xi$ onto $\ago_P$, with respect to the decomposition $\ago_T=\ago_T^P\oplus \ago_P$, does not belong to the lattice $X_*(P)$.
 \end{definition}

\begin{remarque}
   It follows from Definition \ref{def:mcxi} that $\mc^{\xi}\subset \overline{\mc^{\xi}}$. If $\xi$ is in general position, then for any extension $K$ of $k$ and any Hitchin triple $m\in \mc(K)$, the point $\xi$ cannot lie on the boundary of $\overline{\mathcal{C}}_{m_s}$. Hence, for all $\xi$ in general position we have
$$\mc^{\xi}= \overline{\mc^{\xi}}.$$
 \end{remarque}

 \begin{proposition}\label{prop:schamp-ouvert}
   The stacks $\mc^{\xi}$ and $\overline{\mc^{\xi}}$ are open substacks of $\mc$.
 \end{proposition}

 \begin{preuve}
   It suffices to verify that for every affine $k$-scheme $S$ and every $m\in \mc(S)$ the set of $s\in S$ such that $m_s$ belongs to $\mc^{\xi}$, resp. $\overline{\mc^{\xi}}$, is open. Since $\mc$ is locally of finite type, we may assume that $S$ is noetherian, in which case the result follows from the semicontinuity of the function $\mathcal{C}_{m}$, resp. $\overline{\mathcal{C}}_{m}$ (cf. Proposition \ref{prop:scontinue}).
 \end{preuve}

 \begin{proposition}\label{prop:schptf}
If $G$ is semisimple, then the stacks $\mc^{\xi}$ and $\overline{\mc^{\xi}}$ are stacks of finite type over $k$.
 \end{proposition}

 \begin{preuve}
   Let $B\subset G$ be a Borel subgroup containing $T$. Let $\tc_G$ be the stack of $G$-torsors over $C$, and for every cocharacter $\delta\in X_*(T)$, let $\tc_B^\delta$ be the stack of $B$-torsors of degree $\delta$ over $C$. It is well known that the latter is a stack of finite type over $k$. There is a morphism $\tc_B^\delta \to \tc_G$ which sends a $B$-torsor $\ec$ to the $G$-torsor $\ec\times^B G$. Moreover, the morphism $\mc\to\tc_G$ which sends a triple $(\ec,\theta,t)$ to the $G$-torsor $\ec$ is of finite type. It therefore suffices to show that the image of $\mc^\xi$ or $\overline{\mc^{\xi}}$ under this morphism is of finite type. But this image is covered, by the following lemma, by the images of the $\tc_B^\delta$ for $\delta$ ranging over a finite set. Hence it is indeed of finite type.
 \end{preuve}

 \begin{lemme} Let $B\subset G$ be a Borel subgroup containing $T$. There exists a finite set $\mathcal{X}\subset X_*(T)$ such that for every triple $(\ec,\theta,t)\in \overline{\mc^{\xi}}(k)$ there exists a reduction $\ec_B$ of $\ec$ to $B$ satisfying
$$\deg(\ec_B)\in \mathcal{X}.$$
 \end{lemme}

 \begin{preuve}
According to Harder (cf. \cite{Har69} Satz 2.1.1 and \cite{Har74} p.253), there exists a constant $c>0$ such that every $G$-torsor $\ec$ over $C$ admits a reduction $\ec_B$ to $B$ whose degree $\deg(\ec_B)$ belongs to the acute cone
$$\cgo=\{H\in \ago_T \ | \ \al(H)\geq -c \ \forall \al\in \Delta_B \},$$
where $\Delta_B$ is the set of simple roots of $T$ in $B$. Let $N$ be an integer that bounds, for every root of $T$ in $B$, the sum of its coefficients in the basis $\Delta_B$. Let $d$ be an integer satisfying
\begin{equation}
  \label{eq:inegsurd}
  d> \deg(D) +2N c.
\end{equation}

Let $P$ be a parabolic subgroup of $G$ containing $B$ and let $\Delta_P\subset \Delta_B$ be the subset of simple roots in $N_P$. Let $\cgo_P$ be the cone in $\ago_T$ defined by
$$\cgo_P=\{H\in \ago_T \ | \ \al(H)\geq d \ \forall \al\in \Delta_P \text{ and } \al(H)\leq d \ \forall \al\in \Delta_B- \Delta_P \}.$$
As $P$ runs over the standard parabolic subgroups, the cones $\cgo_P$ cover $\ago_T$. It thus suffices to prove the following: there exists a finite set $\mathcal{X}_P\subset X_*(T)$ such that for every triple $(\ec,\theta,t)\in \overline{\mc^{\xi}}(k)$ and every reduction $\ec_B$ of $\ec$ to $B$ whose degree satisfies
$$\deg(\ec_B)\in \cgo\cap \cgo_P,$$
we have $\deg(\ec_B)\in \mathcal{X}_P$.

\begin{lemme}
  Let $(\ec,\theta,t)\in \overline{\mc^{\xi}}(k)$ and let $\ec_B$ be a reduction of $\ec$ to $B$ whose degree satisfies
$$\deg(\ec_B)\in \cgo\cap \cgo_P.$$
Then $\theta$ factors through $\ec_B\times^{B,\Ad} \pgo_D\hookrightarrow \Ad_D(\ec)$.
\end{lemme}

\begin{preuve}
  The adjoint action of $\theta$ induces a morphism of Lie algebra schemes
$$\ec\times^{G,\Ad} \ggo\to  \Ad_D(\ec)=\ec \times^{G,\Ad}\ggo_D.$$
It suffices to prove that this morphism sends $\ec_B\times^B \bgo$ into $\ec_B\times^B \pgo_D$. Consider a filtration of $\ggo$
$$(0)=\bgo_0 \subset \bgo_1 \subset \ldots \subset \bgo_r=\bgo \subset \pgo=\pgo_s \subset \pgo_{s-1} \subset \ldots \subset \pgo_0=\ggo$$
by $B$-stable subspaces so that the successive quotients $\bgo_{i+1}/ \bgo_{i}$ and  $\pgo_{i}/ \pgo_{i+1}$ are $1$-dimensional.

We will show that for all $0\leq i\leq r$ and $1\leq j\leq s$, the adjoint action of $\theta$ sends $\ec_B\times^{B,\Ad}\bgo_{i}$ into $\ec_B\times^{B,\Ad}\pgo_{j,D}$. The case $i=r$ and $j=s$ gives the desired result. We proceed by induction. The case $i=0$, which is trivial, starts the induction. For a given pair $(i,j)$, the induction hypothesis is that the assertion holds for every pair $(i',j')$ with either $i'<i$ and $j'\leq s$, or $i'=i$ and $j'\leq j$. First, assume $j<s$. We prove the assertion for the pair $(i,j+1)$.
By the induction hypothesis, $\theta$ sends $\ec_B\times^{B,\Ad}\bgo_{i-1}$ into $\ec_B\times^{B,\Ad}\pgo_{D}$ and $\ec_B\times^{B,\Ad}\bgo_{i}$ into $\ec_B\times^{B,\Ad}\pgo_{j,D}$. In particular, $\theta$ induces a morphism of line bundles
\begin{equation}
  \label{eq:morphfibdroites}
   \ec_B\times^{B,\Ad}\bgo_i/ \bgo_{i-1} \to \ec_B\times^{B,\Ad}\pgo_{j,D}/ \pgo_{j+1,D}.
 \end{equation}
 If this morphism is nonzero, then the degree of the target must be greater than the degree of the source. There exist characters, say $\al\in \Phi^B_T\cup \{0\}$ and $\beta\in \Phi^{N_P}_T$, such that $T$ acts on the quotients $\bgo_i/ \bgo_{i-1}$ and $\pgo_{j}/ \pgo_{j+1}$ by $\al$ and $-\beta$, respectively. It follows that the degree of the target is $-\beta(\deg(\ec_B))+\deg(D)$ and that of the source is $\al(\deg(\ec_B))$. Hence we must have
$$  (\al+\beta)(\deg(\ec_B)) \leq \deg(D).$$
But $\beta$ is a root of $T$ in $N_P$, so it has a nonzero component on at least one element of $\Delta_P$. Since $\deg(\ec_B)\in \cgo\cap\cgo_P$, we have the lower bound
$$\beta(\deg(\ec_B))\geq d -Nc.$$
Trivially, $\al(\deg(\ec_B))\geq -Nc$, and thus we contradict inequality (\ref{eq:inegsurd}). Consequently, the morphism (\ref{eq:morphfibdroites}) must vanish, which establishes the assertion for the pair $(i,j+1)$. When $j=s$, one must prove that the assertion holds for the pair $(i+1,1)$, which is demonstrated similarly.

\end{preuve}

Let $\ec_P=\ec_B\times^{B} P$ be the $P$-torsor induced from $\ec_B$. By the preceding lemma, $\theta$ factors through $\Ad_D(\ec_P)$. \emph{A priori}, one only knows that $\chi_G(t)$ equals the value of $\chi_{G}(\theta)$ at the point $\infty$, but there is no reason for $\chi_P(t)$ to equal the value of $\chi_{P}(\theta)$ at $\infty$. However, there exists an element $w$ of the Weyl group $W$ such that $\chi_P(w\cdot t)$ equals the value of $\chi_{P}(\theta)$ at $\infty$. It follows that the triple $(\ec_P,\theta,w\cdot t)$ is a reduction of $(\ec,\theta,w\cdot t)$ to $P$. Clearly, by the definitions, $(\ec,\theta,w\cdot t)$ belongs to $\overline{\mc^{w\cdot\xi}}$. Replacing $\xi$ by $w\cdot\xi$ if necessary, we may assume henceforth that $w=1$.

If $P=G$, then the cone $\cgo\cap \cgo_G$ is compact and $\deg(\ec_B)$ belongs to the finite set $X_*(T)\cap \cgo\cap \cgo_G$. If $P\neq G$, then by the $\xi$-semi-stability of $(\ec,\theta,t)$ we have the condition
$$\xi \in -\deg(\ec_P)- \, ^+\ago_P +\ago^P_T.$$
Thus, $\deg(\ec_B)$ belongs to the compact convex
$$\cgo\cap \cgo_P \cap (-\xi - \, ^+\ago_P +\ago^P_T),$$
and hence to a finite set.

 \end{preuve}

\end{paragr}

\begin{paragr}[The Main Theorem.] --- We introduce the following definition.

 \begin{definition}
   Let $f^\xi$, resp. $\overline{f^\xi}$, denote the restriction of the Hitchin morphism $f$ to $\mc^{\xi}$, resp. $\overline{\mc^{\xi}}$.
 \end{definition}

We may now state our main result.

 \begin{theoreme}
   \label{thm:principal}
Let $\xi\in \ago_T$ be in general position in the sense of Definition \ref{def:posgenerale}. Assume that $G$ is semisimple. Then the stack $\mc^\xi$ is a Deligne-Mumford stack, smooth and of finite type over $k$. Moreover, the Hitchin morphism $f^\xi$ is proper.
 \end{theoreme}

 \begin{preuve}
   By Proposition \ref{prop:schamp-ouvert}, $\mc^\xi$ is an open substack of $\mc$. Moreover, $\mc$ is smooth over $k$ (Biswas and Ramanan proved this by a deformation calculation in \cite{BR94}, which was later recapitulated by Ngô; see \cite{Ngo06} Proposition 5.3 and \cite{Ngo08}), hence $\mc^\xi$ is smooth. We have seen in Proposition \ref{prop:schptf} that the stack $\mc^\xi$ is of finite type. Later we will show that $\mc^\xi$ is a Deligne-Mumford stack (see the proof of Theorem \ref{thm:DM}). The properness of $f^\xi$ follows from the valuative criterion combining the results of Theorems \ref{thm:existence} and \ref{thm:separation}, which will be stated later.
 \end{preuve}

\end{paragr}

\begin{paragr}[The $\xi$-Harder-Narasimhan points.] ---  \label{S:HN} We equip the space $\ago_T$ with an inner product $\bg\cdot,\cdot\bd$ invariant under the Weyl group $W^G(T)$, which turns $\ago_T$ into a Euclidean space. Then $\ago_T$ is canonically isomorphic to its dual $\ago_T^*$. Under this isomorphism, a root and its coroot differ only by a strictly positive scalar. Let $\|\cdot\|$ be the Euclidean norm on $\ago_T$ associated to this inner product. Note that the spaces $\ago_M$ and $\ago_T^M$ (for $M\in \lc$) are then orthogonal.

Let $K$ be an extension of $k$ and let $\xi\in \ago_T$.

\begin{definition} Let $m\in \mc(K)$. The $\xi$-Harder-Narasimhan point of $m$ is the unique point $\varrho_m\in \overline{\mathcal{C}}_m$ which satisfies
 $$\|\varrho_m-\xi\|=\min_{X\in \overline{\mathcal{C}}_m} \|X-\xi\|.$$
\end{definition}

\begin{remarque} Since $\overline{\mathcal{C}}_m$ is a nonempty closed convex set (by Proposition \ref{prop:envpecvx}), there exists a unique $\xi$-Harder-Narasimhan point.
\end{remarque}
\medskip
For every parabolic subgroup $P\in \fc$, let $\ago_P^{+}$ be the \emph{acute} open cone in $\ago_P$ defined by
\begin{equation}
  \label{eq:aigu}
  \ago_P^{+}=\{H\in \ago_P \ |\ \al(H)>0 \ \forall \al\in \Delta_P \}.
\end{equation}
Thus, we obtain a partition of $\ago_T$
\begin{equation}
  \label{eq:partition}
  \ago_T=\bigcup_{P\in \fc}\ago_P^{+}.
\end{equation}

The following proposition and its proof are essentially a reformulation of \cite{Beh95} Proposition 3.13.

\begin{proposition} \label{prop:HN} Let $m\in \mc(K)$ and let $M\in \lc$ be such that $f(m)\in \Ac_{M,\el}$. For every $P\in \fc(M)$, let $m_P$ be a reduction of $m$ to $P$. Let $\varrho\in \overline{\mathcal{C}}_m$. The following assertions are equivalent:
  \begin{enumerate}
  \item The point $\varrho$ is the $\xi$-Harder-Narasimhan point of $m$;
 \item There exists a parabolic subgroup $Q\in \fc(M)$ such that
    \begin{enumerate}
    \item $\xi-\varrho\in \ago_Q^+\cap \ago_T^G$;
    \item The projection of $\varrho$ onto $\ago_M^G$ belongs to the projection onto $\ago_M^G$ of the convex envelope of the points
$$-\deg(m_P)$$
for all $P\in \pc(M)$ with $P\subset Q$;
\end{enumerate}
\item There exists a parabolic subgroup $Q\in \fc(M)$ such that
    \begin{enumerate}
    \item $\xi-\varrho\in \ago_Q^+\cap \ago_T^G$;
    \item $\varrho$ belongs to the affine subspace
$$-\deg(m_Q)+\ago_T^Q+\ago_G.$$
    \end{enumerate}
  \end{enumerate}

\end{proposition}
  \begin{preuve} We show that (1) implies (2). Suppose that $\varrho$ is the $\xi$-Harder-Narasimhan point of $m$. Let $Q$ be the unique element of $\fc$ such that $\xi-\varrho\in \ago_Q^+$, cf. the partition (\ref{eq:partition}). Since $\overline{\mathcal{C}}_m$ is invariant under translation by $\ago_T^M+\ago_G$, we have
\begin{equation}
      \label{eq:1ereinclusion}
      \xi-\varrho \in \ago_M^G.
    \end{equation}
It follows that $Q$ must contain $M$. Thus, (2)(a) is verified. We now check (2)(b). Let $P_0\in \pc(M)$ be such that $P_0\subset Q$. By Proposition \ref{prop:envpecvx}, the point $-\deg(m_{P_0})$ belongs to $\overline{\mathcal{C}}_m$. By convexity of $\overline{\mathcal{C}}_m$, for every $0\leq \lambda \leq 1$, the point $(1-\lambda)\varrho -\lambda\deg(m_{P_0})$ also belongs to $\overline{\mathcal{C}}_m$, and the function
$$\| (1-\lambda)\varrho -\lambda\deg(m_{P_0})-\xi\|^2$$
attains its minimum at $\lambda=0$. Differentiating at $\lambda=0$, one obtains the condition
 \begin{equation}
   \label{eq:derivation}
   \bg \varrho+\deg(m_{P_0}), \xi-\varrho\bd \geq 0.
 \end{equation}
Since $\xi-\varrho\in \ago_Q^+\cap \ago_T^G$, we may write
\begin{equation}
  \label{eq:Yvarpi}
\xi-\varrho =\sum_{\varpi \in \hat{\Delta}_Q} y_\varpi \,\varpi
\end{equation}
with $y_\varpi>0$. Condition (\ref{eq:derivation}) then implies
\begin{equation}
  \label{eq:ineg}\sum_{\varpi \in \hat{\Delta}_Q} y_\varpi \,\varpi( \varrho+\deg(m_{P_0}))\geq 0.
\end{equation}
For every $\varpi\in \hat{\Delta}_Q$, we have
$$\varpi( \varrho+\deg(m_{P_0}))= \varpi( \varrho+\deg(m_Q)),$$
and this quantity is negative since $\varrho\in \overline{\mathcal{C}}_m$. The inequality (\ref{eq:ineg}) is therefore possible only if for every $\varpi\in \hat{\Delta}_Q$ we have
\begin{equation}
  \label{eq:0}
\varpi( \varrho+\deg(m_{P_0}))=0.
\end{equation}
For every $P\in \pc(M)$, let $\lambda_P\in \RR$ be such that
\begin{itemize}
\item $0\leq \lambda_P \leq 1$;
\item $\sum_{P\in \pc(M)}\lambda_P=1$;
\item the projection of $\varrho$ onto $\ago_M^G$ equals the projection onto $\ago_M^G$ of the point
$$\sum_{P\in \pc(M)} -\lambda_P \deg(m_P).$$
\end{itemize}
Such real numbers $\lambda_P$ exist by Proposition \ref{prop:envpecvx}. It remains to show that one may choose $\lambda_P$ so that $\lambda_P=0$ if $P\not\subset Q$. By (\ref{eq:0}), we obtain
\begin{equation}
  \label{eq:=}
  \varpi(\deg(m_{P_0}))= \sum_{P\in \pc(M)} \lambda_P \,  \varpi(\deg(m_P))
\end{equation}
for every $\varpi\in \hat{\Delta}_Q$. Lemma \ref{lem:pos} implies that for $P\in \pc(M)$ the difference $\deg(m_P)-\deg(m_Q)$ is a positive linear combination of elements of $\Delta^\vee_P$. Moreover, $\varpi$ takes only positive values on  $\Delta^\vee_P$ and vanishes on $\,\Delta^\vee_P\setminus \Delta^\vee_Q$. Let $P\in \pc(M)$ be adjacent to $P_0$. Then either $\deg(m_P)=\deg(m_{P_0})$, in which case we may take $\lambda_P=0$, or
$$ \deg(m_P)-\deg(m_{P_0})=x_\al \al^\vee$$
with $x_\al>0$ and $\{\al^\vee\}=\Delta_{P_0}^\vee \cap (-\Delta_{P}^\vee)$ (by Lemma \ref{lem:pos}). Equation (\ref{eq:=}) is then possible only if for every $\varpi\in \hat{\Delta}_Q$ we have $\al^\vee\notin \Delta_{Q}^\vee$. But then $P\subset Q$. By induction on the cardinality of $\Delta_{P_0}^\vee \cap (-\Delta_{P}^\vee)$, one sees that for every $P\in \pc(M)$, either $P\subset Q$ or we may assume $\lambda_P=0$. This proves (2).

\medskip
Clearly, (2) implies (3) since for $P\subset Q$ the projection of $\deg(m_P)$ onto $\ago_Q$ equals $\deg(m_Q)$.

\medskip

Finally, we show that (3) implies (1). Suppose there exists a parabolic subgroup $Q\in \fc(M)$ such that conditions (3)(a) and (3)(b) hold. We must show that (1) holds, i.e., that $\varrho$ is the $\xi$-Harder-Narasimhan point of $m$. Let $X\in \overline{\mathcal{C}}_m$. We need to show that $\|X-\xi\|\geq \|\varrho-\xi \|$. We have
\begin{eqnarray}\label{eq:carres}
  \|X-\xi\|^2&=& \|\varrho-\xi \|^2 +2 \bg \xi-\varrho  , \varrho-X  \bd +\|X-\varrho \|^2\\ \nonumber
  &\geq &  \|\varrho-\xi \|^2 +2 \bg \xi-\varrho  , \varrho-X  \bd\\ \nonumber
&\geq&  \|\varrho-\xi \|^2 + \sum_{\varpi \in \hat{\Delta}_Q} 2y_\varpi \,\varpi(\varrho-X).
\end{eqnarray}
where the $y_\varpi$, defined in (\ref{eq:Yvarpi}), are positive. By hypothesis, $\varrho\in -\deg(m_Q)+\ago_T^Q+\ago_G$, and since $X\in \overline{\mathcal{C}}_m$, in particular
$$X\in -\deg(m_Q)-\overline{^+\ago_Q}+\ago_T^Q+\ago_G.$$
Hence,
$$\varrho-X\in  \overline{^+\!\ago_Q}+\ago_T^Q+\ago_G.$$
Therefore, for every $\varpi \in \hat{\Delta}_Q$, we have $\varpi(\varrho-X)\geq 0$, and (\ref{eq:carres}) gives $\|X-\xi\| \geq \|\varrho-\xi \|$, as desired.
\end{preuve}

The preceding proposition immediately implies the following corollary.

\begin{corollaire} \label{cor:HN} Under the hypotheses of Proposition \ref{prop:HN}, let $\varrho\in \overline{\mathcal{C}}_m$ be the $\xi$-Harder-Narasimhan point of $m$. Then there exists a unique $Q\in \fc(M)$ such that $\varrho$ realizes the distance from $\xi$ to the affine subspace $-\deg(m_Q)+\ago_T^Q+\ago_G$.
\end{corollaire}

\end{paragr}

\section{Adelic Description of Hitchin Fibers}\label{sec:description}

\begin{paragr}[The Set $V$.] --- \label{par:adel1}
Let $V$ be a finite set of closed points of $C$ which contains the point $\infty$ and the support of the divisor $D$. The effective divisor $D$ can then be written as the formal sum
$$D=\sum_{v\in V} d_v\, v$$
where, for every $v\in V$, the integer $d_v$ is positive and is zero outside the support of $D$ (thus, in particular, zero at $\infty$). Let $C^V$ be the open subset of $C$ complementary to $V$. Let $k[C^V]$ be the algebra of regular functions on $C^V$.
\end{paragr}

\begin{paragr}[Notations.] --- \label{par:adel2}
For any $k$-algebra $A$, set $C_A=C\times_k \Spec(A)$ and $C_A^V=C^V\times_k \Spec(A)$. Let $A[C^V]=k[C^V]\otimes_k A$. For every $v\in V$, the completion of the quasi-coherent algebra $\oc_{C_A}$ along the divisor $v\times_\kappa A$ is identified with the ring $A[[z_v]]$ via the choice of a uniformizer $z_v$. Let $A((z_v))$ be the ring of formal Laurent series in the variable $z_v$ with coefficients in $A$: it is the localization of $A[[z_v]]$ obtained by inverting $z_v$. We have the following commutative diagram, where the morphisms are the obvious ones:
  \begin{equation}
    \label{eq:diagdescript}
    \xymatrix{
      \Spec(A((z_v))) \ar[d]_{i_{A,v}'}  \ar[r]^{j_{A,v}} & \Spec(A[[z_v]])   \ar[d]^{i_{A,v}} \\
      C_A^V  \ar[r]^{j_A} & C_A
    }
  \end{equation}

For every $k$-scheme $S$, we denote by $S((z_v))$, respectively $S[[z_v]]$, the functor that, to every $k$-algebra $A$, associates the set of points $S(A((z_v)))$, respectively $S(A[[z_v]])$.
\end{paragr}

\begin{paragr}[Formal Descent \`{a} la Beauville-Laszlo and Uniformization.] ---\label{S:BL}
With the notations of the previous paragraphs, we can state the following proposition.

 \begin{proposition} \label{prop:BL}
Let $(X,(g_v)_{v\in V},t)$ be a triple satisfying the following conditions:
   \begin{enumerate}
   \item $X$ is an element of $\ggo(A[C^V])$;
   \item for every $v\in V$, the element $g_v$ belongs to $G((z_v))(A)$ and satisfies
   $$\Ad(g_v^{-1})X\in z_v^{-d_v}\ggo[[z_v]](A);$$
   in particular, $\chi(X)\in \car[[z_\infty]](A)$;
   \item $t$ is an element of $\tgo^{\reg}(A)$ whose characteristic $\chi(t)$ is equal to the reduction modulo $z_\infty$ of $\chi(X)$.
   \end{enumerate}
There exists a triple $m=(\ec,\theta,t)\in \mc(\Spec(A))$ as well as $G$-equivariant isomorphisms
$$\al_{A,v}\ :\ i_{A,v}^*\ec \to G\times_k A[[z_v]]$$
and
$$\beta_A \ :\ j_A^*\ec\to    G\times_k A[C^V]$$
such that
\begin{itemize}
\item the $G$-equivariant automorphism of $G\times_k A((z_v))$ defined by
$$ (i_{A,v}^{'*}) (\beta_{A})\circ (j^{*}_{A,v})(\al_{A,v}^{-1}) $$
is given by left translation by $g_v\in G((z_v))(A)$;
\item $\beta_A$ induces an isomorphism between $\Ad(j_A^*\ec)$ and $\ggo\times_k A[C^V]$ which sends $j_A^*\theta$ to $X$.
\end{itemize}
Moreover, the triple $(m,(\al_{A,v})_{v\in V},\beta_A)$ is uniquely determined up to a unique isomorphism. Every such triple arises from a unique triple $(X,(g_v)_{v\in V},t)$ as above.
  \end{proposition}

  \begin{preuve}
    This is a consequence of the formal descent result of Beauville-Laszlo (cf. \cite{BL95}, where the case of vector bundles is treated).
  \end{preuve}

\begin{sloppypar}
Let us now introduce the category $Cat_V(A)$ whose objects are the triples $(X,(g_v G[[z_v]](A))_{v\in V},t)$ that satisfy conditions 1 to 3 of Proposition \ref{prop:BL} (note that condition 2 on $g_v$ is clearly stable under translation by elements of $G[[z_v]](A)$) and where the set of morphisms between two objects $(X,(g_vG[[z_v]](A))_{v\in V},t)$ and $(X',(g'_vG[[z_v]](A))_{v\in V},t')$ is given by
\begin{itemize}
\item the empty set if $t\not= t'$;
\item the set of $\delta\in G(A[C^V])$ such that $\Ad(\delta)X=X'$ and for every $v\in V$
$$\delta g_v G[[z_v]](A)= g'_vG[[z_v]](A).$$
\end{itemize}
\end{sloppypar}

  \begin{corollaire}\label{cor:BL}
The construction of Proposition \ref{prop:BL} induces an equivalence of categories between
    \begin{itemize}
    \item the category $Cat_V(A)$;
    \item the full subcategory of $\mc(\Spec(A))$ consisting of triples $(\ec,\theta,t)\in \mc(\Spec(A))$ for which the $G$-torsor $\ec$ is trivial on $C_A^V$ and on $\Spec(A[[z_v]])$ for every $v\in V$.
    \end{itemize}
  \end{corollaire}
\end{paragr}

\begin{paragr}[Reductions.] --- \label{S:reducadel}
Let $K$ be an extension of $k$. Let $(a,t)\in \Ac_{G}(K)$. Let $t_a\in \tgo^{\reg}[[z_\infty]](K)$ be the lift of $t$ such that $a$ and $\chi(t_a)$ coincide on $\Spec(K[[z_\infty]])$ (cf. Proposition \ref{prop:caractdesAM}). Let $M\in \lc$ be such that $(a,t)\in \chi^M_G(\Ac_{M,\el}(K))$ (cf. Proposition \ref{prop:reunionAM}). Let $m=(\ec,\theta,t)\in \mc(K)$ be such that  $f(m)=(a,t)$. Let $\al_{K_v}$ and $\beta_K$ be trivializations of $\ec$ respectively on $\Spec(K[[z_v]])$ and on $C_K^V$. We assume that these trivializations satisfy the conditions of Proposition \ref{prop:BL}. Let $(X,(g_v)_{v\in V},t)$ be the triple deduced by Proposition \ref{prop:BL}. We further assume that $X$ is an element of $\mgo(K[C^V])$ such that $X$ and $t_a$ are conjugate under $M((z_\infty))(K)$.

  \begin{proposition} \label{prop:reducadel}
Let $P\in \fc(M)$. For every $v\in V$, let $p_v\in P((z_v))(K)$ be such that
$$g_v\in p_v\, G[[z_v]](K).$$
Then there exists a triple $(\ec_P,\theta_P,t)$ which is a reduction of $m$ to $P$ and trivializations $\al_{K,v}^P$ and $\beta_K^P$ of $\ec_P$ respectively on $\Spec(K[[z_v]])$ and on $C_K^V$ such that
\begin{itemize}
\item the $P$-equivariant automorphism of $G\times_k K((z_v))$ defined by
$$ (i_{K,v}^{'*}) (\beta_{K}^P)\circ (j^{*}_{K,v})(\al_{K,v}^P)^{-1} $$
is given by left translation by $p_v\in P((z_v))(K)$;
\item $\beta_K^P$ induces an isomorphism between $\Ad(j_K^*\ec_P)$ and $\pgo\times_k K[C^V]$ which sends $j_K^*\theta_P$ to $X$.
\end{itemize}
  \end{proposition}

  \begin{preuve}
From the relation $\Ad(g_v^{-1})X\in z_v^{-d_v}\ggo[[z_v]](K)$ and the fact that $X\in \mgo(K[C^V])$, it follows that
   $$\Ad(p_v^{-1})X\in z_v^{-d_v}\pgo[[z_v]](K).$$
In particular, $\chi_P(X)\in \car_M[[z_\infty]](K)$. Since we have assumed that $X$ and $t_a$ are conjugate under $M(F_\infty)$, it follows that $\chi_P(X)=\chi_P(t_a)$. Thus, $(X,(p_v)_{v\in V}, t)$ is a triple satisfying conditions 1 to 3 of Proposition \ref{prop:BL} relative to $P$. Applying Proposition \ref{prop:BL} to the group $P$ yields a triple $(\ec_P,\theta_P,t)$. One verifies that this is indeed a reduction of $m$ to $P$.
  \end{preuve}

\end{paragr}

\begin{paragr}[Functions $H_P$.] --- \label{S:HP}
We continue with the notations of the previous paragraph.

  \begin{definition} \label{def:HP}
Let $P\in \fc(T)$ and $v\in V$, and define
$$H_P\ : \ G((z_v))(K) \to \ago_{P}$$
to be the map satisfying, for every $\mu\in X^*(P)$ and $g\in G((z_v))(K)$,
$$\mu(H_P(g))=-\val_v(\mu(p)),$$
where $\val_v$ is the usual valuation on $K((z_v))$ and $p$ is an element of $P((z_v))(K)$, uniquely defined up to left translation by an element of $P[[z_v]](K)$, such that
$$g\in p\, G(K[[z_v]]).$$
The existence of such an element $p$ is ensured by the Iwasawa decomposition.

More generally, for a family $(g_v)_{v\in V}$ of elements of $G((z_v))(K)$, we set
$$H_P((g_v)_{v\in V})=\sum_{v\in V} H_P(g_v).$$
This definition makes sense when $V$ is infinite provided that $g_v$ belongs to $G[[z_v]]$ for all $v$ outside a finite set. In particular, the function $H_P$ is well defined on the $G$-points valued in the adeles of $F$.
\end{definition}

We now state a corollary to Proposition \ref{prop:reducadel}.

\begin{corollaire}\label{cor:reducadel}
Under the assumptions and notations of \S \ref{S:reducadel} and Proposition \ref{prop:reducadel}, we have the following equality for every $P\in \fc(M)$ and every reduction $m_P$ of $m$ to $P$:
$$\deg(m_P)=H_P\bigl( (g_v)_{v\in V}\bigr).$$
\end{corollaire}

\begin{preuve}
Let $P\in\fc(M)$. Proposition \ref{prop:reducadel} shows that there exists a reduction $(\ec_P,\theta_P,t)$ of $m$ to $P$ such that the $P$-torsor $\ec_P$ is defined by the gluing data $p_v\in P((z_v))(K)$ where, for every $v\in V$, we have $g_v\in p_v\, G[[z_v]]$. Therefore, for every character $\mu\in X^*(P)$, the family $(\mu(p_v))_{v\in V}$ serves as gluing data for the line bundle $\mu(\ec_P)$. One verifies the equality
$$\deg\bigl(\mu(\ec_P)\bigr)=-\sum_{v\in V} \val_v\bigl(\mu(p_v)\bigr)=\mu\Bigl(H_P\bigl( (g_v)_{v\in V}\bigr)\Bigr),$$
which yields the result.
\end{preuve}

\end{paragr}

\begin{paragr}[Adelic Description of Hitchin Fibers.] --- \label{S:description-adelic}
We continue with the notations of the previous paragraphs except that now $V$ denotes the set of \emph{all} closed points of $C$. Let $F$ be the function field of the curve $C$. Let $\xi\in \ago_T$ and $(a,t)\in \Ac(k)$.

\medskip

 \begin{definition}\label{def:Cata}
Let $\mathcal{X}_{(a,t)}$ be the set of pairs $(X,(g_v)_{v\in V})$ satisfying
\begin{enumerate}
\item $X$ is a semisimple, $G$-regular element of $\ggo(F)$ whose characteristic $\chi(X)$ is equal to the restriction of $a$ at the generic point of $C$;
\item for every $v\in V$, the element $g_v$ belongs to $G((z_v))(k)/G[[z_v]](k)$ and satisfies the following three conditions:
  \begin{enumerate}
  \item for all $v\in V$ outside a finite set, $g_v$ is the trivial class;
  \item $\Ad(g_v^{-1})X\in z_v^{-d_v}\ggo[[z_v]](k)$;
  \item the characteristic $\chi(t)$ is equal to the reduction modulo $z_\infty$ of $\chi(X)$.
  \end{enumerate}
\end{enumerate}
  \end{definition}

  \begin{remarque}
The condition 2.(b) implies that the restriction of $\chi(X)$ to $\Spec(k((z_\infty)))$ belongs to $\car[[z_\infty]](k)$. Therefore, condition 2.(c) makes sense.
  \end{remarque}

Let $t_a\in \tgo^{\reg}[[z_\infty]](k)$ be the unique lift of $t$ such that $\chi(t_a)$ equals the restriction of $a$ to $\Spec(k[[z_\infty]])$ (cf. Proposition \ref{prop:caractdesAM}). Let $M\in \lc$ be such that $(a,t)\in \chi^M_G(\Ac_{M,\el}(k))$ (cf. Proposition \ref{prop:reunionAM}).

\begin{definition}\label{def:Cata2}
Let $\mathcal{X}^\xi_{(a,t)}$ be the subset of $\mathcal{X}_{(a,t)}$ consisting of those pairs $(X,(g_v)_{v\in V})$ which satisfy

\begin{enumerate}
\item[(bis)] $X$ is a semisimple, elliptic, $G$-regular element of $\mgo(F)$ that is conjugate, via an element of $M((z_\infty))(k)$, to $t_a$;
\item
\begin{enumerate}
\item[(d)] the projection of $\xi$ onto $\ago_M^G$ belongs to the projection onto $\ago_M^G$ of the convex hull of the points
$$-H_P\bigl((g_v)_{v\in V}\bigr)=-\sum_{v\in V} H_P(g_v).$$
\end{enumerate}
\end{enumerate}
\end{definition}

For every $\delta\in G(F)$ and every pair $(X,(g_v)_{v\in V})$ in $\mathcal{X}_{(a,t)}$, the pair $(\Ad(\delta)X,(\delta g_v)_{v\in V})$ also belongs to $\mathcal{X}_{(a,t)}$. This gives rise to a left action of $G(F)$ on $\mathcal{X}_{(a,t)}$. The induced action of $M(F)$ preserves $\mathcal{X}_{(a,t)}^\xi$. Conditions 1.(bis) and 2.(d) above are preserved by the action of $M(F)$: the first trivially so, and the second because
$$ H_P\bigl((mg_v)_{v\in V}\bigr)=H_P\bigl((g_v)_{v\in V}\bigr)$$
for every $m\in M(F)$. (More precisely, for every $\mu\in X^*(P)$ and every $m\in M(F)$,
$$ \mu\Bigl(H_P\bigl((mg_v)_{v\in V}\bigr)\Bigr)=\mu\Bigl(H_P\bigl((g_v)_{v\in V}\bigr)\Bigr)-\sum_{v\in V} \val_v\bigl(\mu(m)\bigr),$$
and the sum over $V$ is zero by the product formula.)

\medskip

When an abstract group $G$ acts on a set $\mathcal{X}$, we denote by
\begin{equation}
  \label{eq:groupoidequot}
  [G\backslash \mathcal{X}]
\end{equation}
the quotient groupoid whose objects are the elements of $\mathcal{X}$ and for any two objects $x$ and $x'$, the set of morphisms from $x$ to $x'$ is the set of $g\in G$ such that $g\cdot x=x'$.

\begin{proposition}\label{prop:fibre}
The construction of Proposition \ref{prop:BL} induces an equivalence of categories between the quotient groupoid  $[M(F)\backslash \mathcal{X}^\xi_{(a,t)}]$ and the fiber $\overline{f^\xi}^{-1}(a,t)$.
\end{proposition}

\begin{preuve}
Let $m=(\ec,\theta,t)\in \overline{\mc^\xi}(k)$ such that $f(m)=(a,t)$. Since the field $k$ is algebraically closed, the $G$-torsor $\ec$ admits a generic trivialization (cf. Lemma \ref{lem:Sen}). Fix such a trivialization to deduce a generic trivialization of $\Ad_D(\ec)$, hence an element $X\in \ggo(F)$ which is the image of $\theta$ under this trivialization. Let $Y$ be a semisimple, $G$-regular, and elliptic element in $\mgo(F)$ such that
  \begin{itemize}
  \item $\chi_G(Y)$ equals $a_\eta$, the restriction of $a$ at the generic point $\eta$ of $C$;
  \item $Y$ is conjugate to $t_a$ by an element of $M((z_\infty))(k)$.
  \end{itemize}
Such a $Y$ exists (cf. Corollary \ref{cor:caractdesAM}). It follows from Lemma \ref{lem:cjsep} that $X$ and $Y$ are conjugate under $G(F)$. After possibly changing the trivialization of $\ec$, we may (and will) assume that $X=Y$.

A choice of trivializations of $\ec$ on the formal neighborhoods $\Spec(k[[z_v]])$ gives rise, via the construction of Proposition \ref{prop:BL}, to elements $g_v\in G((z_v))$. The class of $g_v$ modulo $G[[z_v]]$ does not depend on this choice. Thus, we obtain a pair $(X,(g_v)_{v\in V})$ in $\mathcal{X}_{(a,t)}$ which furthermore satisfies condition 1.(bis) of Definition \ref{def:Cata}. Using Proposition \ref{prop:reducadel} and the fact that the considered generic trivialization extends to an open subset of $C$, one sees that
$$\deg(m_P)=\sum_{v\in V} H_P(g_v)$$
for every $P\in \pc(M)$. By Proposition \ref{prop:envpecvx}, the $\xi$-semistability condition translates into condition 2.(d) of Definition \ref{def:Cata}. Hence, $(X,(g_v)_{v\in V})$ belongs to $\mathcal{X}^\xi_{(a,t)}$. The reader is invited to verify that this construction indeed gives an equivalence of categories as stated.

\end{preuve}

\end{paragr}

\section{Valuative Criterion: Existence}\label{sec:existence}

\begin{paragr}[Statement.] ---
Let $\xi$ be an \emph{arbitrary} element of $\ago_T$. The goal of this section is to prove the existence part of the valuative criterion for properness for the morphism $\overline{f^\xi}$, which is summarized by the following theorem.

\begin{theoreme}\label{thm:existence}
Let $\kappa$ be an algebraically closed extension of $k$. Let $R$ be a discrete valuation ring, complete and with residue field $\kappa$. Let $K$ be the field of fractions of $R$ and $\overline{K}$ an algebraic closure of $K$.

Let $(a,t)\in \Ac(R)$ and let $m_K\in \overline{\mc^\xi}(K)$ be such that $f(m_K)=(a,t)$.

Then there exists a finite extension $K'\subset \overline{K}$ of $K$ and an element $m\in \overline{\mc^\xi}(R')$, where $R'$ is the integral closure of $R$ in $K'$, such that
\begin{enumerate}
\item the image of $m$ in $\overline{\mc^\xi}(K')$ is isomorphic to that of $m_K$;
\item $f(m)=(a,t)$.
\end{enumerate}
\end{theoreme}

The proof of Theorem \ref{thm:existence} will occupy us until the end of this section.
\end{paragr}

\begin{paragr}[A few words on the proof.] --- In \cite{Lan75} (Section 3, Theorem), Langton showed that the moduli stack of semistable vector bundles on a smooth projective curve satisfies the existence part of the valuative criterion of properness. By adapting Langton's arguments, Nitsure (cf. \cite{Nit91}, Section 6) proved Theorem \ref{thm:existence} in the case $G=GL(n)$ and $\xi=0$. Faltings succeeded in extending this proof to any reductive group $G$ in the case of a base field of characteristic $0$ and for $\xi=0$ (cf. \cite{Fal93}, Theorem II.4). However, his proof does not seem to extend to the case of a base field of positive characteristic. In \cite{Hei08}, Heinloth proved the analogue of Langton's theorem for the moduli stack of $G$-bundles on a smooth projective curve over a base field of characteristic zero or ``not too small''. Our proof partly follows Heinloth's arguments, which themselves were inspired by Langton's.

Let us sketch the broad outline of our proof. The notations are those of Theorem \ref{thm:existence}. Let $C_K$ and $C_\kappa$ be respectively the generic and the special fiber of the curve $C\times_k R$ over $R$. Let $(a,t)\in \Ac(R)$ and let $m_K\in \overline{\mc^\xi}(K)$ be a $\xi$-semi-stable Hitchin triple over $C_K$ of characteristic $(a,t)$. The existence statements that follow sometimes require replacing $K$ by a finite extension, which we shall not recall each time. Such a triple $m_K$ extends to a triple $m_R$ over the whole curve $C_R$ (cf. \S\ref{S:existpourf}). The special fiber $m_\kappa$ of the extension is not necessarily $\xi$-semi-stable. However, one may choose the extension $m_R$ so that the special fiber $m_\kappa$ is as little ``$\xi$-unstable'' as possible, by which we mean that the distance $d$ from $\xi$ to the convex set $\overline{\mathcal{C}}_{m_\kappa}$ (defined in Section \ref{sec:cvx}) is minimal (cf. \S\ref{S:d}). If it is zero, $m_\kappa$ is $\xi$-semi-stable. Otherwise, there exists a maximal parabolic subgroup $Q$ of $G$ such that the face of $\overline{\mathcal{C}}_{m_\kappa}$ associated to $Q$ contains the point at which this distance is attained (cf. \S\ref{S:YHN}). The extension $m_R$ is given by ``gluing data'' (cf. \S\ref{S:trivadm}; the result of \cite{DS95} plays a crucial role there and requires the reduction to the semisimple case of \S\ref{S:reducsemisimple}). By conjugating these data by a suitable $K$-point of the center of $Q$, one obtains another triple $m'_R$ which is isomorphic to $m_K$ on the generic fiber and whose face of $\overline{\mathcal{C}}_{m'_\kappa}$ associated to the parabolic subgroup $Q^-$ opposite to $Q$ is equal to the face of $\overline{\mathcal{C}}_{m_\kappa}$ associated to $Q$. By minimality of $d$, this forces the convex set $\overline{\mathcal{C}}_{m'_\kappa}$ to be degenerate: the faces of $\overline{\mathcal{C}}_{m'_\kappa}$ associated to $Q$ and to $Q^-$ lie in the same affine hyperplane. But this implies that the reduction to $Q$ of $m_\kappa$ lifts to a reduction to $Q$ of the triple $m_{R/(\pi)^2}\in \mc(R/(\pi)^2)$ deduced from $m_R$ by the base change $R\to R/(\pi)^2$, where $(\pi)\subset R$ is the maximal ideal. One may iterate this process by changing the $K$-point of the center of $Q$ (cf. \S\ref{S:cstr-de-triv}). One then obtains a reduction of $m$ to $Q$ over the formal completion $\varinjlim_n C\times_k R/\pi^n$. By a theorem of Grothendieck, this reduction to $Q$ algebraizes, and this then contradicts the $\xi$-stability of $m_K$ (cf. \S\ref{S:lacontradiction}).
\end{paragr}

\begin{paragr}[Notations.] --- \label{S:anneauR}
From now on, we fix the objects $R$, $\kappa$, $K$, and $\overline{K}$ as in the statement of Theorem \ref{thm:existence}. By choosing a uniformizer $\pi$, we identify $R$ with the ring $\kappa[[\pi]]$.
\end{paragr}

\begin{paragr}[Theorem \ref{thm:existence} for the stack $\mc$.] --- \label{S:existpourf}
First, we show that Theorem \ref{thm:existence} holds when we replace the stack $\overline{\mc^\xi}$ by the entire stack $\mc$. This is well known. We give an argument here for the sake of completeness. We recall the assumptions of Theorem \ref{thm:existence} except that we do not assume that the point $m_K\in \mc(K)$ is $\xi$-semistable. The point $m_K$ can be written explicitly as a Hitchin triple $(\ec_K,\theta_K,t)$. Let $\eta$ be the generic point of $C_K=C\times_k K$. After possibly replacing $K$ by a finite extension, we may and do assume that the bundle $\ec_K$ is trivial at the point $\eta$ (cf. Lemma \ref{lem:Sen}). Fix such a trivialization; one then obtains a trivialization of $\Ad_D(\ec)$ at $\eta$, and hence an element $Y\in \ggo(F)$, where $F$ is the function field of $C_K$, which is the image of $\theta$ under this trivialization. Note that $\chi(Y)$ is equal to the restriction $a_\eta$ of $a$ to $\eta$.

Let $B$ be the local ring of the generic point of the special fiber of $C_R=C\times_k \Spec(R)$. This is a discrete valuation ring with field of fractions $F$. By taking a Kostant section (for example), one sees that there exists $X\in \ggo(B)$ such that $\chi(X)=a|_{\Spec(B)}$. Since $\chi(X)=\chi(Y)$, Lemma \ref{lem:cjsep} implies that, after possibly replacing $K$ by a finite extension, we may assume that there exists $g\in G(F)$ such that $\Ad(g)X=Y$. One may then, with this element $g\in G(F)$, glue the bundle $\ec_K$ and the section $\theta$ with the pair $(G\times_k B,X)$ over $\Spec(B)$. Thus one obtains a pair $(\ec_U,\theta_U)$ that extends $(\ec_K,\theta_K)$ to an open subset $U$ of $C_R$ which contains all points of codimension $\leq 1$. But every $G$-bundle over such an open set $U$ extends uniquely to a $G$-bundle over all of $C_R$ (cf. \cite{CS79}, Theorem 6.13). One thus obtains a $G$-bundle $\ec$ over $C_R$ extending $\ec_K$. The bundle $\Ad_D(\ec)$ possesses a section over $U$, namely $\theta_U$. But since $U$ contains all points of codimension $\leq 1$ and $\Ad_D(\ec)$ is affine over $C$, this section extends uniquely over $C_R$ (cf. \cite{Gro67}, Corollary 20.4.12). We have thus obtained a pair $(\ec,\theta)$. The points $\chi(\theta)(\infty_R)$ and $\chi(t)$ are two points of $\car(R)$. They are in fact equal since their images in $\car(K)$ are equal. Thus the triple $(\ec,\theta,t)\in \mc(R)$, hence proving Theorem \ref{thm:existence} for $\mc$.
\end{paragr}

\begin{paragr}[Reduction to the semisimple case.] --- \label{S:reducsemisimple}
In this paragraph, we show that Theorem \ref{thm:existence} for every semisimple group implies Theorem \ref{thm:existence} for every reductive group.

Let $G$ be a connected reductive group over $k$ with Lie algebra $\ggo$ and let $T$ be a maximal torus. Let $G_{\der}$ be the derived subgroup of $G$ and $\ggo_{\der}$ its Lie algebra. The decomposition
$$\ggo=\ggo_{\der} \oplus \zgo$$
induces a decomposition
\begin{equation}
  \label{eq:carG1}
  \car_G=\car_{G_{\der}}\oplus \zgo.
\end{equation}

Let $A_G$ be the connected center of $G$ and $\zgo$ its Lie algebra. Let $G'=G/A_G$. This is a semisimple group with Lie algebra $\ggo'=\ggo/\zgo$ and which admits $T'=T/A_G$ as a maximal torus. The obvious morphism $\ggo \to \ggo'$ induces isomorphisms $\ggo_{\der}\to \ggo'$ and
\begin{equation} \label{eq:carG2}
  \car_{G_{\der}}\simeq \car_{G'}.
\end{equation}
The projection of $\car_G$ onto $\car_{G_{\der}}$ defined by (\ref{eq:carG1}), composed with the isomorphism (\ref{eq:carG2}) and the obvious morphism $\tgo \to \tgo'=\tgo/\zgo$, induces a morphism $\Ac_G \to \Ac_{G'}$. The obvious bijection $\lc^G(T) \to \lc^{G'}(T')$ is compatible with the decomposition of Proposition \ref{prop:reunionAM}.

Let $S$ be an affine $k$-scheme and let $m=(\ec,\theta,t)\in\mc_G(S)$. Let $\ec'$ be the $G'$-bundle obtained by pushing forward $\ec$ via the morphism $G\to G'$. The section $\theta$ pushed forward by the morphism $\ggo\to \ggo'$ yields a section $\theta'$ of $\Ad_D(\ec')$. Let $t'$ be the image of $t$ under the morphism $\tgo\to\tgo'=\tgo/\zgo$. The triple $m'=(\ec',\theta',t')$ belongs to $\mc_{G'}(S)$. Thus, one obtains a morphism $\mc_G\to \mc_{G'}$ which sends $m$ to $m'$. It fits into the commutative diagram
$$\xymatrix{\mc_G \ar[d]^{f_G} \ar[r] & \mc_{G'} \ar[d]^{f_{G'}} \\ \Ac_G  \ar[r] & \Ac_{G'}}.$$

Using the morphism $P \to P/A_G$, it is clear that every reduction of $m$ to a parabolic subgroup $P$ gives naturally a reduction of $m'$ to $P/A_G$. If $S$ is the spectrum of a field, one obtains a natural bijection between the reductions of $m$ and those of $m'$ compatible with the bijection $\fc^G(T) \to \fc^{G'}(T')$. The obvious surjective map $\ago_T \to \ago_{T'}$ induces an isomorphism $\ago_T^G\simeq \ago_{T'}$. It is then evident from the definitions that $m$ is $\xi$-semistable if and only if $m'$ is $\xi^G$-semistable, where $\xi^G$ is the projection of $\xi$ onto $\ago_T^G$.

Returning to the situation of Theorem \ref{thm:existence}, let $(a,t)\in \Ac_G(R)$ and let $m_K=(\ec_K,\theta_K,t)\in \overline{\mc^\xi_G}(K)$ such that $f_G(m_K)=(a,t)$. Let $m_K'=(\ec_K',\theta_K',t')\in \overline{\mc^\xi_{G'}}(K)$ be the image of $m_K$ under the morphism $\mc_G\to \mc_{G'}$. Let $(a',t')=f_{G'}(m'_K)$. The point $(a',t')$ belongs to $\Ac_{G'}(R)$ since it is the image of $(a,t)$. Theorem \ref{thm:existence} for the semisimple group $G'$ implies that, after possibly replacing $K$ by a finite extension $K'$ in $\overline{K}$ and $R$ by its integral closure in $K'$, we may assume that $m_K'$ extends to a triple $m'=(\ec',\theta',t')\in \overline{\mc^\xi_{G'}}(R)$ extending $m_K$. After again possibly replacing $K$ and $R$ as above, we may assume that the $G'$-bundle $\ec'$ is trivial on the spectrum of $B$, the local ring of the point of $C_R$ defined by $(\pi)$ (cf. \cite{DS95}, Theorem 2). Fix such a trivialization. One then obtains a trivialization of $\Ad_D(\ec')$ on $\Spec(B)$; this sends $\theta'$ to a point $Y'\in \ggo'(B)\simeq\ggo_{\der}(B)$ such that $\chi_{G'}(Y')=a'|_{\Spec(B)}$. Note that $a|_{\Spec(B)}=a'|_{\Spec(B)}\oplus Z$ with $Z\in \zgo(B)$ according to (\ref{eq:carG1}) and (\ref{eq:carG2}); hence we obtain a point $Y=Y'+Z\in \ggo(B)$. Any trivialization of $\ec_K$ over an open set $U$ of $C_K$ gives trivializations of $\ec'_K$ and $\Ad_D(\ec'_K)$ over the same open set as well as sections $X$ and $X'$ of $\ggo\times_k U$ and $\ggo'\times_k U$, respectively, deduced from $\theta$ and $\theta'$. The gluing condition for the trivializations of $(\ec',\theta')$ over $\Spec(B)$ and $U$ translates into the existence of an element $g'\in G'(F)$ such that $\Ad(g')X'=Y'$, where $F$ is the function field of $C_R$. Any lift $g\in G(F)$ of $g'$ then satisfies $\Ad(g)X=Y$ by the construction of $Y$. But the obstruction to such a lift lives in $H^1(F,A_G)$. After possibly replacing $K$ by an extension, we may assume that this group is trivial (cf. \cite{Ser94}, Chapter III \S2), and hence such a lift exists. From such a lift, one deduces, as in Paragraph \ref{S:existpourf}, an extension $m=(\ec,\theta,t)\in \mc_G(R)$ of $m_K$. By construction, the image of $(\ec,\theta,t)$ in $\mc_{G'}(R)$ coincides with $m'$ on an open subset of $C_R$ containing all points of codimension $\le 1$. By an argument already mentioned in Paragraph \ref{S:existpourf}, the image of $(\ec,\theta,t)$ is then (isomorphic to) $m'$. It follows that $m \in \overline{\mc^\xi_{G}}(R)$ as desired.
\end{paragr}

\begin{paragr}[The points $m_K$, $(a,t)$ and $(\bar{a},\bar{t})$.] ---
From now on and until the end of the proof of Theorem \ref{thm:existence}, assume that $G$ is \emph{semisimple}. We work under the hypotheses of Theorem \ref{thm:existence}. Let $(a,t)\in \Ac(R)$ and let $m_K\in \overline{\mc^\xi}(K)$ be such that $f(m_K)=(a,t)$. Let $(\bar{a},\bar{t})\in \Ac(\kappa)$ be the point obtained from $(a,t)$ by base change $\Spec(\kappa)\to \Spec(R)$.
\end{paragr}

\begin{paragr}[The real number $d$ and the point $m=(\ec,\theta,t)\in \mc(R)$.] --- \label{S:d}
For every finite extension $K'\subset\overline{K}$ of $K$, let $m_{K'}\in \overline{\mc^\xi}(K')$ be the point obtained from $m_K$ by base change $\Spec(K')\to \Spec(K)$. Let $R'$ be the integral closure of $R$ in $K'$. Let $m_{R'}\in \mc(R')$ be an extension of $m_{K'}$ to $C\times_k R'$. The ring $R'$ is a discrete valuation ring with residue field $\kappa$. Let $m_{\kappa}$ be the point of $\mc(\kappa)$ obtained from $m_{R'}$ by base change $\Spec(\kappa)\to \Spec(R')$. Let
$$d(\overline{\mathcal{C}}_{m_{\kappa}},\xi)$$
denote the distance from $\xi$ to the convex set $\overline{\mathcal{C}}_{m_{\kappa}}$ defined in Section \ref{sec:cvx}.

\begin{lemme}\label{lem:d}
The set of real numbers $d(\overline{\mathcal{C}}_{m_{\kappa}},\xi)$, as $R'$ and $m_{R'}$ vary as above, has a minimum.
\end{lemme}

\begin{preuve}
This set is nonempty (cf. \S \ref{S:existpourf}) and is bounded below by $0$. In fact, it is even discrete and closed in $\ago_T$ since it is contained in the set of distances from $\xi$ to the subspaces $X+\ago_T^Q+\ago_G$ for $X\in X_*(Q)$ and $Q\in \fc$ (cf. Corollary \ref{cor:HN}). It therefore contains its infimum.
\end{preuve}

Let $d\geq 0$ be the smallest element of the set considered in Lemma \ref{lem:d}. After possibly replacing $K$ by a finite extension, we may and do assume that there exists
$$m=(\ec,\theta,t)\in \mc(R)$$
whose image in $\mc(K)$ is the starting point $m_K$ and such that the distance from $\xi$ to the convex set $\overline{\mathcal{C}}_{m_\kappa}$ is $d$, where $m_\kappa\in \mc(\kappa)$ is the image of $m$. From now on, we \emph{fix} such an $m$.

Since $m$ is $\xi$-semistable over $C_K$, it will be $\xi$-semistable if and only if it is also $\xi$-semistable over $C_\kappa$. Thus, we have the equivalences:
$$m\in \overline{\mc^{\xi}}(R) \Leftrightarrow m_\kappa \in \overline{\mc^{\xi}}(\kappa)  \Leftrightarrow d=0.$$
Thus, Theorem \ref{thm:existence} holds for the point $m_K$ if and only if $d=0$. In what follows, we assume that
$$d>0$$
and we will show that this leads to a contradiction.
\end{paragr}

\begin{paragr}[The Levi subgroup $L$ and the parabolic subgroups $Q_0$ and $Q$.] --- \label{S:YHN}
Let $L\in \lc$ be such that $f(m_{\kappa})\in \chi^L_G(\Ac_{L,\el}(\kappa))$ (cf. Proposition \ref{prop:reunionAM}). Let $\varrho$ be the Harder-Narasimhan $\xi$-point of $m_\kappa$. According to Proposition \ref{prop:HN}, there exists a parabolic subgroup $Q_0\in \fc(L)$ such that
\begin{itemize}
    \item $\xi-\varrho\in \ago_{Q_0}^+\cap \ago_T^G$;
    \item $\varrho$ belongs to the affine subspace
$$-\deg(m_{\kappa,Q_0})+\ago_T^{Q_0}.$$
\end{itemize}
where $m_{\kappa,Q_0}$ is a reduction of $m_\kappa$ to $Q_0$ (here $\ago_G=\{0\}$ since $G$ is semisimple). By hypothesis, the distance from $\varrho$ to $\overline{\mathcal{C}}_{m_\kappa}$ is nonzero; hence the parabolic subgroup $Q_0$ is proper. In the remainder, fix $Q\in \fc(Q_0)$ maximal among the proper parabolic subgroups of $G$ that contain $Q_0$. Let $Q^-$ be the parabolic subgroup opposite to $Q$, in the sense that $Q^- \cap Q$ is the unique Levi factor of $Q$ containing $T$.
\end{paragr}

\begin{paragr}[Some notations: $V$, $i_{R,v}$, $j_{R,v}$, etc.] ---
Let $V$ be a finite set of closed points of $C_\kappa$ which contains the support of the divisor
$$D_\kappa=D\times_k \kappa=\sum_{v\in V} d_v \, v,$$
as well as the point $\infty_\kappa$. Let $C_\kappa^V=C_\kappa-V$ and let $\kappa[C^V]$ denote the algebra of regular functions on $C_\kappa^V$. We recall the notations of Section \ref{sec:description}, in particular diagram (\ref{eq:diagdescript}) for the curve $C_\kappa=C\times_k \kappa$, where the morphisms $j_A$, $i_{A,v}$, etc. are defined for every point $v\in V$ and for every $\kappa$-algebra.
\end{paragr}

\begin{paragr}[The point $t_a$.] ---
Let $t_a$ be the unique point of $\tgo^{G\textrm{-}\reg}[[z_\infty]](R)$ which satisfies the following two conditions:
  \begin{itemize}
  \item the reduction modulo $z_\infty$ of $t_a$ is the point $t\in \tgo^{G\textrm{-}\reg}(R)$;
  \item $\chi_G(t_a)=\chi_D(\theta)\circ i_{R,\infty}$.
  \end{itemize}
The existence and uniqueness of $t_a$ follow from the fact that $\chi$ induces an \'{e}tale morphism from $\tgo^{G\textrm{-}\reg}$ onto $\cgo^{G\textrm{-}\reg}$. Let
$$\bar{t}_a\in \tgo^{G\textrm{-}\reg}[[z_\infty]](\kappa)$$
be the reduction modulo $\pi$ of $t_a$. Note that we have $\bar{t}_a=\bar{t}_{\bar{a}}$, where $t_{\bar{a}}$ is the unique point of $\tgo^{G\textrm{-}\reg}[[z_\infty]](\kappa)$ with reduction $\bar{t}$ and characteristic $\bar{a}$.
\end{paragr}

\begin{paragr}[Trivialization of $m$ on an open subset of the special fiber.] ---
Let $m_\kappa=(\ec_\kappa,\theta_\kappa,\bar{t})\in \mc(\kappa)$ be the point obtained from $m$ by base change $\Spec(\kappa)\to \Spec(R)$. We have the following lemma.

\begin{lemme} \label{lem:trivouvert}
After possibly adding a finite number of points to $V$, we are in the following situation. There exists a trivialization of the $G$-bundle $j_\kappa^*(\ec_\kappa)$ over $C_\kappa^V$
$$  \beta_\kappa\ :\ j_\kappa^*(\ec_\kappa)\to G\times_{k} C_\kappa^V,$$
such that the induced isomorphism of group schemes
$$\iota_\kappa \ :\ \Aut_G(j_\kappa^*(\ec_\kappa)) \to G\times_{k}C_\kappa^V$$
satisfies the following condition: the element $d\iota_\kappa(j_\kappa^*\theta)$, where $d\iota_\kappa$ is the derivative of $\iota_\kappa$, satisfies:
\begin{enumerate}
\item $d\iota_\kappa(j_\kappa^*\theta)$ belongs to $\lgo(\kappa[C_\kappa^V])$;
\item $d\iota_\kappa(j_\kappa^*\theta)$ is conjugate to $\bar{t}_\infty$ by an element of $L((z_\infty))(\kappa)$.
\end{enumerate}
\end{lemme}

\begin{preuve}
Let $F$ be the function field of $C_\kappa$. Since $\kappa$ is algebraically closed, the $G$-bundle $\ec_\kappa$ is trivial at the generic point of $C_\kappa$ (cf. Lemma \ref{lem:Sen}). By choosing a trivialization, one also obtains a generic trivialization of $\Ad_D(\ec_\kappa)$ which sends $\theta$ to some element $X\in \ggo(F)$. Let $(\bar{a},\bar{t})=f(m_\kappa)$. This element belongs to $\Ac_{L,\el}(\kappa)$ by the definition of $L$ (cf. \S\ref{S:YHN}). According to Corollary \ref{cor:caractdesAM}, there exist $X'$ semisimple and $G$-regular in $\lgo(F)$ and $l\in L((z_\infty))(\kappa)$ such that
\begin{enumerate}
\item[(A)] $\chi_G(X')=\bar{a}|_{\Spec(F)}$;
\item[(B)] $\Ad(l)X'=\bar{t}_\infty$.
\end{enumerate}
Note that we have $\chi_G(X)=\chi_G(X')$ by the definition of $\bar{a}$. Lemma \ref{lem:cjsep} implies that, after changing the generic trivialization of $\ec_\kappa$, we may assume that $X$ satisfies conditions (A) and (B) above.

We then obtain the lemma since any generic trivialization of $\ec_\kappa$ extends to an open subset of $C_\kappa$ and the element $X$ likewise extends.
\end{preuve}
\end{paragr}

\begin{paragr}[Trivialization of $m$ on formal disks of the special fiber.] ---
Fix a trivialization
\begin{equation}
  \label{eq:betakappa}
  \beta_\kappa\ :\ j_\kappa^*(\ec_\kappa)\to G\times_{k} C_\kappa^V,
\end{equation}
which satisfies the conditions of Lemma \ref{lem:trivouvert}.

For every $v\in V$, let
\begin{equation}
  \label{eq:alphakappav}
  \al_{\kappa,v} \ : \ i_{\kappa,v}^*(\ec_\kappa) \to G\times_k \kappa[[z_v]]
\end{equation}
be a $G$-equivariant isomorphism. Such an isomorphism exists because it amounts to giving a section of the torsor $i_{\kappa,v}^*(\ec_\kappa)$ over $\Spec(\kappa[[z_v]])$. Such a section exists at least over the special point and, by smoothness, extends to $\Spec(\kappa[[z_v]])$.

\begin{lemme} \label{lem:triv_en_v}
One may choose the trivializations $ \al_{\kappa,v}$ so that the triple
\begin{equation}
  \label{eq:tripletadele}
  (\bar{X},(\bar{g}_v)_{v\in V},\bar{t})
\end{equation}
associated to the point $m_\kappa=(\ec_\kappa,\theta_\kappa,\bar{t})\in \mc(\kappa)$ and the trivializations $(\al_{\kappa,v},\beta_\kappa)$ via the bijection of Proposition \ref{prop:BL} satisfies the following conditions:
  \begin{enumerate}
  \item  $\bar{g}_v\in Q((z_v))(\kappa)$;
  \item $\al_{\kappa,\infty}(i_{\kappa,\infty}^*\theta)=\bar{t}_a$;
  \item $\bar{g}_\infty \in L((z_\infty))(\kappa)$.
  \end{enumerate}
\end{lemme}

\begin{preuve}
For any choice of $ \al_{\kappa,v}$, we have
$$\bar{g}_v\in G((z_v))(\kappa).$$
The Iwasawa decomposition
$$G((z_v))(\kappa)=Q((z_v))(\kappa) \cdot G[[z_v]](\kappa)$$
shows that one can always modify $ \al_{\kappa,v}$ so that
\begin{equation}
  \label{eq:gbinQ}
  \bar{g}_v\in Q((z_v))(\kappa),
\end{equation}
which gives assertion 1.

We have
$$\Ad(\bar{g}_\infty)^{-1}\bar{X}\in \ggo[[z_\infty]](\kappa).$$
The element $\chi(\bar{X})$, after base change to $\Spec(k[[z_\infty]])$, belongs to $\car[[z_\infty]](\kappa)$ and its reduction modulo $z_\infty$ coincides with that of $\chi(\bar{t}_\infty)$. Since $\chi$ is \'{e}tale over $\car^{\reg}$, we have
$$\chi(\Ad(\bar{g}_\infty)^{-1}\bar{X})=\chi(\bar{X})=\chi(\bar{t}_\infty).$$
Lemma \ref{lem:section} below shows that one can modify $ \al_{\kappa,\infty}$ so that
\begin{equation}
   \label{eq:cj_t_X}
   \Ad(\bar{g}_\infty)^{-1}\bar{X}=\bar{t}_\infty.
 \end{equation}
But by Lemma \ref{lem:trivouvert}, $\bar{X}$ and $\bar{t}_\infty$ are conjugate under $L((z_\infty))(\kappa)$. Hence, we also have $\bar{g}_\infty\in L((z_\infty))(\kappa)$.
\end{preuve}

\begin{lemme}\label{lem:section}
Let $Y \in \ggo[[z_\infty]](R)$ be such that $\chi_G(Y)=\chi_G(t_a)$. Then there exists $g\in G[[z_\infty]](R)$ such that $\Ad(g)Y=t_a$.\\
The same statement holds when replacing $R$ by $\kappa$ and $t_a$ by $\bar{t}_a$.
\end{lemme}

\begin{preuve}
Let $\tc$ be the $R[[z_\infty]]$-scheme defined for every $R[[z_\infty]]$-algebra $B$ by
$$\tc(B)=\{h\in G(B)\  | \ \Ad(h)(Y)=t_a\}.$$

We show that the scheme $\tc$ is smooth over $\Spec(R[[z_\infty]])$. First, note that the centralizers of $Y$ and $t_a$ in $G\times_k R[[z_\infty]]$ are subgroup schemes that are tori; this follows from the fact that $\theta_\infty\in \tgo^{G\textrm{-}\reg}(R[[z_\infty]])$ and $\chi_G(t_a)=\chi_G(Y)$. Consequently, locally in the \'{e}tale topology, these centralizers are conjugate. Since $Y$ and $t_a$ have the same characteristic, these elements are locally conjugate in the \'{e}tale topology. Thus, the scheme $\tc$ has local sections for the \'{e}tale topology. It is therefore a torsor under the torus $T\times_k R[[z_\infty]]$, and hence smooth.\\
By Lemma \ref{lem:cjsep}, the scheme $\tc\times_{R[[z_\infty]]} \kappa$ has sections. By the smoothness of $\tc$ over $R[[z_\infty]]$ and by the completeness of $R[[z_\infty]]\simeq \kappa[[\pi,z_\infty]]$, this section lifts to a section $h\in G[[z_\infty]](R)$ of $\tc$.

The second statement is proved similarly.
\end{preuve}
\end{paragr}

\begin{paragr}[Admissible trivializations of $(\ec,\theta,t)$.] --- \label{S:trivadm}
From now on, we fix trivializations $(\al_{\kappa,v},\beta_\kappa)$ of $m_\kappa$ that satisfy Lemmas \ref{lem:trivouvert} and \ref{lem:triv_en_v}. Let $C^V_R=C^V_\kappa \times_\kappa R$. Let $R[C^V]=\kappa[C^V]\otimes_\kappa R$.

\begin{definition} \label{def:trivadm}
An \emph{admissible trivialization} of $m=(\ec,\theta,t)$ is the data of $G$-equivariant isomorphisms
\begin{equation}
  \label{eq:betaR}
  \beta_R\ :\ j_R^*\ec\to G\times_{k} R[C^V],
\end{equation}
and for each $v\in V$
\begin{equation}
  \label{eq:alphaRv}
  \al_{R,v} \ : \ i_{R,v}^*\ec \to G\times_k R[[z_v]],
\end{equation}
which satisfy the following conditions:
\begin{enumerate}
\item the isomorphism $\beta_R$ specializes to the isomorphism $\beta_\kappa$ in (\ref{eq:betakappa});
\item the isomorphism $\al_{R,v}$ specializes to the isomorphism $\al_{\kappa,v}$ in (\ref{eq:alphakappav});
\item $\al_{R,\infty}(i_{R,\infty}^*\theta)=t_a$.
\end{enumerate}
\end{definition}

\begin{remarque}
Let $(X,(g_v)_{v\in V},t)$ be the triple associated to $m\in \mc(R)$ and the trivializations (\ref{eq:betaR}) and (\ref{eq:alphaRv}) by Proposition \ref{prop:BL}. Then these trivializations form an admissible trivialization of $m$ if and only if
\begin{itemize}
\item the reduction modulo $\pi$ of the triple $(X,(g_v)_{v\in V},t)$ is the triple $(\bar{X},(\bar{g}_v)_{v\in V},\bar{t})$ defined in (\ref{eq:tripletadele});
\item $\Ad(g_\infty)^{-1}X=t_a$.
\end{itemize}
\end{remarque}

\begin{lemme}
After possibly replacing $K$ by a finite extension and $R$ by its integral closure in that extension, we may assume that there exist admissible trivializations of $(\ec,\theta,t)$.
\end{lemme}

\begin{preuve}
After possibly replacing $K$ by a finite extension, we may and do assume that $j^*_R\ec$ is trivial over $C_R^V$ (since $G$ is semisimple, this is possible by Theorem 3 of \cite{DS95}). Hence, the $G$-bundle $j^*_R\ec$ has sections over $C^V_R$. This yields an isomorphism $\beta_R$ as in (\ref{eq:betaR}). By translating this section by an element of $G(\kappa[C^V])$, we may assume that $\beta_R$ satisfies condition 1 above.\\
For each $v\in V$, by using the smoothness of $\ec$ over $C_R$, one sees that the bundle $i_{R,v}^*\ec$ has sections over $\Spec(R[[z_v]])$. Hence, one obtains isomorphisms $\al_{R,v}$ as in (\ref{eq:alphaRv}). As before, we may assume that $\al_{R,v}$ satisfies condition 2.\\
By Proposition \ref{prop:BL}, one deduces from $\beta_R$ and the $\al_{R,v}$ a triple $(X,(g_v)_{v\in V},t)$ as above. Let
$$Y=\Ad(g_\infty^{-1})X.$$
This is an element of $\ggo[[z_\infty]](R)$ satisfying $\chi_G(Y)=\chi_G(X)=\chi_G(t_a)$. Lemma \ref{lem:section} shows that $Y$ and $t_a$ are conjugate under $G[[z_\infty]](R)$. This allows us to conclude.
\end{preuve}
\end{paragr}

\begin{paragr}[Unipotent subgroups.] ---
Let $\Phi^+ = \Phi^{N_Q}_T$ and $\Phi^- = \Phi^{N_{Q^-}}_T$ be the respective sets of roots of $T$ in the unipotent radicals $N_Q$ and $N_{Q^-}$. The connected component $A_Q$ of the center of $M_Q$ is a torus of dimension $1$. Let
\begin{equation}
    \label{eq:lambda}
    \lambda : \Gm \to A_Q
\end{equation}
be the unique isomorphism that satisfies $\alpha(\lambda)>0$ for every $\alpha\in \Phi^+$. For every integer $i>0$, define
$$\Phi^+_i=\{\alpha\in \Phi^+ \mid \alpha(\lambda)=i\},$$
and set
$$\Phi^-_i=-\Phi^+_i$$
(with additive notation).

For every $\alpha\in \Phi^+\cup\Phi^-$, let $N_\alpha$ be the associated unipotent subgroup and let
\begin{equation}
  \label{eq:zeta}
  \zeta_\alpha \ : \ \mathbb{G}_a \to N_\alpha
\end{equation}
be the one-parameter subgroup associated to $\alpha$. For every integer $i>0$, let $N_i$ be the unipotent subgroup generated by the subgroups $N_\alpha$ such that $\alpha(\lambda)\le -i$. This yields a decreasing sequence of unipotent subgroups
$$ \ldots \subset N_i \subset \ldots \subset N_1 = N_{Q^-},$$
and for $i$ sufficiently large, we have
$$N_i=\{1\}.$$
One verifies the following commutator relation:
$$[N_i,N_j]\subset N_{i+j}.$$
\end{paragr}

\begin{paragr}[Auxiliary condition on $V$.] ---
We begin with the following lemma.

\begin{lemme}\label{lem:Adn}
There exists a nonempty finite set $V_0$ of closed points of $C_\kappa$ such that for every $\Gamma(C_\kappa-V_0,\oc_{C_\kappa})$-algebra $B$ and every integer $i\ge 1$, the map from $N_{i}(B)$ to $\ngo_{i}(B)$ defined by
$$n\in  N_{i}(B)\mapsto \Ad(n^{-1})\bar{X}-\bar{X}$$
is bijective.
\end{lemme}

\begin{preuve}
Let $F$ be the function field of $C\times_k \kappa$. We leave to the reader the task of proving that the morphism $n \mapsto \Ad(n^{-1})\bar{X}-\bar{X}$ induces an isomorphism from $N_{i}\times_k F$ onto $\ngo_{i}\times_k F$. Hence, it extends to an isomorphism on a Zariski open subset of $C_\kappa$, whose complement is a finite set $V_0$. The lemma follows.
\end{preuve}

After enlarging $V$ if necessary, we henceforth assume that $V$ contains the finite set $V_0$ given by Lemma \ref{lem:Adn}.
\end{paragr}

\begin{paragr}[Constructions of certain admissible trivializations.] --- \label{S:cstr-de-triv}
Let $r$ be the least common multiple of the integers $i$ such that $\Phi^+_i\neq 0$. Let $R'=R[\pi^{\frac{1}{r}}]$ and let $K'$ be the field of fractions of $R'$. For every integer $j$, define $z_j\in A_Q(K')$ by
$$z_j=\lambda\Bigl(\pi^{\frac{j}{r}}\Bigr),$$
where $\lambda$ is the isomorphism $\Gm \to A_Q$ defined in (\ref{eq:lambda}).

Let $(\beta_R,(\al_{R,v})_{v\in V})$ be an admissible trivialization of $m=(\ec,\theta,t)$ (cf. Definition \ref{def:trivadm}). Let $(X,(g_v)_{v\in V},t)$ be the triple deduced from $m$ via Proposition \ref{prop:BL}.

\begin{definition}
  \label{def:j}
Define
$$j(\beta_R,(\al_{R,v}))\in \NN\cup \{\infty\}$$
to be the supremum of the set of integers $j\in \NN$ satisfying
\begin{enumerate}
\item $\Ad(z_j)X \in \ggo(R'[C^V])$;
\item for every $v\in V$,
$$z_j g_v z_j^{-1} \in G((z_v))(R').$$
\end{enumerate}
\end{definition}

\begin{proposition}\label{prop:cruciale}
For every integer $j\in \NN$, there exists an admissible trivialization $(\beta_R,(\al_{R,v}))$ of $m=(\ec,\theta,t)$ such that
$$j(\beta_R,(\al_{R,v}))\geq j.$$
\end{proposition}

\begin{preuve}
Assume by contradiction that the proposition fails. Let $j\in \NN$ be the smallest integer such that
$$j\geq j(\beta_R,(\al_{R,v}))$$
for every admissible trivialization $(\beta_R,(\al_{R,v}))$ of $m$. Since we have $\bar{X}\in \lgo(\kappa[C^V])$ and $\bar{g}_v\in Q((z_v))(\kappa)$ (cf. assertion 1 of Lemma \ref{lem:trivouvert} and assertion 1 of Lemma \ref{lem:triv_en_v}), we have
$$j\geq 1.$$

Let $(\beta_R,(\al_{R,v}))$ be an admissible trivialization of $m$ such that
$$j=j(\beta_R,(\al_{R,v})).$$

Let $(X,(g_v)_{v\in V},t)$ be the triple deduced from $m$ via Proposition \ref{prop:BL}.

Let $\val_{\pi}$ denote the valuation on $R$, normalized by $\val_\pi(\pi)=1$. This valuation extends to $R'$ and to $R'[C^V]$ in the obvious way. For every $\alpha\in \Phi^{\pm}$, the $k$-vector space $\ngo_\alpha$ is one-dimensional. For any vector $X_\alpha\in \ngo_\alpha(R'[C^V])$, let $\val_\pi(X_\alpha)$ denote the valuation of the coefficient of $X_\alpha$ with respect to a nonzero vector of $\ngo_{\alpha}(k)$.

According to the decomposition into eigenspaces for $A_Q$
$$(\bigoplus_{\alpha\in \Phi^{\pm}}\ngo_\alpha) \oplus \mgo_Q,$$
write
\begin{equation}
  \label{eq:decompoX}
  X=\sum_{\alpha\in \Phi^\pm}X_\alpha  + X_{M_Q}.
\end{equation}
Since the reduction modulo $\pi$ of $X$ is $\bar{X}\in \lgo(\kappa[C^V])$, we have
\begin{equation}
  \label{eq:valX}
  \val_\pi(X_\alpha)\geq 1
\end{equation}
for every $\alpha\in \Phi^{\pm}$. Moreover, the reduction modulo $\pi$ of $X_{M_Q}$ is $\bar{X}$.

Since the order of the Weyl group is invertible in $k$ (cf. \S2.1), we have the decomposition
$$\Ad(z_j)X=\sum_{\alpha\in \Phi^\pm}\pi^{\frac{j\alpha(\lambda)}{r}} X_\alpha  + X_{M_Q}.$$
By condition 1 in Definition \ref{def:j}, for every $\alpha\in \Phi^{\pm}$,
\begin{equation}
  \label{eq:inegvaluation}
  \frac{j\alpha(\lambda)}{r}+\val_\pi(X_\alpha)\geq 0.
\end{equation}
For $\alpha\in \Phi^+$, this inequality is strict.

\begin{lemme} \label{lem:existstricte}
There exists an admissible trivialization $(\beta_R,(\al_{R,v}))$ of $m$ for which
\begin{itemize}
\item $j=j((\beta_R,(\al_{R,v})))$;
\item the inequality (\ref{eq:inegvaluation}) is strict for every $\alpha\in \Phi^{-}$.
\end{itemize}
\end{lemme}

\begin{preuve}
Start with an admissible trivialization $(\beta_R,(\al_{R,v}))$ of $m$ for which there exists some $\alpha\in \Phi^-$ such that
\begin{equation}
  \label{eq:egal=valXal}
  \frac{j\alpha(\lambda)}{r}+\val_\pi(X_\alpha)=0.
\end{equation}
Let $i$ be the smallest integer such that there exists $\alpha\in \Phi^-_i$ satisfying (\ref{eq:egal=valXal}). Let
$$Y=\Ad(z_j)X\in \ggo(R'[C^V])$$
and let $\Yb$ be the reduction modulo $\pi^{\frac{1}{r}}$ of $Y$. Then $\Yb$ belongs to $\bar{X}+\ngo_i(\kappa[C^V])$. By Lemma \ref{lem:Adn}, there exists $n_i\in N_i(\kappa[C^V])$ such that
\begin{equation}
  \label{eq:XbYb}
  \Ad(n_i^{-1}) \bar{X} = \Yb.
\end{equation}
For every $\alpha\in \Phi_i^-$, choose $u_\alpha\in \kappa[C^V]$ and let $n_{i+1}\in N_{i+1}(\kappa[C^V])$ such that
\begin{equation}
  \label{eq:n}
  n_i=n_{i+1}n,
\end{equation}
with
$$n=\prod_{\alpha\in \Phi_i^-}\zeta_\alpha(u_\alpha).$$
Since $n_{i+1}\in N_{i+1}(\kappa[C^V])$, we have
$$\Ad(n_{i+1}^{-1})\bar{X}\in \bar{X}+\ngo_{i+1}(\kappa[C^V]).$$
Comparing with (\ref{eq:XbYb}) and (\ref{eq:n}), we obtain
\begin{equation}
  \label{eq:AdnYb}
  \Ad(n)Y\in \bar{X}+\ngo_{i+1}(\kappa[C^V]).
\end{equation}
Note that
$$z_j^{-1}nz_j=\prod_{\alpha\in \Phi_i^-}\zeta_\alpha\Bigl( \pi^{ \frac{ji}{r}} u_\alpha\Bigr)$$
belongs to $N_i(R[C^V])$ and its reduction modulo $\pi$ is trivial. Indeed, by assumption, there exists $\alpha\in \Phi_i^-$ for which (\ref{eq:egal=valXal}) holds. Hence, by (\ref{eq:valX}),
$$\frac{ji}{r}=-\frac{j\alpha(\lambda)}{r}=\val_\pi(X_\alpha)\in \NN^*.$$

The triple $(\Ad(z_j^{-1}nz_j)X,(z_j^{-1}nz_jg_v)_{v\in V},t)$ defines, via Proposition \ref{prop:BL}, a trivialization of $m$, say $(\beta_R',(\al'_{R,v}))$, which is still admissible. Clearly by Definition \ref{def:j} we have
$$j(\beta_R',(\al'_{R,v}))\geq j$$
(and by the maximality of $j$ we even have equality). Replacing $(\beta_R,(\al_{R,v}))$ by $(\beta_R',(\al'_{R,v}))$, and using (\ref{eq:AdnYb}), we reduce to the case where $\Yb\in \bar{X}+\ngo_{i+1}(\kappa[C^V])$. By induction, we may even assume that $\Yb=\bar{X}$. But then the inequality (\ref{eq:inegvaluation}) is strict.
\end{preuve}

Now let $(\beta_R,(\al_{R,v}))$ be an admissible trivialization of $m$ satisfying the conclusions of Lemma \ref{lem:existstricte}. Let $(X,(g_v)_{v\in V},t)$ be the triple deduced from $m$ by Proposition \ref{prop:BL}. Thus, we have
\begin{equation}
  \label{eq:firstineg}
   \min_{\alpha\in \Phi^-}\Bigl(-\frac{r}{\alpha(\lambda)}\val_\pi(X_\alpha)\Bigr)>j.
\end{equation}
Note that, by the definition of $r$, the left-hand side of this inequality is an integer.

For each $v\in V$, let $B_v$ be the ring obtained from $R((z_v))$ by localizing at the prime ideal generated by $\pi$. It is a discrete valuation ring with residue field $\kappa((z_v))$. Denote again by $\val_\pi$ the valuation normalized by $\val_\pi(\pi)=1$. The point $g_v\in G((z_v))(R)$ then induces a morphism from the trait $\Spec(B_v)$ into $G$. Since the reduction modulo $\pi$ of $g_v$ is $\bar{g}_v$ and $\bar{g}_v\in Q((z_v))(\kappa)$ (cf. (\ref{eq:gbinQ})), the image of the special point by this morphism lies in the open subset $N_{Q} M_Q N_{Q^-}$. It follows that the image of the entire trait is contained in the open set $N_{Q} M_Q N_{Q^-}$. Therefore, for each $v\in V$, there exist $x_v \in M_Q(B_v)$ and elements $b_{\alpha,v}\in B_v$ for each $\alpha\in \Phi^\pm$ such that
\begin{equation}
  \label{eq:decompogv}
  g_v=\Bigl(\prod_{\alpha\in \Phi^+} \zeta_\alpha(b_{\alpha,v})\Bigr) \cdot x_v \cdot \Bigl(\prod_{\alpha\in \Phi^-} \zeta_\alpha(b_{\alpha,v})\Bigr).
\end{equation}
Let $\xb_v\in M_Q((z_v))(\kappa)$ be the reduction modulo $\pi$ of $x_v$. Since the reduction $\bar{g}_v$ of $g_v$ modulo $\pi$ belongs to $Q((z_v))(\kappa)$, we have $\val_\pi(b_{\alpha,v})\geq 1$ for $\alpha\in \Phi^-$ and
\begin{equation}
  \label{eq:inMN_Q}
 \bar{g}_v\in \xb_v N_Q((z_v))(\kappa).
\end{equation}
We have
\begin{equation}
  \label{eq:z_jg}
  z_j g_v z_j^{-1} =\Bigl(\prod_{\alpha\in \Phi^+} \zeta_\alpha\Bigl(\pi^{\frac{j\alpha(\lambda)}{r}}b_{\alpha,v}\Bigr)\Bigr) \cdot x_v \cdot \Bigl(\prod_{\alpha\in \Phi^-} \zeta_\alpha\Bigl(\pi^{\frac{j\alpha(\lambda)}{r}}b_{\alpha,v}\Bigr)\Bigr)
\end{equation}
and this element belongs to $G((z_v))(R')$ by the definition of $j$. Thus, for every $\alpha\in \Phi^{\pm}$, we have
\begin{equation}
  \label{eq:inegvaluation2}
  \frac{j\alpha(\lambda)}{r}+\val_\pi(b_{\alpha,v})\geq 0.
\end{equation}
The inequality (\ref{eq:inegvaluation2}) is strict for $\alpha\in \Phi^+$.

\begin{lemme} \label{lem:existstricte2}
There exists an admissible trivialization $(\beta_R,(\al_{R,v}))$ of $m$ such that:
  \begin{enumerate}
  \item $j=j((\beta_R,(\al_{R,v})))$;
  \item the inequality (\ref{eq:firstineg}) holds;
  \item the inequality (\ref{eq:inegvaluation2}) is strict for every $\alpha\in \Phi^{-}$.
  \end{enumerate}
\end{lemme}

\begin{preuve}
Let $(\beta_R,(\al_{R,v}))$ be an admissible trivialization of $m$ satisfying the conclusions of Lemma \ref{lem:existstricte} (i.e., the first two assertions). Let $(X,(g_v)_{v\in V},t)$ be the corresponding triple from Proposition \ref{prop:BL}. We now use the notations introduced above, in particular in (\ref{eq:decompogv}) and (\ref{eq:inMN_Q}). Define
$$h_v =z_j g_v z_j^{-1}\in G((z_v))(R[\pi^{\frac{1}{r}}])$$
and
$$Y=\Ad(z_j)X \in \ggo(R[C^V][\pi^{\frac{1}{r}}]).$$

\begin{lemme} \label{lem:m'}
The following assertions hold:
  \begin{enumerate}
  \item The triple $(Y, (h_v)_{v\in V},t)$ is a triple satisfying the conditions of Proposition \ref{prop:BL} for the ring $R'$.
  \item Let $m'\in \mc(R')$ be the point associated to this triple by Proposition \ref{prop:BL}. Then, after base change to $K'$, the points $m'$ and $m$ become isomorphic.
  \end{enumerate}
\end{lemme}

\begin{preuve}
For assertion 1, one needs to check that the triple $(Y, (h_v)_{v\in V},t)$ satisfies:
  \begin{itemize}
   \item $Y \in \ggo(R'[C^V])$;
   \item For every $v\in V$, $h_v \in G((z_v))(R')$ and
$$\Ad(h_v^{-1})Y\in z_v^{-d_v}\ggo[[z_v]](R');$$
\item $t$ is an element of $\tgo^{\reg}(R')$ whose characteristic $\chi(t)$ equals the reduction modulo $z_\infty$ of $\chi_G(Y)$.
   \end{itemize}
The first condition and the fact that $h_v \in G((z_v))(R')$ follow from the definition of $j$. Hence,
$$\Ad(h_v^{-1})Y \in  \ggo((z_v))(R').$$
In fact, we can replace $\ggo((z_v))$ by $\ggo[[z_v]]$ because on one hand,
$$\Ad(h_v^{-1})Y=\Ad(z_j) (\Ad(g_v^{-1})X),$$
and on the other hand,
$$\Ad(g_v^{-1})X\in z_v^{-d_v}\ggo[[z_v]](R).$$

Similarly, the third condition holds since $X$ and $Y$, being conjugate, have the same characteristic.

For assertion 2, note that $z_j\in G(K')$ and we conclude using Corollary \ref{cor:BL}.
\end{preuve}

\begin{lemme} \label{lem:P-}
Continuing with the notations of Lemma \ref{lem:m'}. Let $m'_\kappa\in \mc(\kappa)$ be the point obtained from $m'$ by base change $\Spec(\kappa)\to\Spec(R)$. Then:
  \begin{enumerate}
  \item $f(m')\in \Ac_{L,\el}(\kappa)$;
  \item For every $P\in \fc(L)$ and every reduction $m'_{\kappa,P}$ of $m'_\kappa$ to $P$, we have
  $$\deg(m'_{\kappa,P})=H_P((\hb_v)_{v\in V}),$$
  where $\hb_v\in G((z_v))(\kappa)$ is the reduction modulo $\pi^{\frac{1}{r}}$ of $h_v$;
  \item For every $P\in \fc(L)$ such that $P\subset Q$, let $m_{\kappa,P}$ be a reduction of $m_\kappa$ to $P$ and $m'_{\kappa,P^-}$ a reduction of $m'_\kappa$ to $P^-=(M_Q\cap P)N_{Q^-}$; then we have the equality
  $$\deg(m'_{\kappa,P^-})=\deg(m_{\kappa,P}).$$
  \end{enumerate}
\end{lemme}

\begin{preuve}
The triple $(\Yb, (\hb_v)_{v\in V},\bar{t})$ is the reduction modulo $\pi^{\frac{1}{r}}$ of the triple $(Y, (h_v)_{v\in V},t)$. Therefore, the point $m'_\kappa$ is obtained from the triple $(\Yb, (\hb_v)_{v\in V},\bar{t})$ via the bijection of Proposition \ref{prop:BL}. The inequality (\ref{eq:firstineg}) implies that $\Yb$ equals $\bar{X}$ as defined in (\ref{eq:tripletadele}). It follows that $f(m'_\kappa)=f(m_\kappa)$, and this point belongs to $\Ac_{L,\el}(\kappa)$ (cf. \S \ref{S:YHN}), which gives assertion 1. Recall that $\bar{X}\in \lgo(\kappa[C^V])$ is conjugate to $\bar{t}_a$ via $\bar{g}_\infty\in L((z_\infty))(\kappa)$ (cf. assertions 2 and 3 of Lemma \ref{lem:triv_en_v}). Then assertion 2 follows from Corollary \ref{cor:reducadel}.

To prove assertion 3, let $P\in \fc(L)$ such that $P\subset Q$. Let $P^-\in \fc(L)$ be defined by $P^-=(M_Q\cap P)N_{Q^-}$. Then $P^-\subset Q^-$. From equality (\ref{eq:z_jg}) and the fact that inequality (\ref{eq:inegvaluation2}) is strict for $\alpha\in \Phi^+$, we deduce
$$\hb_v\in N_{Q^-}((z_v))(\kappa)\,\xb_v,$$
and
$$\deg(m'_{\kappa,P^-})=H_{P^-}((\hb_v)_{v\in V})=H_{P^-\cap M_Q}((\xb_v)_{v\in V}).$$

But equality (\ref{eq:inMN_Q}) also implies that
$$\deg(m_{\kappa,P})=H_{P}((\gb_v)_{v\in V})=H_{P\cap M_Q}((\xb_v)_{v\in V}).$$
Since $P\cap M_Q=P^{-}\cap M_Q$, we have
$$\deg(m'_{\kappa,P^-})=\deg(m_{\kappa,P}).$$
\end{preuve}

According to the notations in Paragraph \ref{S:YHN} and the equivalences in Proposition \ref{prop:HN}, $\varrho$ is the $\xi$-Harder-Narasimhan point of $m_\kappa$ and $Q_0$ is the parabolic subgroup in $\fc(L)$ satisfying
\begin{itemize}
\item[(A)] $\xi-\varrho \in \ago_{Q_0}^+\cap \ago_T^G$;
\item[(B)] the projection of $\varrho$ onto $\ago_L^G$ belongs to the projection of the convex hull of the points $-\deg(m_{\kappa,P})$ for $P\in \pc(L)$ with $P\subset Q_0$ and for every reduction $m_{\kappa,P}$ of $m_\kappa$ to $P$.
\end{itemize}

Assertion (B) above and assertion 3 of Lemma \ref{lem:P-} show that the projection of $\varrho$ onto $\ago_L^G$ belongs to the projection onto $\ago_L^G$ of the convex hull of the points $-\deg(m'_{\kappa,P^-})$ for $P\in \fc^{Q_0}(L)$ (with the notations of Lemma \ref{lem:P-}). Proposition \ref{prop:envpecvx} then implies that the point $\varrho$ belongs to the convex set $\overline{\mathcal{C}}_{m'_\kappa}$ associated to $m'_\kappa$.

It follows that we have the following inequality between distances
$$d(\overline{\mathcal{C}}_{m'_\kappa},\xi)\leq d(\varrho_{m_\kappa},\xi)=d(\overline{\mathcal{C}}_{m_\kappa},\xi)=d.$$
But by the definition of $d$ (cf. \S\ref{S:d}) we necessarily have equality in the above inequality. Hence $\varrho$ is also the $\xi$-Harder-Narasimhan point of $m'_\kappa$. Since $\xi-\varrho\in \ago_{Q_0}^+\cap \ago_T^G$, the parabolic subgroup satisfying assertion (2) of Proposition \ref{prop:HN} for $m'_\kappa$ is necessarily $Q_0$. The projection of $\varrho$ onto $\ago_L^G$ is therefore equal to the projection onto $\ago_L^G$ of the convex hull of the points $-\deg(m'_{\kappa,P})$ associated to the reductions $m'_{\kappa,P}$ of $m'_\kappa$ to the subgroups $P\in \pc(M)$ with $P\subset Q_0$. In particular, the projection of $\varrho$ onto the line $\ago_{Q}^G$ equals the point $-\deg(m'_{\kappa,Q})$, where $m'_{\kappa,Q}$ is a reduction of $m'_\kappa$ to $Q$ (since $\ago_G=\{0\}$ as $G$ is semisimple). Assertion (B) also implies that this projection equals $-\deg(m_{\kappa,Q})$. Hence,
$$\deg(m'_{\kappa,Q})=\deg(m_{\kappa,Q}),$$
which, combined with assertion 3 of Lemma \ref{lem:P-}, gives
$$\deg(m'_{\kappa,Q})=\deg(m'_{\kappa,Q^-}).$$
Equivalently, by assertion 2 of Lemma \ref{lem:P-},
\begin{equation}
  \label{eq:HQ}
  H_{Q^-}((\hb_v)_{v\in V})=H_{Q}((\hb_v)_{v\in V}).
\end{equation}
Lemma \ref{lem:reducM} below then implies that for every $v\in V$, we have
\begin{equation}
  \label{eq:reduchbv}
  \hb_v \in M_Q((z_v))(\kappa) \, G[[z_v]](\kappa).
\end{equation}

Now consider first the point $v=\infty_\kappa$. By condition 3 defining an admissible trivialization, we have $\Ad(g_\infty^{-1})X=t_a$, hence
$$\Ad(h_\infty^{-1})Y=t_a.$$
Inequality (\ref{eq:firstineg}) implies that $\Yb=\bar{X}$. Reducing modulo $\pi^{\frac{1}{r}}$, we obtain
$$\Ad(\hb_\infty^{-1})\bar{X}=\bar{t}_a.$$
But $\bar{X}$ and $\bar{t}_a$ are conjugate by $\bar{g}_\infty\in L((z_\infty))(\kappa)$ (cf. assertion 3 of Lemma \ref{lem:triv_en_v}). Hence, we also have $\hb_\infty\in L((z_\infty))(\kappa)$. Consequently, inequality (\ref{eq:inegvaluation2}) is strict for $v=\infty_\kappa$ and for every $\alpha\in \Phi^\pm$.

Now suppose $v\neq \infty_\kappa$. Also, suppose that there exists $\alpha\in \Phi^-$ such that inequality (\ref{eq:inegvaluation2}) is an equality, i.e.,
\begin{equation}
  \label{eq:egvaluation2}
  \frac{j\alpha(\lambda)}{r}+\val_\pi(b_{\alpha,v})= 0.
\end{equation}
Let $i$ be the largest integer such that there exists $\alpha\in \Phi^-_i$ satisfying (\ref{eq:egvaluation2}). Then
\begin{equation}
  \label{eq:inmN_i}
  \hb_v\in \xb_v N_i((z_v))(\kappa).
\end{equation}
Comparing with (\ref{eq:reduchbv}), we deduce that
\begin{equation}
  \label{eq:inmN_i2}
  \hb_v\in \xb_v N_i[[z_v]](\kappa).
\end{equation}
For every $\alpha\in \Phi_i^-$, choose $u_\alpha\in \kappa[[z_v]]$ and define
$$n=\prod_{\alpha\in \Phi_i^-}\zeta_\alpha(u_\alpha),$$
where the one-parameter groups $\zeta_\alpha$ are those defined in (\ref{eq:zeta}). such that
$$\hb_v n \in \xb_v N_{i+1}[[z_v]](\kappa).$$
Let $n'=z_j^{-1}nz_j$. Then,
$$n'=\zeta_\alpha\Bigl( \pi^{\frac{ij}{r}} u_\alpha\Bigr).$$
Equality (\ref{eq:egvaluation2}) implies that $\frac{ij}{r}$ is a strictly positive integer. Consequently, $n'$ belongs to $N_Q((z_v))(R)$ and its reduction modulo $\pi$ is trivial. Replacing $g_v$ by $g_v n'$, we see that we reduce to the case where
$$\hb_v\in \xb_v N_{i+1}[[z_v]](\kappa),$$
and by induction to the case where $\hb_v=\xb_v$. But in that case, inequality (\ref{eq:inegvaluation2}) is strict for every $\alpha\in \Phi^{\pm}$. This completes the proof of Lemma \ref{lem:existstricte2}.
\end{preuve}

Now let $(\beta_R,(\al_{R,v}))$ be an admissible trivialization of $m$ satisfying the conditions of Lemma \ref{lem:existstricte2}. Let $(X,(g_v)_{v\in V},t)$ be the triple deduced from $m$ via Proposition \ref{prop:BL}. Define
$$j_1=\min\Bigl(\min_{\alpha\in \Phi^-}\Bigl(-\frac{r}{\alpha(\lambda)}\val_\pi(X_\alpha)\Bigr),\min_{v\in V,\alpha\in \Phi^-}\Bigl(-\frac{r}{\alpha(\lambda)}\val_\pi(b_{\alpha,v})\Bigr)\Bigr).$$
This is an integer satisfying $j_1>j$, with
$$\Ad(z_{j_1})X\in \ggo(R[C^V])$$
and for every $v\in V$,
$$z_{j_1}g_v z_{j_1}^{-1}\in G((z_v))(R).$$
This is clearly the desired contradiction, and it completes the proof of Proposition \ref{prop:cruciale}.
\end{preuve}

The following lemma was used in the previous proof (in the proof of Proposition \ref{prop:cruciale}).

\begin{lemme}\label{lem:reducM}
Let $Q$ and $Q^-$ be two opposite parabolic subgroups in $\fc^G$. For every $v\in V$, let $g_v\in G((z_v))(\kappa)$. The equality
\begin{equation}
     H_Q((g_v)_{v\in V})=H_{Q^-}((g_v)_{v\in V})
\end{equation}
implies that for every $v\in V$, there exist $m_v\in M_Q((z_v))(\kappa)$ and $k_v\in G[[z_v]](\kappa)$ such that
$$g_v=m_v\,k_v.$$
\end{lemme}

The proof of Lemma \ref{lem:reducM} will use the following lemma.

\begin{lemme}\label{lem:borel-adj}
Let $B$ and $B'$ in $\fc^G$ be two adjacent Borel subgroups. Let $\alpha$ be the unique root of $T$ in $G$ which is positive for $B$ and negative for $B'$. Let $\zeta_\alpha$ be an isomorphism from $\mathbb{G}_a$ onto the root subgroup $N_\alpha$ corresponding to the root $\alpha$. For every $g\in G((z))(\kappa)$, there exist $t\in T((z))(\kappa)$, $u\in (N_B\cap N_{B'})(\kappa((z)))$, $x\in \kappa((z))$ and $h\in G[[z]](\kappa)$ such that $g=tu\zeta_\alpha(x)h$. For any such decomposition, we have
$$H_{B'}(g)-H_B(g)=\sup(-\val(x),0)\,\alpha^\vee.$$
\end{lemme}

\begin{preuve}
This lemma is well known. We give the proof for the reader's convenience. The existence of the decomposition of $g$ follows from the Iwasawa decomposition and from the equality $N_B=(N_B\cap N_{B'})N_\alpha$. With the notations of the statement, write $g=tu\zeta_\alpha(x)h$. We have
$$H_{B'}(g)-H_B(g)=H_{B'}(\zeta_\alpha(x)).$$
If $\val(x)\geq 0$, then $H_{B'}(\zeta_\alpha(x))=0$. Suppose now that $\val(x)<0$. The following equality in $SL_2((z))(\kappa)$
$$\begin{pmatrix} 1 & x\\ 0 & 1\end{pmatrix}=\begin{pmatrix} x & 0\\ 0 & x^{-1}\end{pmatrix}\begin{pmatrix} 1 & 0\\ x & 1\end{pmatrix}\begin{pmatrix} x^{-1} & 1\\ -1 & 0\end{pmatrix}$$
shows that
$$\zeta_\alpha(x)\in \alpha^\vee(x)\, N_{B'}((z))(\kappa)\, G[[z]](\kappa)$$
and hence
$$H_{B'}(\zeta_\alpha(x))=H_{B'}(\alpha^\vee(x))=-\val(x)\alpha^\vee,$$
which is what had to be proved.
\end{preuve}

We can now prove Lemma \ref{lem:reducM}.

\begin{preuve}
Let us first treat the case where $V=\{v\}$ is a singleton. To lighten the notations, we omit the index $v$. Let $g\in G((z))(\kappa)$. We shall show that $-H_Q(g)+H_{Q^-}(g)$ is a sum with \emph{positive} coefficients of coroots $\alpha^\vee$ in $\Delta_Q^\vee$, and that this sum is zero if and only if there exist $m\in M_Q((z))(\kappa)$ and $k\in G[[z]](\kappa)$ such that $g=mk$. By the Iwasawa decomposition, we may and do assume that $g\in N_Q((z))(\kappa)$.

Let $B\subset Q$ and $B^-\subset Q^-$ be two opposite Borel subgroups of $\fc^G$. Let $n=|\Phi^G|/2$ and let $B_0=B, B_1,\ldots,B_n=B^-$ be a minimal chain of Borel subgroups of $\fc^G$ such that $B_i$ and $B_{i+1}$ are adjacent for $0\leq i<n$. By Lemma \ref{lem:borel-adj}, the difference
$$H_{B^-}(g)-H_B(g)=\sum_{i=0}^{n-1}\bigl(H_{B_{i+1}}(g)-H_{B_i}(g)\bigr)$$
is a linear combination with positive coefficients of coroots of $T$ in $B$. Since $H_{Q^-}(g)-H_Q(g)$ is the image of $H_{B^-}(g)-H_B(g)$ under the canonical morphism $\ago_B\to \ago_Q$, the difference $H_{Q^-}(g)-H_Q(g)$ is a linear combination with positive coefficients of elements of $\Delta_Q^\vee$. Suppose from now on that
$$H_{Q^-}(g)-H_Q(g)=0.$$
Number the roots $(\alpha_i)_{1\leq i\leq n}$ of $T$ in $B$ so that for every $0\leq i<n$
$$\Phi^B\cap \Phi^{B_i}=\{\alpha_{i+1},\ldots,\alpha_n\}.$$
Let $I\subset \{1,\ldots,n\}$ be such that $\Phi^{N_Q}=\{\alpha_i \mid i\in I\}$. For every $i\in I$ we therefore have
\begin{equation}
  \label{eq:HBi}
  H_{B_i}(g)-H_{B_{i-1}}(g)=0.
\end{equation}
For every $1\leq i\leq n$, let $\zeta_i$ be an isomorphism from $\mathbb{G}_a$ onto the root subgroup corresponding to the root $\alpha_i$. We have assumed $g\in N_Q((z))(\kappa)$. There therefore exist $x_i\in \kappa((z))$ for $1\leq i\leq n$ such that $x_i=0$ if $i\notin I$ and
$$g=\zeta_n(x_n)\zeta_{n-1}(x_{n-1})\cdots \zeta_1(x_1).$$
Note that for $1\leq i<n$ we have
$$\zeta_n(x_n)\zeta_{n-1}(x_{n-1})\cdots \zeta_{i+1}(x_{i+1})\in (N_{B_i}\cap N_{B_{i-1}})(\kappa((z))).$$
Using Lemma \ref{lem:borel-adj} and equality (\ref{eq:HBi}), one sees by induction on $i$ that $\val(x_i)\geq 0$ for every $1\leq i\leq n$. In particular $g\in N_Q[[z]](\kappa)$, as desired. This concludes the proof when $V$ is a singleton.

Let us return to the general case of a finite nonempty set $V$. We have just seen that for every $v\in V$ the difference $-H_Q(g_v)+H_{Q^-}(g_v)$ is a combination with positive coefficients of elements of $\Delta_Q$. It follows that the equality
$$ H_Q((g_v)_{v\in V})=H_{Q^-}((g_v)_{v\in V})$$
holds if and only if for every $v\in V$
$$ H_Q(g_v)=H_{Q^-}(g_v).$$
We are thus reduced to the singleton case.
\end{preuve}
\end{paragr}

\begin{paragr}[Where the Contradiction is Obtained.] --- \label{S:lacontradiction}
For every integer $n$, let
$$R_n=R/ \pi^n R.$$
Let $n\geq 1$ be an integer. Let $(\beta_R, \al_{R,v})$ be an admissible trivialization of $m$ satisfying
\begin{equation}
  \label{eq:j>n}
  j(\beta_R, \al_{R,v})\geq n,
\end{equation}
(cf. Proposition \ref{prop:cruciale}). Let $(X,(g_v)_{v\in V},t)$ be the triple associated to $m$ and this trivialization via Proposition \ref{prop:BL}. Let $(X_n,(g_{n,v})_{v\in V},t_n)$ be the triple obtained by reduction modulo $\pi^n$. This triple satisfies the following:
\begin{enumerate}
\item $X_n\in \qgo(R_n[C^V])$;
\item For every $v\in V$, $g_{n,v}\in Q((z_v))(R_n)$;
\item $\Ad(g_{n,\infty})^{-1}X_n= t_{a,n}$, where $t_{a,n}$ is the reduction modulo $\pi^n$ of $t_a$.
\end{enumerate}

From this triple, by Proposition \ref{prop:BL} applied to the group $Q$, one deduces a triple $m_{Q,n}=(\ec_{Q,n},\theta_{Q,n},t_n)$ consisting of a $Q$-bundle $\ec_{Q,n}$ over
$$C_n=C\times_k R_n,$$
with a section $\theta_{Q,n}$ of $\Ad_D(\ec_{Q,n})$ such that $\chi_Q(\theta_{Q,n}(\infty_{R_n}))=\chi_Q(t_{n})$. Moreover, this triple is a reduction to $Q$ of the point $m_n\in \mc(R_n)$ obtained by base change to $R_n$.

The algorithm used in the proof of Proposition \ref{prop:cruciale} consists of conjugating $X$ (and translating $g_v$ on the right accordingly), respectively translating $g_v$ on the left, by elements of the form $\zeta_\alpha(u_\alpha)$ (where $u_\alpha\in R[C^V]$, respectively in $R((z_v))$, with $\pi$-adic valuation at least $n$), so as to obtain a triple $(X',(g'_v)_{v\in V},t)$ associated to an admissible trivialization whose associated $j$ is greater than $n+1$. In particular, the reduction modulo $\pi^n$ of $(X',(g'_v)_{v\in V},t)$ is equal to $(X_n,(g_{n,v})_{v\in V},t_n)$. Let $(X_{n+1},(g_{n+1,v})_{v\in V},t_{n+1})$ be the reduction modulo $\pi^{n+1}$ of $(X',(g'_v)_{v\in V},t)$. Let $m_{Q,n+1}=(\ec_{Q,n+1},\theta_{Q,n+1},t_{n+1})$ be the reduction to $Q$ of $m_{n+1}\in \mc(R_{n+1})$ deduced as above. The base change of $m_{Q,n+1}$ to $C_n$ recovers $m_{Q,n}$.

By induction, we obtain for every integer $n$ a reduction $m_{Q,n}=(\ec_{Q,n},\theta_{Q,n},t_n)$ to $Q$ of the point $m_n=(\ec_n,\theta_n,t_n) \in \mc(R_n)$ deduced from $m\in \mc(R)$ by base change, such that the image of $m_{n+1}$ in $\mc(R_n)$ is isomorphic to $m_n$. The $Q$-equivariant morphism
$$\ec_{Q,n} \to \ec_{Q,n}\times^G_k Q \simeq \ec_n$$
induces, after taking the quotient, a section
$$\sigma_n \ : \ C_n \to \ec_n/Q.$$
When $n$ varies, the sections obtained form a commutative diagram
$$\xymatrix{\ec_{n}/Q \ar[r]  & \ec_{n+1}/Q \\ C_{n} \ar[r]  \ar[u]^{\sigma_{n}} & C_{n+1} \ar[u]_{\sigma_{n+1}} } $$
where the horizontal morphisms are the canonical transition maps. They form an inductive system $\hat{\sigma}=(\sigma_n)_{n\geq 1}$. Let
$$\hat{C}=\varinjlim_{n} C_n$$
and
$$\widehat{\ec/Q}=\varinjlim_{n}\ec_n/Q.$$
The inductive system $\hat{\sigma}$ gives an element, still denoted $\hat{\sigma}$, in
$$\Hom_{\Spf(R)}(\hat{C},\widehat{\ec/Q}).$$
Since $C_R$ is proper over $\Spec(R)$ and $\ec/Q$ is separated and of finite type over $\Spec(R)$, by Grothendieck's algebraization theorem (Theorem 5.4.1 of \cite{Gro61}), the natural map
$$\Hom_{\Spec(R)}(C_R,\ec/Q)\to \Hom_{\Spf(R)}(\hat{C},\widehat{\ec/Q})$$
is a bijection. Therefore, there exists $\sigma\in \Hom_{\Spec(R)}(C_R,\ec/Q)$ which maps to $\hat{\sigma}$ under the above bijection. By again using the algebraization theorem, we see that $\sigma$ is in fact a section of $\ec/Q$. Let $\ec_Q$ be the $Q$-bundle over $C_R$ defined by
$$\ec_Q=C_R\times_{\sigma,\ec/Q}\ec.$$
By construction, $\ec_Q$ is a reduction of $\ec$ to $Q$.

Let
$$\widehat{\Ad_D(\ec_Q)}=\varinjlim_{n}\Ad_D(\ec_{Q,n}).$$
For every integer $n\geq 1$, the section $\theta_n$ of $\Ad_D(\ec_n)$ factors through the section $\theta_{Q,n}$ of $\Ad_D(\ec_{Q,n})$. By construction, the sections $\theta_n$ form an inductive system and hence an element, denoted $\hat{\theta}$, in
$$\Hom_{\Spf(R)}(\hat{C},\widehat{\Ad_D(\ec_Q)}).$$
By the aforementioned algebraization theorem, it follows that $\theta$ factors through a section $\theta_Q$ of $\Ad_D(\ec_Q)$.

By construction, the reduction modulo $\pi^n$ of $\chi_Q(\theta_{Q})(\infty_R)$ equals $\chi_Q(\theta_{Q,n})(\infty_{R_n})=\chi_Q(t_{n})$. It follows that
 $$\chi_Q(\theta_{Q})(\infty_R)=\chi_Q(t).$$
Therefore, the triple $m_Q=(\ec_Q,\theta_Q,t)$ is a reduction of $m$ to $Q$. Let $m_{Q,K}$ and $m_{Q,\kappa}$ be the reductions to $Q$ of $m_K$ and $m_\kappa$, respectively, obtained by base change. By the flatness of $\ec_Q$ over $C_R$, we have
$$\deg(m_{Q,K})=\deg(m_{Q,\kappa}).$$
But $m_K$ is $\xi$-semistable; hence,
$$\xi\in -\deg(m_{Q,K})-\, \overline{^+ \ago_Q}+ \ago_T^Q.$$
In Paragraph \ref{S:YHN}, we introduced the parabolic subgroup $Q_0\subset Q$ which satisfies $\xi-\varrho\in \ago_{Q_0}^+$ and $\varrho\in -\deg(m_{Q_0,\kappa})+\ago_T^{Q_0}$. Therefore,
$$\xi \in -\deg(m_{Q_0,\kappa})+\ago_{Q_0}^+ + \ago_T^{Q_0}.$$
It follows that
$$\xi\in \Bigl(-\deg(m_{Q,K})-\, \overline{^+ \ago_Q}+ \ago_T^Q \Bigr) \cap \Bigl(-\deg(m_{Q_0,\kappa})+\ago_{Q_0}^+ + \ago_T^{Q_0}\Bigr).$$
Since $\deg(m_{Q,K})=\deg(m_{Q,\kappa})$, we deduce that
$$\Bigl(-\, \overline{^+ \ago_Q}+ \ago_T^Q \Bigr) \cap \Bigl(\ago_{Q_0}^+ + \ago_T^{Q_0}\Bigr)\neq\emptyset.$$
Projecting onto $\ago_Q$, we find
$$-\overline{^+ \ago_Q} \cap \ago_Q^+\neq\emptyset,$$
which is impossible (see the end of \S\ref{S:d}). This is the desired contradiction.
\end{paragr}

% ---------- sec09 ----------
\section{Separation of the Morphism $f^\xi$}\label{sec:separation}

\begin{paragr}[Statement.] --- \label{S:sepenonce}
The aim of this section is to prove the following theorem.

\begin{theoreme}\label{thm:separation}
Suppose $G$ is semisimple. For every $\xi\in \ago_T$, the morphism
$$f^\xi \ : \ \mc^\xi \to \Ac$$
is separated.
\end{theoreme}

\begin{remarque}
The statement and the proof of Theorem \ref{thm:separation} are inspired by a theorem of Langton (cf. \cite[\S 3, Theorem]{Lan75}). Faltings gave a proof of Theorem \ref{thm:separation} for $\xi=0$ and when the base field is of characteristic zero (cf. \cite[Theorem II.4]{Fal93}).
\end{remarque}

In all that follows, we suppose $G$ \emph{semisimple}. By the valuative criterion of separation, it is equivalent to prove the following statement.

\begin{proposition}\label{prop:separation}
Let $\kappa$ be an algebraically closed extension of $k$. Let $R$ be a complete discrete valuation ring, with fraction field $K$ and residue field $\kappa$. Let $m$ and $m'$ be two elements of $\mc^\xi(R)$, and let $m_K$ and $m'_K$ be the elements of $\mc^\xi(K)$ deduced from them by base change. Suppose that these objects satisfy the following two conditions:
\begin{itemize}
\item $f^\xi(m)=f^\xi(m')$;
\item there exists an isomorphism $\phi_K \ : \ m'_K\to m_K$.
\end{itemize}
Then there exists a unique isomorphism $\phi \ : \ m'\to m$ which extends $\phi_K$.
\end{proposition}
\end{paragr}

\begin{paragr}[Where We Reduce to a Problem on a Trait.] --- \label{S:sur-un-trait}
In all that follows, let $m=(\ec,\theta,t)$ and $m'=(\ec',\theta',t)$ be two triples in $\mc^\xi(R)$ such that $f^\xi(m)=f^\xi(m')$. Let $m_K=(\ec_K,\theta_K,t)$ and $m'_K=(\ec'_K,\theta'_K,t)$ be the triples of $\mc^\xi(K)$ deduced from them by base change. Let $\phi_K \ : \ m'_K\to m_K$ be an isomorphism. We again denote by $\phi_K$ the underlying $G$-equivariant isomorphism
$$\phi_K \ : \ \ec'_K \to \ec_K.$$
The isomorphism deduced from it
$$\Ad_D(\phi_K) \ : \ \Ad_D(\ec'_K) \to \Ad_D(\ec_K)$$
sends $\theta'_K$ to $\theta_K$. We seek to extend $\phi_K$ to a $G$-equivariant isomorphism
$$\phi \ : \ \ec' \to \ec$$
in such a way that the isomorphism deduced from it
$$\Ad_D(\phi) \ : \ \Ad_D(\ec') \to \Ad_D(\ec)$$
sends $\theta'$ to $\theta$. Let $\pi$ be a uniformizer of $R$ and let $B$ be the local ring of $C_R$ at the point of codimension $1$ defined by the ideal $(\pi)$. It is therefore a discrete valuation ring whose residue field is the function field $\kappa(C)$ of the curve $C_\kappa$ and whose fraction field is the function field $F$ of the curve $C_K$.

\begin{lemme}\label{lem:prolongB}
In order for $\phi_K$ to extend to $C_R=C\times_k R$ it is necessary and sufficient that it extend to $\Spec(B)$. Moreover, such extensions, if they exist, are unique.
\end{lemme}

\begin{preuve}
A $G$-equivariant isomorphism from $\ec'$ onto $\ec$ is nothing but a global section of the $C_R$-scheme
$$\mathrm{Isom}_G(\ec',\ec)=(\ec'\times_{C_R}\ec)/G$$
where $G$ acts diagonally on the right on $\ec'\times_{C_R}\ec$. The isomorphism $\phi_K$ defines a section $\sigma_K$ of $\mathrm{Isom}_G(\ec',\ec)$ over $C_K$. Any extension of this section to $C_R$ is necessarily unique, whence the uniqueness in the statement. Suppose that $\phi_K$ extends to $\Spec(B)$. The section $\sigma_K$ then extends to $\Spec(B)$, and consequently to an open subset of $C_R$ which contains all the points of codimension $\leq 1$. Since $\mathrm{Isom}_G(\ec',\ec)$ is affine over $C_R$ --- it is even a torsor under $\Aut_G(\ec)$ --- such a section extends automatically to $C_R$ (cf. \cite[Corollary 20.4.12]{Gro67}).
\end{preuve}

Note that if $\phi_K$ extends to an isomorphism $\phi$, the isomorphism $\Ad_D(\phi)$ necessarily sends $\theta'$ to $\theta$, since the sections $\Ad_D(\phi)(\theta')$ and $\theta$, which coincide on an open subset of $C_R$, are in fact equal.
\end{paragr}

\begin{paragr}[Study on the Trait $\Spec(B)$.] --- \label{S:etude-trait}
For every finite extension $K'$ of $K$, let $R'$ be the integral closure of $R$ in $K'$. Let $F'$ be the function field of the relative curve $C_{R'}=C\times_k R'$ and let $B'$ be the integral closure of $B$ in $F'$. In the proof of Lemma \ref{lem:prolongB}, we saw that $\phi_K$ is interpreted as a section of the affine $C_R$-scheme $\mathrm{Isom}_G(\ec',\ec)$, and that extending $\phi_K$ amounts to extending this section to $\Spec(B)$. Since $B$ is normal, $\phi_K$ extends to $\Spec(B)$ if and only if the base change of $\phi_K$ to $K'$ extends to $B'$. In what follows, we may always, if need be, replace $K$ by $K'$.

Let $(a,t)\in \Ac_G(R)$ be defined by $(a,t)=f^\xi(m)$. Let $t_a\in \tgo^{G\text{-}\reg}[[z_\infty]](R)$ be the unique element whose characteristic is the restriction of $a$ to $\Spec(R[[z_\infty]])$ and whose reduction modulo $z_\infty$ is $t$.

Let $M\subset L$ be the two Levi subgroups in $\lc$ such that $(a,t)$ belongs, in the generic fiber, to $\chi^L_G(\Ac_{L,\el}(K))$ and, in the special fiber, to $\chi^M_G(\Ac_{M,\el}(\kappa))$ (cf. Proposition \ref{prop:reunionAM}). Let $(a_L,t)\in \Ac_L(R)$ be the unique point such that $(a,t)=\chi^L_G((a_L,t))$ (cf. Proposition \ref{prop:immersionfermee}).

\begin{lemme}\label{lem:situation}
After possibly replacing $K$ by a large enough finite extension $K'$ and $R$ by $R'$, we are in the following situation:
\begin{enumerate}
\item there exists a point $X\in \lgo(B)$ which satisfies:
  \begin{enumerate}
  \item $\chi_L(X)$ coincides with $a_L$ on $\Spec(B)$;
  \item the reduction $\bar{X}$ of $X$ modulo $\pi$ belongs to $\mgo(\kappa(C))$ and is conjugate to the reduction modulo $\pi$ of $t_a$ by an element of $M((z_\infty))(\kappa)$;
  \end{enumerate}
\item there exist trivializations of the $G$-torsors $\ec$ and $\ec'$ over $\Spec(B)$ such that the trivializations of $\Ad_D(\ec)$ and $\Ad_D(\ec')$ deduced from them send $\theta$ and $\theta'$ to $X$.
\end{enumerate}
\end{lemme}

\begin{preuve}
The Kostant section shows that a point $X\in \lgo(B)$ satisfying (1)(a) exists. By Proposition \ref{prop:caractdesAM} and Lemma \ref{lem:cjsep}, we may suppose that, after possibly conjugating $X$ by an element of $L(\kappa(C))$, assertion (1)(b) is also satisfied.

If the extension $K'$ is large enough, the torsors $\ec$ and $\ec'$ are trivial over $\Spec(B')$ (Theorem 2 of \cite{DS95}). After possibly replacing $R$ by $R'$, we may suppose that the torsors $\ec$ and $\ec'$ are trivial over $\Spec(B)$ and even, by Lemma \ref{lem:conjB} below, that there exist trivializations which satisfy assertion (2).
\end{preuve}

\begin{lemme}\label{lem:conjB}
Let $a\in \car^{\reg}(B)$. Let $X$ and $Y$ be two elements of $\ggo(B)$ such that
$$\chi_G(X)=\chi_G(Y).$$
There exists a finite extension $K'$ of $K$ such that $X$ and $Y$ are conjugate by an element of $G(B')$, where $B'$ is the integral closure of $B$ in the function field of $C_{K'}$.
\end{lemme}

\begin{preuve}
For every $B$-algebra $A$, let
$$\Tc(A)=\{g\in G(A)\mid \Ad(g)X=Y\}.$$
Recall that the fraction field of $B$ is the function field $F$ of $C_K$ and that its residue field is the function field $\kappa(C)$ of $C_\kappa$. By Lemma \ref{lem:cjsep}, the set $\Tc(\kappa(C))$ is nonempty and, after possibly replacing $K$ by a finite extension, we may suppose that the same holds for $\Tc(F)$. We deduce that the functor $\Tc$ is represented by a torsor, still denoted $\Tc$, under the group scheme $T_X$ which centralizes $X$. Now, in view of the hypothesis $a\in \car^{\reg}(B)$, this group scheme is a torus scheme. Since the torsor $\Tc$ is trivial in the generic fiber $\Spec(F)$, it is trivial globally (cf. for example \cite[Proposition 2.2]{CS78}). It therefore possesses a section over $\Spec(B)$.
\end{preuve}

From now on we fix $X\in \ggo(B)$ and trivializations of $\ec$ and $\ec'$ over $\Spec(B)$ which satisfy the assertions of Lemma \ref{lem:situation}. We then obtain a trivialization of $\phi_K$ which is identified with left translation by an element $\delta\in G(F)$ which satisfies
$$\Ad(\delta^{-1})X=X.$$

It follows from Lemma \ref{lem:prolongB} and the above considerations that we have the following lemma.

\begin{lemme}\label{lem:deltaB}
The morphism $\phi_K$ extends if and only if $\delta\in G(B)$.
\end{lemme}

Since the restriction of $a$ defines an element of $\car^{\reg}(B)$, the reduction modulo $\pi$ of $X$ is still regular semisimple. The centralizer of $X$ in $G\times_k B$ is therefore a torus scheme over $B$. Since $B$ is normal, it is known that such a torus scheme is isotrivial. In other words, the torus $T_{X,F}$ splits over a separable and unramified extension. We deduce that there exist $\lambda\in X_*(T_X)$ and an element $h\in T_X(F)\cap G(B)$ such that
$$\delta=\lambda(\pi)h.$$
Since $h\in G(B)$ centralizes $X$, we may compose the trivialization of $\ec$ with $h$ without affecting condition (2) of Lemma \ref{lem:situation}. After possibly changing the trivialization, we may and shall suppose that $h=1$. Note that $\lambda$ is necessarily fixed under $\Gal(F_s/F)$. The characteristic $\chi_L(X)$ is equal to the restriction of $a_L$ to $\Spec(B)$. Since $(a_L,t)$ defines, in the generic fiber, a point of $\Ac_{L,\el}(K)$, the torus $T_{X,F}$ is an elliptic $F$-subtorus of $L$ (cf. Corollary \ref{cor:caractdesAM} and Lemma \ref{lem:cjsep}), which translates into the equality
$$X_*(T_X)^{\Gal(F_s/F)}=X_*(A_L)$$
where $A_L$ is the connected center of $L$. Thus $\lambda\in X_*(A_L)$. The following lemma, combined with Lemma \ref{lem:deltaB}, shows the existence of the extension of $\phi_K$.

\begin{lemme}\label{lem:lambda=0}
Under the above hypotheses, we have $\lambda=0$.
\end{lemme}
\end{paragr}

\begin{paragr}[Proof of Lemma \ref{lem:lambda=0}.] --- \label{S:preuvelambda}
We continue with the notations of the previous paragraph. Note first that the lemma is obvious if $L=G$, since in that case $X_*(A_G)=0$. In what follows, we therefore suppose that $L\neq G$. We argue by contradiction, supposing $\lambda\neq 0$, and we seek a contradiction. Let $P$, respectively $\bar{P}$, be the parabolic subgroup in $\fc(L)$ defined by the following condition: the roots $\al$ of $T$ in $P$ satisfy the inequality $\al(\lambda)\geq 0$, respectively $\al(\lambda)\leq 0$. These subgroups are opposite, in the sense that the intersection $P\cap \bar{P}$ is a common Levi subgroup. Note that this intersection contains the Levi subgroup $M$ defined in \S\ref{S:etude-trait}, and that these parabolic subgroups are proper since $\lambda$ is not zero. Let $r$ be the largest integer $\al(\lambda)$ as $\al$ runs over the set $\Phi^G_T$ of roots of $T$ in $G$. Let
$$\ggo=\bigoplus_{\al\in X^*(T)}\ggo_\al$$
be the decomposition of $\ggo$ into eigenspaces for the action of $T$. For every integer $i$ satisfying $-r\leq i\leq r$, let
$$\ggo_i=\bigoplus_{\{\al\in X^*(T)\mid \al(\lambda)=i\}}\ggo_\al$$
and
$$\ggo_+=\bigoplus_{\{-r<i\leq r\}}\ggo_i.$$
The subspaces $\ggo_+$ and $\ggo_r$ are stable under $P$, and $\ggo_{-r}$ is stable under $\bar{P}$.

The inclusion
\begin{equation}
  \label{eq:inclusionpir}
  \pi^r\Ad(\lambda(\pi))\ggo(B)\subset \ggo(B)
\end{equation}
is an inclusion between two sub-$B$-modules of $\ggo(F)$, free and of maximal rank, each of which provides an extension to $C_R$ of the vector bundle $\Ad_D(\ec_K)$ over $C_K$ (cf. \cite[Proposition 6]{Lan75}). The one associated to $\ggo(B)$ is none other than $\Ad_D(\ec)$, and we denote the other by $\vc$. In what follows, we denote by $\ec_\kappa$ and $\vc_\kappa$ the $G$-torsor and the vector bundle over $C_\kappa=C\times_k \kappa$ deduced from $\ec$ and $\vc$ by base change. From the inclusion (\ref{eq:inclusionpir}) we deduce a morphism of vector bundles
\begin{equation}
  \label{eq:VversAdD}
  \vc\to \Ad_D(\ec).
\end{equation}
Let $\wc$ and $\Qc$ be the kernel and the image of the morphism $\vc_\kappa\to \Ad_D(\ec_\kappa)$. Note that $\wc$ is locally a direct factor of $\vc_\kappa$. This is not the case for $\Qc$, which is simply a locally free subsheaf of $\Ad_D(\ec_\kappa)$. We have
$$\deg(\vc_\kappa)=0.$$
Indeed, by flatness we have $\deg(\vc_\kappa)=\deg(\vc_K)$, and on the other hand $\deg(\vc_K)=\deg(\Ad_D(\ec_K))=0$. Since $\deg(\vc_\kappa)=0$, we have
\begin{equation}
  \label{eq:degW=-degQ}
  \deg(\wc)=-\deg(\Qc).
\end{equation}

Let
$$\Ad_D(\ec')\to \vc$$
be the unique isomorphism which, in the generic fiber, is the composite of the isomorphism
$$\Ad_D(\phi_K) \ : \ \Ad_D(\ec'_K)\to \Ad_D(\ec_K)$$
with the homothety of ratio $\pi^r$, and which on $\Spec(B)$ coincides with the isomorphism
$$\ggo(B)\to \pi^r\Ad(\lambda(\pi))\ggo(B)$$
given by $\pi^r\Ad(\lambda(\pi))$. Composing this isomorphism with the morphism $\vc\to \Ad_D(\ec)$ defined in (\ref{eq:VversAdD}), we obtain a morphism of vector bundles over $C_R$
\begin{equation}
  \label{eq:AdE'versAdE}
  \Ad_D(\ec')\to \Ad_D(\ec),
\end{equation}
which gives, in the special fiber, a morphism
\begin{equation}
  \label{eq:AdE'kversAdEk}
  \Ad_D(\ec'_\kappa)\to \Ad_D(\ec_\kappa),
\end{equation}
whose kernel we denote by $\wc'$. By construction, we have $\wc'\simeq \wc$ and the image of this morphism is the vector bundle $\Qc$. In view of (\ref{eq:degW=-degQ}), we have
\begin{equation}
  \label{eq:degW'=-degQ}
  \deg(\wc')=-\deg(\Qc).
\end{equation}

We fixed in \S\ref{S:etude-trait} trivializations of $\ec$ and $\ec'$ over $\Spec(B)$ which satisfy the assertions of Lemma \ref{lem:situation}. Through these choices, the parabolic subgroup $P$ of $G$ defines a generic reduction of $\ec'_\kappa$ to $P$ which extends automatically to a reduction, denoted $\ec'_{P,\kappa}$, of $\ec'_\kappa$ to $P$ over $C_\kappa$. Likewise, we denote by $\ec_{\bar{P},\kappa}$ the reduction of $\ec_\kappa$ to $\bar{P}$ which extends the generic reduction given by $\bar{P}\subset G$.

\begin{lemme}\label{lem:WQ}
With the above notations, we have:
\begin{enumerate}
\item an equality $\ec'_{P,\kappa}\times^{P,\Ad}_k \ggo_+=\wc'$;
\item an injection $\Qc\hookrightarrow \ec_{\bar{P},\kappa}\times^{\bar{P},\Ad}_k \ggo_{-r}$ between vector bundles of the same rank.
\end{enumerate}
\end{lemme}

\begin{preuve}
Since $\ec'_{P,\kappa}\times^{P,\Ad}_k \ggo_+$ and $\wc'$ are vector subbundles of $\Ad_D(\ec'_\kappa)$, it suffices to check 1 at the generic point of $C_\kappa$. Likewise, since $\Qc$ and $\ec_{\bar{P},\kappa}\times^{\bar{P},\Ad}_k \ggo_{-r}$ are locally free subsheaves of $\Ad_D(\ec_\kappa)$, and since $\ec_{\bar{P},\kappa}\times^{\bar{P},\Ad}_k \ggo_{-r}$ is even locally a direct factor, it suffices to check 2 at the generic point of $C_\kappa$. Now, in the trivializations of $\ec$ and $\ec'$ over $\Spec(B)$ that we fixed in \S\ref{S:etude-trait}, the morphism (\ref{eq:AdE'versAdE}) is the morphism $\ggo(B)\to \ggo(B)$ given by $\pi^r\Ad(\lambda(\pi))$. In particular, upon reduction modulo $\pi$, this morphism is the projection of $\ggo(\kappa(C))$ onto $\ggo_{-r}(\kappa(C))$ with kernel $\ggo_+(\kappa(C))$. The lemma is then clear.
\end{preuve}

\begin{lemme}\label{lem:mu}
Let $\mu\in X^*(P)$ be the determinant of the adjoint action of $P$ on $\ggo_+$. We have the inequality
\begin{equation}
  \label{eq:inegmu}
  \deg(\mu(\ec'_{P,\kappa}))\geq \deg(\mu(\ec_{\bar{P},\kappa})).
\end{equation}
\end{lemme}

\begin{remarque}
The character $\mu$ is a highest weight in the representation of $G$ on $\bigwedge^{\dim(\ggo_+)}\ggo$. It is therefore written in the basis $\hat{\Delta}_P$ with positive coefficients, at least one of which is nonzero.
\end{remarque}

\begin{preuve}
By assertion (1) of Lemma \ref{lem:WQ}, we have
$$\deg(\wc')=\deg(\ec'_{P,\kappa}\times^{P,\Ad}_k \ggo_+)=\deg(\mu(\ec'_{P,\kappa})).$$
By assertion (2) of Lemma \ref{lem:WQ}, we have
$$\deg(\Qc)\leq \deg(\ec_{\bar{P},\kappa}\times^{\bar{P},\Ad}_k \ggo_{-r})=-\deg(\mu(\ec_{\bar{P},\kappa})).$$
We may then conclude, since $\deg(\wc')=-\deg(\Qc)$ (cf. (\ref{eq:degW'=-degQ})).
\end{preuve}

We defined in \S\ref{S:etude-trait} a Levi subgroup $M$. Since $P$ and $\bar{P}$ both contain $M$, it follows from Proposition \ref{prop:Preduc} and from its proof that the torsors $\ec'_{P,\kappa}$ and $\ec_{\bar{P},\kappa}$ are the first factors of triples $m'_{P,\kappa}$ and $m_{\bar{P},\kappa}$ which are reductions to $P$ and $\bar{P}$ of $m'_\kappa$ and $m_\kappa$. Now, by hypothesis, these Hitchin triples are $\xi$-stable. We therefore have, on the one hand,
$$\xi\in -\deg(m'_{P,\kappa})- \,^+\ago_P+\ago_T^M$$
(recall that $\ago_G=\{0\}$ since $G$ is semisimple), whence (cf. the remark following Lemma \ref{lem:mu})
$$\mu(\xi)<-\mu(\deg(m'_{P,\kappa}))=-\deg(\mu(\ec'_{P,\kappa}))$$
and, on the other hand,
$$\xi\in -\deg(m_{\bar{P},\kappa})- \,^+\ago_{\bar{P}}+\ago_T^M$$
which implies the inequality
$$-\mu(\xi)<\mu(\deg(m_{\bar{P},\kappa}))=\deg(\mu(\ec_{\bar{P},\kappa})).$$
We therefore have
$$\deg(\mu(\ec'_{P,\kappa}))<-\mu(\xi)<\deg(\mu(\ec_{\bar{P},\kappa}))$$
which contradicts inequality (\ref{eq:inegmu}) of Lemma \ref{lem:mu}. This was the contradiction sought at the beginning of this paragraph.
\end{paragr}

% ---------- sec10 ----------
\section{$\mc^\xi$ is a Deligne--Mumford Stack}\label{sec:DM}

\begin{paragr}[Statement.] --- \label{S:DMenonce}
The object of this section is to prove the following theorem.

\begin{theoreme}\label{thm:DM}
For every semisimple group $G$ and every parameter $\xi$ in general position (cf. Definition \ref{def:posgenerale}), the stack $\mc^\xi_G$ is a Deligne--Mumford stack.
\end{theoreme}

The proof is to be found in the following paragraph.
\end{paragr}

\begin{paragr}[Proof.] --- \label{S:DMpreuve}
We fix a semisimple group $G$ and a parameter $\xi$ in general position. Let $(\ec,\theta,t)\in \mc^\xi(k)$ and let $\Aut_G(\ec,\theta)$ be the group subscheme of $\Aut_G(\ec)$ which centralizes $\theta$. The functor which associates to a $k$-scheme $S$ the group of sections over $C\times_k S$ of the group scheme $\Aut_G(\ec,\theta)\times_k S$ is represented by a group scheme over $k$. It remains to see that this group is unramified, or again, by a differential criterion, that its tangent space over $k$ is trivial. Now this tangent space admits the following description: it is the space of sections $\varphi\in H^0(C,\Ad(\ec))$ which commute with $\theta$, that is to say such that the bracket $[\theta,\varphi]$ is zero. The vanishing of the tangent space therefore results from the following proposition.

\begin{proposition}\label{prop:phi=0}
For every section $\varphi\in H^0(C,\Ad(\ec))$ such that $[\theta,\varphi]=0$ we have
$$\varphi=0.$$
\end{proposition}

\begin{preuve}
Let $\varphi$ be as above. Note that, by the obvious lemma below, the characteristic of $\varphi$ is constant.

\begin{lemme}\label{lem:H0car}
The space of sections $H^0(C,\car_G)$ is identified with $\car_G(k)$.
\end{lemme}

Let $M\in \lc$ be such that $f(\ec,\theta,t)\in \chi^M_G(\Ac_{M,\el}(k))$ and let $F$ be the function field of the curve $C$.

\begin{lemme}\label{lem:trivgenDM}
There exists a generic trivialization of $\ec$ which satisfies the following two assertions:
\begin{enumerate}
\item the identification of $\Ad_D(\ec)$ with $\ggo(F)$ which results from it sends $\theta$ to an element $X\in \mgo(F)$ which is semisimple, $G$-regular and elliptic in $M$, and which satisfies assertion (1) of Definition \ref{def:Cata};
\item the identification of $\Ad(\ec)$ with $\ggo(F)$ which results from it sends $\varphi$ to an element $Y\in \tgo(k)$.
\end{enumerate}
\end{lemme}

\begin{preuve}
We saw in the proof of Proposition \ref{prop:fibre} that there exists a generic trivialization of $\ec$ which satisfies assertion (1). In the identification of $\Ad(\ec)$ with $\ggo(F)$, the section $\varphi$ is sent to an element of $\ggo(F)$, denoted $Z$. Since $Z$ commutes with $X$ and since the centralizer of $X$ in $\ggo$ is contained in $\mgo$, we deduce that $Z$ belongs to $\mgo(F)$. The characteristic $\chi_M(Z)\in \car_M(F)$ is sent by $\chi^M_G$ to the characteristic of $\varphi$, which by the previous lemma is an element of $\car_G(k)$. It follows that $\chi_M(Z)$ belongs to $\car_M(k)$ and that we can find an element $Y\in \tgo(k)$ such that $\chi_M(Y)=\chi_M(Z)$. Lemma \ref{lem:cjsep} shows that $Y$ and $Z$ are conjugate by an element of $M(F)$. The lemma is then obvious.
\end{preuve}

From now on we fix a generic trivialization of $\ec$ which satisfies Lemma \ref{lem:trivgenDM}. Let $X\in \mgo(F)$ and $Y\in \tgo(k)$ be the elements deduced from it. Let $L$ be the centralizer of $Y$ in $G$. One checks that $L$ is an element of $\lc^G(T)$ which contains $M$. Let $Q$ be a parabolic subgroup of $G$ with Levi $L$. We shall note the following lemma, whose proof is left to the reader.

\begin{lemme}\label{lem:NQiso}
The morphism $N_Q\to \ngo_Q$ given by $n\mapsto \Ad(n^{-1})Y-Y$ is an isomorphism.
\end{lemme}

Let $V$ be the set of closed points of $C$. We resume the notations of Section \ref{sec:description}. From the generic trivialization of the torsor $\ec$ and from the existence of a trivialization on a formal neighborhood of every point $v\in V$, we deduce a family $(g_v)_{v\in V}$ such that:
\begin{enumerate}
\item $g_v\in G((z_v))(k)/G[[z_v]](k)$ is trivial except for finitely many $v$;
\item $\Ad(g_v^{-1})Y\in \ggo[[z_v]](k)$;
\item the projection of $\xi$ onto $\ago^G_M$ belongs to the projection onto $\ago^G_M$ of the convex hull of the points
$$-H_P((g_v)_{v\in V})=-\sum_{v\in V}H_P(g_v).$$
\end{enumerate}
Condition (2) expresses the fact that $\varphi$ is a global section of $\Ad(\ec)$, and condition (3) the $\xi$-stability of the triple $(\ec,\theta,t)$ (cf. condition (d) of Definition \ref{def:Cata}).

\begin{lemme}\label{lem:relevL}
For every $v\in V$, the class $g_v$ lifts to an element of $L((z_v))(k)$.
\end{lemme}

\begin{preuve}
Let $Q$ be a parabolic subgroup of $G$ with Levi $L$. By the Iwasawa decomposition, the element $g_v$ lifts to an element $l_vn_v$ with $l_v\in L((z_v))(k)$ and $n_v\in N_Q((z_v))(k)$. The fact that $L$ centralizes $Y$ together with relation (1) above then implies that
$$\Ad(n_v^{-1})Y\in \ggo[[z_v]](k)$$
whence
$$\Ad(n_v^{-1})Y-Y\in \ngo_Q[[z_v]](k).$$
By Lemma \ref{lem:NQiso}, we then have $n_v\in N_Q[[z_v]](k)$. Whence the lemma.
\end{preuve}

The previous lemma implies that the triple $(\ec,\theta,t)$ possesses a ``reduction'' to the Levi subgroup $L$. It also implies that the projection onto $\ago^G_L$ of the convex hull in 3 above is reduced to a point, which is necessarily a point of the lattice $X_*(L)$ defined in \S\ref{S:espacesa}. The $\xi$-stability entails that the projection of $\xi$ onto $\ago^G_L$ belongs to this lattice. If $L\neq G$, this contradicts the fact that $\xi$ is in general position. We therefore have $L=G$ and $Y=0$, whence $\varphi=0$.
\end{preuve}
\end{paragr}

% ---------- sec11 ----------
\section{Counting Rational Points in a Fiber}\label{sec:comptage}

\begin{paragr}[The Main Theorem.] --- \label{S:thm-comptage}
Here is the situation that will occupy us until the end of this section. Recall that in \S \ref{S:lacourbe} we fixed a curve $C$ over $k$ as well as a divisor $D$ and a point $\infty\in C(k)$. We assume that $C$ arises by base change from a curve $C_0$ over $\Fq$. Let $\tau$ be the Frobenius automorphism of $k$ given by $x\mapsto x^q$. We again denote by $\tau$ the automorphism $1\times\tau$ of $C=C_0\times_{\Fq}k$. We assume that $D$ and $\infty$ are fixed under $\tau$.

Let $F$ be the function field of $C$. The automorphism $\tau$ of $C$ induces an automorphism of $F$, still denoted by $\tau$, whose set of fixed points is precisely the field $F_0$ of the curve $C_0$. Let $V$, respectively $V_0$, be the set of closed points of $C$, respectively of $C_0$. For every $v_0\in V_0$, let $F_{0,v_0}$ be the completion of $F_0$ at $v_0$; we simply denote by $\tau$ the automorphism $\tau\otimes 1$ of
$$F\,\hat{\otimes}_{F_0} F_{0,v_0}=\prod_{v\in V,\ v|v_0} F_v$$
where $F_v$ is the completion of $F$ at $v$. We identify $F_v$ with $k((z_v))$ and $F_{0,v_0}$ with $\Fq((z_{v_0}))$ by the choice of uniformizers. Let $\AAA$ be the ring of adeles of $F$. This ring is naturally equipped with an action of $\tau$ for which the subring of points fixed under $\tau$ is identified with the ring $\AAA_0$ of adeles of $F_0$.

We assume that the group $G$ is \emph{semi-simple} and that the pair $(G,T)$ (cf. \S \ref{sec:morcar}) arises from an analogous pair $(G_0,T_0)$ defined over $\Fq$ such that the torus $T_0$ is split over $\Fq$. In this case, all the Levi or parabolic subgroups of $G$ which contain $T$ arise from analogous objects defined over $\Fq$ relative to $(G_0,T_0)$. We again denote by $\tau$ the automorphism of $G=G_0\times_{\Fq}k$ given by $1\times\tau$. From this we deduce an automorphism, still denoted by $\tau$, of the groups of points of $G$ with values in $F$, $\AAA$ or $F\otimes_{F_0}F_{0,v_0}$ for every $v_0\in V_0$.

Since the divisor $D$ on $C$ is fixed by $\tau$, it descends to a divisor $D_0$ on $C_0$ of the form
$$\sum_{v\in V_0} d_v v.$$
Let $\mathbf{1}_{D_0}$ be the function on $\ggo(\AAA_0)$ which is the characteristic function of the set
$$\prod_{v\in V_0} z_v^{-d_v}\ggo[[z_v]](\Fq).$$
To lighten the notation a little, we set
$$K=\prod_{v\in V_0} G[[z_v]](\Fq),$$
which is an open and compact subgroup of $G(\AAA_0)$. Moreover, the function $\mathbf{1}_{D_0}$ is invariant under the adjoint action of $K$.

For every $\delta\in M(F)$, let $M^\delta(F_0)\subset M(F)$, respectively $\mgo^\delta(F_0)\subset\mgo(F)$, be the subgroup of points fixed under $\Int(\delta)\circ\tau$, respectively the Lie subalgebra of points fixed under $\Ad(\delta)\circ\tau$.

The Frobenius automorphism $\tau$ also acts on the spaces $\Ac$ and $\Ac_M$ for every $M\in\lc$. Let $(a,t)\in\Ac(k)$ be an element fixed by $\tau$. There exist $M\in\lc$ and $(a_M,t)\in\Ac_{M,\el}(k)$ such that
$$(a,t)=\chi_G^M((a_M,t)).$$
By Propositions \ref{prop:immersionfermee} and \ref{prop:reunionAM}, such elements exist and are unique: in particular they are fixed under $\tau$.

Here is the main theorem of this section. It expresses the number of rational points of a fiber of $\overline{f^\xi}$ in terms of \emph{weighted orbital integrals} of Arthur.

\begin{theoreme}\label{thm:comptage}
Let $\xi$ be in general position. Let $M\in\lc$ be a Levi subgroup, $(a_M,t)\in\Ac_{M,\el}(k)$ fixed by $\tau$ and $(a,t)=\chi_G^M((a_M,t))$. The cardinality of the groupoid of rational points of the fiber of $\overline{f^\xi}$ at $(a,t)$ is equal to
$$\vol(\ago_M/X_*(M))^{-1}\cdot\vol(a,t)\cdot\sum_h\ \sum_X\ J_M^G(\Ad(h^{-1})X,\mathbf{1}_{D_0})$$
where:
\begin{itemize}
\item $h$ runs through the set of double cosets
$$h\in M(F)\backslash M(\AAA)/M(\AAA_0)$$
which satisfy
$$\delta=h\tau(h)^{-1}\in M(F);$$
\item $X$ runs through a system of representatives of the set of $M^\delta(F_0)$-conjugacy classes in the set
$$\{X\in\mgo^\delta(F_0)\mid \chi_M(X)=a_{M|\Spec(F_0)}\}$$
of those $X\in\mgo^\delta(F_0)$ whose characteristic $\chi_M(X)$ is equal to the restriction of $a_M$ to $\Spec(F_0)$;
\item $J_M^G(\Ad(h^{-1})X,\mathbf{1}_{D_0})$ is the weighted orbital integral of Arthur defined in (\ref{eq:JM});
\item the volume $\vol(a,t)$ is defined in \S \ref{S:mesures-haar}, line (\ref{eq:volat}).
\end{itemize}
\end{theoreme}

\begin{remarques}
The weighted orbital integrals considered above depend on choices of Haar measures on $\ago_M$ (cf. \S \ref{S:poids-arthur}) and on certain tori (cf. \S \ref{S:mesures-haar}). Since these choices also enter into the factors $\vol(\ago_M/X_*(M))$ and $\vol(a,t)$, the sum that appears in Theorem \ref{thm:comptage} does not depend on any choice.

The cardinality of a $\xi$-stable Hitchin fiber does not depend on the choice of generic $\xi$. As one of the referees of this article notes, it would be interesting to have an explanation of a geometric nature. It is clear that the notion of $\xi$-stability depends only on the connected component of $\xi$ in the ``generic'' open set. Ng\^{o} has constructed an action of a Picard stack on a Hitchin fiber $\mc_a$ (cf. \cite{Ngo06}). Via a degree morphism, this stack also acts on $\ago_T$ in such a way that if the action sends $\xi$ to $\xi'$, the truncation $\mc_a^\xi$ is sent to $\mc_a^{\xi'}$. However, the action of this stack is in general not transitive on the set of connected components of the generic open set. One would have to analyze what happens when one ``crosses a wall'' (this is probably related to the works \cite{DH98} and \cite{Tha96}).

When the cohomology set $H^1(\Gal(F_0^{\mathrm{sep}}/F_0),G(F_0^{\mathrm{sep}}))$ is trivial (where $F_0^{\mathrm{sep}}$ is a separable closure of $F_0$), for example when the group $G$ is semi-simple simply connected, one can show that the set of double cosets $h$ reduces to a single element. In particular, Theorem \ref{thm:comptage} does indeed give Theorem \ref{thm:slncomptage} of the introduction. The proof of Theorem \ref{thm:comptage} rests on the adelic description of a Hitchin fiber, of which one takes the points ``fixed by Frobenius''. When the above $H^1$ is trivial, the Hitchin triples over the curve $C_0$ are generically trivial and, for the counting, one could just as well use an adelic description of the Hitchin triples over $C_0$.
\end{remarques}

Before beginning the proof of Theorem \ref{thm:comptage}, which will occupy us until the end of this section, we shall give some reminders on the cardinality of a quotient groupoid and on its fixed points under an automorphism.
\end{paragr}

\begin{paragr}[Fixed Points Under an Automorphism.] --- \label{S:points-fixes}
Let $\xc$ be a set and $G$ an abstract group acting on the left on $\xc$. Let $\tau$ be an automorphism of $G$. Let there be a bijection of $\xc$, again denoted by $\tau$, compatible with the automorphism of $G$ in the sense that one has
$$\tau(g\cdot x)=\tau(g)\cdot\tau(x).$$
By definition, the groupoid of points of $[G\backslash\xc]$ fixed under $\tau$ is the groupoid $[G\backslash\xc^\tau]$ where:
\begin{itemize}
\item $\xc^\tau$ is the set of pairs $(x,g)\in\xc\times G$ such that
$$g\cdot\tau(x)=x;$$
\item $G$ acts on the left on $\xc^\tau$ by
$$h\cdot(x,g)=(hx,hg\tau(h)^{-1}).$$
\end{itemize}
\end{paragr}

\begin{paragr}[Cardinality of a Groupoid.] --- \label{S:cardinal-groupoide}
Let $\xc$ be a set and $G$ an abstract group acting on the left on $\xc$. We assume that for every $x\in\xc$, the stabilizer $\stab_G(x)$ of $x$ in $G$ is finite. We also assume that the set $G\backslash\xc$ of orbits of $\xc$ under $G$ is finite. The cardinality of the quotient groupoid $[G\backslash\xc]$ (cf. (\ref{eq:groupoidequot}) of \S \ref{S:description-adelic}) is by definition
$$\card([G\backslash\xc])=\sum_{x\in G\backslash\xc}\frac{1}{|\stab_G(x)|}.$$
\end{paragr}

\begin{paragr}[First Step.] --- \label{S:premiere-etape}
In what follows, we fix $M\in\lc$ and pairs $(a_M,t)$ and $(a,t)$ as in the statement of Theorem \ref{thm:comptage}. Let $\xi\in\ago_T$ be an \emph{arbitrary} element. The fiber of the Hitchin morphism $\overline{f^\xi}$ above $(a,t)$ is identified, by Proposition \ref{prop:fibre}, with the quotient groupoid $[M(F)\backslash\xc^\xi_{(a,t)}]$ where the set $\xc^\xi_{(a,t)}$ is the one of Definition \ref{def:Cata} of paragraph \ref{S:description-adelic}. Our hypotheses imply that $\tau$ induces a bijection of this set compatible with the action of $\tau$ on $M(F)$ (in the sense of \S \ref{S:points-fixes}). We are looking for an expression for the number of points over $\Fq$ of this fiber, that is to say the cardinality of the groupoid of points of $[M(F)\backslash\xc^\xi_{(a,t)}]$ fixed under $\tau$, in terms of weighted orbital integrals. Since the Hitchin morphism $\overline{f^\xi}$ is of finite type (cf. Proposition \ref{prop:schptf}), the required finiteness hypotheses (cf. \S \ref{S:cardinal-groupoide}) are indeed satisfied.

Following paragraph \ref{S:points-fixes}, we introduce the set $\xc^{\xi,\tau}_{(a,t)}$ formed of the triples $(X,(g_v)_{v\in V},\delta)$ such that $(X,(g_v)_{v\in V})\in\xc^\xi_{(a,t)}$ and $\delta\in M(F)$ satisfy the relations
\begin{equation}
\label{eq:Adtau}
\Ad(\delta)\tau(X)=X
\end{equation}
and
\begin{equation}
\label{eq:deltatau}
\delta\tau((g_v)_{v\in V})=(g_v)_{v\in V}.
\end{equation}

The group $M(F)$ acts on the first two factors by the action described in \S \ref{S:description-adelic} and by $\tau$-conjugation on the second factor (that is to say an element $\gamma\in M(F)$ acts on $\delta\in M(F)$ by $\gamma\delta\tau(\gamma)^{-1}$). The first step of the counting consists in describing a system of representatives of the orbits of $\xc^{\xi,\tau}_{(a,t)}$ under the action of $M(F)$. To this end, we introduce a few additional notations.

For every $g\in G(\AAA)$, let $\mathbf{1}_{M,g}$ be the function on $\ago_M$ which is the characteristic function of the convex hull of the points $-H_P(g)$ for $P\in\pc(M)$. Definition \ref{def:HP} of \S \ref{S:HP}, applied to the group $M$ and to its parabolic subgroup $P=M$, gives a morphism
$$H_M:M(\AAA)\to\ago_M.$$
Let $\xi_M$ be the projection of $\xi$ onto $\ago_M$ according to the decomposition $\ago_T=\ago_M\oplus\ago_T^M$.

\begin{lemme}\label{lem:representants}
Consider the following sets:
\begin{enumerate}
\item a system of representatives in $M(\AAA)$ of the set of double cosets
$$h\in M(F)\backslash M(\AAA)/M(\AAA_0)$$
which satisfy
$$h\tau(h)^{-1}\in M(F);$$
\item for every $\delta\in M(F)$, a system of representatives in $\mgo^\delta(F_0)$ of the set of $M^\delta(F_0)$-conjugacy classes in the set
$$\{X\in\mgo^\delta(F_0)\mid\chi_M(X)=a_{M|\Spec(F_0)}\}$$
of those $X\in\mgo^\delta(F_0)$ whose characteristic $\chi_M(X)$ is equal to the restriction of $a_M$ to $\Spec(F_0)$;
\item for all $h\in M(\AAA)$ and $\delta\in M(F)$ such that $\delta=h\tau(h)^{-1}$ and every semi-simple $G$-regular $X\in\mgo^\delta(F_0)$, a system of representatives in $G(\AAA_0)$ of the set of double cosets
$$\Int(h^{-1})T_X(F_0)\backslash G(\AAA_0)/K,$$
where $T_X(F_0)$ is the centralizer of $X$ in $M^\delta(F_0)$, which satisfy the two conditions:
\begin{enumerate}
\item[(a)] $\mathbf{1}_{D_0}(\Ad(g^{-1})\Ad(h^{-1})X)=1$;
\item[(b)] $\mathbf{1}_{M,g}(\xi_M+H_M(h))=1$.
\end{enumerate}
\end{enumerate}
The map
$$(X,g,h)\mapsto(X,hgG(\oc),h\tau(h)^{-1})$$
from the set formed of the triples $(X,g,h)$ such that $h$, $X$ and $g$ belong to the sets described respectively in (1), (2) and (3) (for $\delta=h\tau(h)^{-1}$) to $\xc^{\xi,\tau}_{(a,t)}$ is injective and its image is a system of representatives of the set of classes $M(F)\backslash\xc^{\xi,\tau}_{(a,t)}$.

Moreover, the cardinality of the stabilizer in $M(F)$ of the image of $(X,g,h)$ is given by the cardinality of the finite set
$$(\Int(h^{-1})T_X(F_0))\cap gKg^{-1}.$$
\end{lemme}

\begin{preuve}
We must show that the map in question induces a bijective map on the quotient $M(F)\backslash\xc^{\xi,\tau}_{(a,t)}$. Let us first show surjectivity. Let us start from a triple $(X,(g_v)_{v\in V},\delta)$ in $\xc^{\xi,\tau}_{(a,t)}$. Recall that $g_v$ denotes a class in $G((z_v))(k)/G[[z_v]](k)$. Let $g\in G(\AAA)$ be an element such that for every $v$ the component at $v$ lifts $g_v$. By relation (\ref{eq:deltatau}) above, one has
$$g^{-1}\delta\tau(g)\in\prod_{v\in V}G[[z_v]](k).$$
By Lang's theorem (applied to the (connected) quotients of the pro-algebraic $k$-group $G[[z_v]]$), one knows that such an element can be written $x\tau(x^{-1})$ for a certain $x\in\prod_{v\in V}G[[z_v]](k)$. After possibly replacing $g$ by $gx$, we may and will assume that one has
\begin{equation}
\label{eq:tauconj}
g^{-1}\delta\tau(g)=1
\end{equation}
in other words the $\tau$-conjugacy class of $\delta$ in $G(\AAA)$ is trivial. This implies that the $\tau$-conjugacy class of $\delta$ in $M(\AAA)$ is trivial (we have not found a reference: an argument will be given in the following article). There therefore exists $h\in M(\AAA)$ such that
$$\delta=h\tau(h)^{-1}.$$
After possibly letting $M(F)$ act on the triple $(X,(g_v)_{v\in V},\delta)$, one may assume that $h$ belongs to the set described in (1).

By relation (\ref{eq:Adtau}), the element $X$ belongs to $\mgo^\delta(F_0)$. Let $t_a\in\tgo^{\reg}[[z_\infty]](k)$ be the unique lift of $t$ whose characteristic is $a$. In the same way, one can define a point $t_{a_M}$ associated with the pair $(a_M,t)$. But by uniqueness of $t_a$, one has $t_a=t_{a_M}$. Thus the characteristic $\chi_M(t_a)$ is $a_M$. Condition (1)(bis) of Definition \ref{def:Cata2} implies that $X$ is conjugate to $t_a$ by an element of $M((z_\infty))(k)$, whence
$$\chi_M(X)=\chi_M(t_a)=a_M.$$
After possibly letting $M^\delta(F_0)$ act on the triple $(X,(g_v)_{v\in V},\delta)$, which changes neither $\delta$ nor condition (\ref{eq:tauconj}), one may assume that $X$ belongs to the set described in (2).

Let us set $g_0=h^{-1}g$. Then relation (\ref{eq:tauconj}) translates into
$$g_0\in G(\AAA_0).$$
After possibly letting $T_X(F_0)$ (which is contained in $M^\delta(F_0)$) act on the triple $(X,(g_v)_{v\in V},\delta)$, one may assume moreover that $g_0$ belongs to the set described in (3). Conditions (2)(b) and (2)(d) of Definitions \ref{def:Cata} and \ref{def:Cata2} translate into conditions (3)(a) and (b) above. We have thus proved the desired surjectivity. We leave it to the reader to verify injectivity.

The centralizer of $(X,(g_v)_{v\in V},\delta)=(X,hgG(\oc),h\tau(h)^{-1})$ in $M(F)$ is the group
$$\{t\in T_X(F_0)\mid (tg_v)_{v\in V}=(g_v)_{v\in V}\}=T_X(F_0)\cap\prod_{v\in V}g_vG[[z_v]](k)g_v^{-1}$$
$$\simeq(\Int(h^{-1})T_X(F_0))\cap gKg^{-1}$$
which gives the last assertion.
\end{preuve}
\end{paragr}

\begin{paragr}[Haar Measures.] --- \label{S:mesures-haar}
We equip $G(\AAA_0)$ with the Haar measure normalized by
$$\vol(K)=1.$$
More generally, for every semi-standard parabolic subgroup $P$ of $G$, the groups $M_P(\AAA_0)$ and $N_P(\AAA_0)$ are equipped with the Haar measures normalized by
$$\vol(M_P(\AAA_0)\cap K)=1\quad\text{and}\quad\vol(N_P(\AAA_0)\cap K)=1.$$

Let $X_0\in\mgo(F_0)$ whose characteristic is equal to the restriction of $a_M$ to $\Spec(F_0)$. Such an $X_0$ exists by the Kostant section relative to $M$. We equip $T_{X_0}(\AAA_0)$ with a Haar measure. Let
$$T_{X_0}(\AAA_0)^1=T_{X_0}(\AAA_0)\cap\Ker(H_M);$$
this is an open subgroup of $T_{X_0}(\AAA_0)$ which we equip with the induced measure. We equip $T_{X_0}(F_0)$ with the counting measure and $T_{X_0}(F_0)\backslash T_{X_0}(\AAA_0)^1$ with the quotient measure.

\begin{lemme}\label{lem:volfini}
The volume of the quotient below is finite
\begin{equation}
\label{eq:volfini}
\vol(T_{X_0}(F_0)\backslash T_{X_0}(\AAA_0)^1)<\infty.
\end{equation}
\end{lemme}

\begin{preuve}
By Corollary \ref{cor:caractdesAM}, there exists an element $X\in\mgo(F)$ which is semi-simple, $G$-regular and elliptic in $\mgo(F)$. By Lemma \ref{lem:cjsep}, $X$ and $X_0$ are conjugate under $M(F)$. One deduces from this that the largest $F_0$-split subtorus of $T_{X_0}$ is contained in the connected center of $M\times F_0$ (and in fact equal to it). Thus $T_{X_0}$ is $F_0$-elliptic in $M$ and it is well known that the quotient $T_{X_0}(F_0)\backslash T_{X_0}(\AAA_0)^1$ is then compact (cf. for example \cite[Theorem 4.11]{Spr79}), whence the finiteness of the volume.
\end{preuve}

Let $\delta\in M(F)$ and $h\in M(\AAA)$ be such that $\delta=h\tau(h)^{-1}$. Let $X\in\mgo^\delta(F_0)$ of characteristic $a_M$. The condition on the characteristic implies that there exists $m\in M(F)$ such that $X=\Ad(m^{-1})X_0$ (cf. Lemma \ref{lem:cjsep}). Let
$$Y=\Ad(h^{-1})X=\Ad((hm)^{-1})X_0.$$
One notes that $Y\in\mgo(\AAA_0)$. Let $T_Y(\AAA_0)$ be the centralizer of $Y$ in $M(\AAA_0)$. The automorphism $\Int(hm)^{-1}$ induces an isomorphism from $T_{X_0}(\AAA_0)$ onto $T_Y(\AAA_0)$. We equip the latter with the Haar measure obtained by transport under this automorphism. Let us set, somewhat abusively,
$$T_Y(F_0)=\Int(hm)^{-1}T_{X_0}(F_0).$$
This group is discrete and one has, by our choices of measures,
\begin{equation}
\label{eq:volTY}
\vol(T_Y(F_0)\backslash T_Y(\AAA_0)^1)=\vol(T_{X_0}(F_0)\backslash T_{X_0}(\AAA_0)^1),
\end{equation}
where we set $T_Y(\AAA_0)^1=T_Y(\AAA_0)\cap\Ker(H_M)$. We denote by
\begin{equation}
\label{eq:volat}
\vol(a,t)
\end{equation}
the above volume.
\end{paragr}

\begin{paragr}[Lattices in $\ago_M$.] --- \label{S:reseaux}
Let $G_{\scnx}$ be the simply connected covering of $G$. Let $L\in\lc^G(T)$ be a semi-standard Levi subgroup. Let $L_{\scnx}$ be the inverse image of $L$ in $G_{\scnx}$. Let $L_{\der}$ be the derived group of $L$ and $L_{\SCNX}$ the simply connected covering of $L_{\der}$. One should not confuse $L_{\scnx}$ and $L_{\SCNX}$. The group $L_{\SCNX}$ is in fact the derived group of $L_{\scnx}$. We again denote by $T_?$ the torus obtained as the inverse image of $T$ in $?$, where $?$ may be any one of the three groups $L_{\der}$, $L_{\SCNX}$ or $L_{\scnx}$. Let
$$\pi_1(L_{\der})=\mathrm{coker}(X_*(T_{L_{\SCNX}})\to X_*(T_{L_{\der}})).$$

The restriction morphism $X^*(L)\to X^*(T)$ gives dually a morphism $X_*(T)\to X_*(L)$, where we have set $X_*(L)=\Hom_\ZZ(X^*(L),\ZZ)$.

\begin{lemme}\label{lem:diagramme}
One has a commutative diagram with exact rows and columns.
$$\xymatrix{
& 0\ar[d] & 0\ar[d] & 0\ar[d] & \\
0\ar[r] & X_*(T_{L_{\SCNX}})\ar[d]\ar[r] & X_*(T_{G_{\scnx}})\ar[d]\ar[r] & X_*(L_{\scnx})\ar[d]\ar[r] & 0\\
0\ar[r] & X_*(T_{L_{\der}})\ar[d]\ar[r] & X_*(T)\ar[d]\ar[r] & X_*(L)\ar[d]\ar[r] & 0\\
0\ar[r] & \pi_1(L_{\der})\ar[d]\ar[r] & \pi_1(G)\ar[d]\ar[r] & X_*(L)/X_*(L_{\scnx})\ar[d]\ar[r] & 0\\
& 0 & 0 & 0 &
}$$
\end{lemme}

\begin{preuve}
Let us treat the exactness of the second row. It is a matter of a sequence of free $\ZZ$-modules whose dual sequence is the exact sequence
$$0\longrightarrow X^*(L)\longrightarrow X^*(T)\longrightarrow X^*(T_{L_{\der}})\longrightarrow 0.$$
Only surjectivity is less obvious. It follows from the isomorphisms
$$T/T_{L_{\der}}=Z_L^0/(Z_L^0\cap L_{\der})=L/L_{\der},$$
where $Z_L^0$ is the connected center of $L$, and from the equality $X^*(L)=X^*(L_{\der})$.

The exactness of the first row is deduced from that of the second upon replacing $L$ by $L_{\scnx}$. The injectivity of the first two vertical arrows is well known. We recall that $X_*(T_{G_{\scnx}})$ is the sub-$\ZZ$-module of $X_*(T)$ generated by the coroots of $T$ in $G$. The rest of the diagram then follows from the snake lemma and from the injectivity of
$$\pi_1(L_{\der})\longrightarrow\pi_1(G)$$
that is to say from the inclusion
$$X_*(T_{L_{\der}})\cap X_*(T_{G_{\scnx}})\subset X_*(T_{L_{\SCNX}}).$$
This last is obvious since $X_*(T_{L_{\SCNX}})$ is of finite index in $X_*(T_{L_{\der}})$ and since on the other hand $X_*(T_{L_{\SCNX}})$ is a direct factor in $X_*(T_{G_{\scnx}})$.
\end{preuve}

Lemma \ref{lem:diagramme} shows that the projection $\ago_T\to\ago_M$ maps $X_*(T)$ surjectively onto $X_*(M)$. On the other hand, it maps the sub-$\ZZ$-module $X_*(T_{G_{\scnx}})$ surjectively onto the sub-$\ZZ$-module $X_*(M_{\scnx})$ of $X_*(M)$.
\end{paragr}

\begin{paragr}[The Weight $\mathrm{w}_M^\xi$.] --- \label{S:poids-wM}
For every $g\in G(\AAA)$, we introduce the weight
$$\mathrm{w}_M^\xi(g)=|\{\mu\in X_*(M)\mid\mathbf{1}_{M,g}(\xi_M+\mu)=1\}|$$
which is the (finite) number of points of the set $\xi_M+X_*(M)$ which belong to the convex hull of the points $-H_P(g)$ for $P\in\pc(M)$.

\begin{lemme}\label{lem:poids-invariant}
The weight function $g\in G(\AAA)\mapsto\mathrm{w}_M^\xi(g)$ is left invariant under $M(\AAA)$.

Let $h$ and $X$ be two elements respectively of the sets described in (1) and (2) of Lemma \ref{lem:representants} for $\delta=h\tau(h)^{-1}$. Let $Y=\Ad(h)^{-1}X$. One has an exact sequence (with the notations of \S \ref{S:mesures-haar})
$$1\longrightarrow T_Y(\AAA_0)^1\longrightarrow T_Y(\AAA_0)\stackrel{H_M}{\longrightarrow}X_*(M)\longrightarrow 0.$$
\end{lemme}

\begin{preuve}
The left action of $h\in M(\AAA)$ translates the convex hull of the points $-H_P(g)$ by the vector $H_M(h)$. Since the latter belongs to $X_*(M)$, this does not change $\mathrm{w}_M^\xi(g)$.

In the exact sequence, only the surjectivity on the right is not obvious. Let $X_0\in\mgo(F_0)$ of characteristic $a_M$. Since the map $H_M$ is invariant under the adjoint action of $M(\AAA)$ and since there exists an element of $M(\AAA)$ which conjugates $Y$ and $X_0$ on the one hand and $T_Y(\AAA_0)$ and $T_{X_0}(\AAA_0)$ on the other hand (cf. \S \ref{S:mesures-haar}), one may assume that one has $Y=X_0$. Combining Lemma \ref{lem:cjsep} and Proposition \ref{prop:caractdesAM}, one shows that $X_0$ and $t_a$ are conjugate by $m\in M((z_\infty))(k)$ ($t_a\in\tgo^{\reg}[[z_\infty]](k)$ being the element defined in Proposition \ref{prop:caractdesAM}). One will note that one even has $t_a\in\tgo^{\reg}[[z_\infty]](\Fq)$. It follows that the automorphism $\Int(m)$ then induces an isomorphism between $T_{X_0}(\Fq((z_\infty)))$ and $T(\Fq((z_\infty)))$. Since the map $H_M$ restricted to $T(\Fq((z_\infty)))$ is surjective, the same holds for its restriction to $T_{X_0}(\Fq((z_\infty)))$ and \emph{a fortiori} to $T_{X_0}(\AAA_0)$.
\end{preuve}
\end{paragr}

\begin{paragr}[Weighted Orbital Integrals $J_M^\xi$.] --- \label{S:JMxi}
Let $h$ and $X$ be two elements respectively of the sets described in (1) and (2) of Lemma \ref{lem:representants} for $\delta=h\tau(h)^{-1}$. Let $Y=\Ad(h^{-1})X$. We have equipped in \S \ref{S:mesures-haar} the centralizer $T_Y(\AAA_0)$ of $Y$ in $G(\AAA_0)$ with a Haar measure. We then introduce the following weighted orbital integral
\begin{equation}
\label{eq:JMxi}
J_M^\xi(Y,\mathbf{1}_{D_0})=\int_{T_Y(\AAA_0)\backslash G(\AAA_0)}\mathbf{1}_{D_0}(\Ad(g^{-1})Y)\mathrm{w}_M^\xi(g)\,dg.
\end{equation}
We have equipped $T_Y(\AAA_0)\backslash G(\AAA_0)$ with the quotient measure. The integral converges since the integrand has compact support. To lighten the notation, we omit the function $\mathbf{1}_{D_0}$ and we set
$$J_M^\xi(Y)=J_M^\xi(Y,\mathbf{1}_{D_0}).$$

\begin{remarque}
This weighted orbital integral is not the one that Arthur usually uses. The link with Arthur's will be made in \S \ref{S:comparaison}.
\end{remarque}
\end{paragr}

\begin{paragr}[A First Count.] --- \label{S:premier-comptage}
Here is a first expression for the cardinality of a fiber.

\begin{proposition}\label{prop:premier-comptage}
One has the equality
$$\card([M(F)\backslash\xc^{\xi,\tau}_{(a,t)}])=\vol(a,t)\sum_h\ \sum_X\ J_M^\xi(\Ad(h^{-1})X)$$
where $h$ and $X$ run respectively through the sets described in (1) and (2) of Lemma \ref{lem:representants} for $\delta=h\tau(h)^{-1}$.
\end{proposition}

\begin{preuve}
By the formula given in \S \ref{S:cardinal-groupoide} and Lemma \ref{lem:representants}, the cardinality of the groupoid $[M(F)\backslash\xc^{\xi,\tau}_{(a,t)}]$ can be written as the sum over the elements $h$ and $X$ of the sets described respectively in (1) and (2) of Lemma \ref{lem:representants} (for $\delta=h\tau(h)^{-1}$) of
$$\sum_{g\in T_Y(F_0)\backslash G(\AAA_0)/K}\frac{1}{|T_Y(F_0)\cap gKg^{-1}|}\mathbf{1}_{D_0}(\Ad(g^{-1})Y)\mathbf{1}_{M,g}(\xi_M+H_M(h)),$$
where we have set $Y=\Ad(h^{-1})X$ and $T_Y(F_0)$ is the group defined in \S \ref{S:mesures-haar}. By our choice of measure on $G(\AAA_0)$, this expression can also be written
$$\int_{T_Y(F_0)\backslash G(\AAA_0)}\mathbf{1}_{D_0}(\Ad(g^{-1})Y)\mathbf{1}_{M,g}(\xi_M+H_M(h))\,dg,$$
where the centralizer $T_Y(\AAA_0)$ of $Y$ in $G(\AAA_0)$ is equipped with the Haar measure of \S \ref{S:mesures-haar}. Since $g\mapsto\mathbf{1}_{D_0}(\Ad(g^{-1})Y)$ is right invariant under $T_Y(\AAA_0)$, the preceding expression can be written
$$\int_{T_Y(\AAA_0)\backslash G(\AAA_0)}\mathbf{1}_{D_0}(\Ad(g^{-1})Y)\int_{T_Y(F_0)\backslash T_Y(\AAA_0)}\mathbf{1}_{M,tg}(\xi_M+H_M(h))\,dt\,\frac{dg}{dt}.$$
Using the equality
$$\mathbf{1}_{M,tg}(\xi_M+H_M(h))=\mathbf{1}_{M,g}(\xi_M+H_M(h)+H_M(t)),$$
our choices of measures and Lemma \ref{lem:poids-invariant}, one sees that one has
$$\int_{T_Y(F_0)\backslash T_Y(\AAA_0)}\mathbf{1}_{M,tg}(\xi_M+H_M(h))\,dt=\vol(T_Y(F_0)\backslash T_Y(\AAA_0)^1)\cdot\mathrm{w}_M^\xi(g).$$
The result follows.
\end{preuve}
\end{paragr}

\begin{paragr}[A Formula for the Weight $\mathrm{w}_M^\xi$.] --- \label{S:formule-poids}
The aim of this section is to give an analytic formula for the weight $\mathrm{w}_M^\xi$. This will then allow us to compare this weight with the one Arthur considers in his works. For this, we follow, up to a few details, Arthur (cf. \cite[\S 6]{Art91}). We introduced in \S \ref{S:reseaux} the sub-$\ZZ$-module $X_*(M_{\scnx})$ of $X_*(M)$. Let $P\in\pc(M)$. The set $\Delta_P^\vee$ defined in \S \ref{S:deltachapeau} provides a basis of $X_*(M_{\scnx})$. Every $\la\in\ago_M$ can be written in a unique way
\begin{equation}
\label{eq:decomplambda}
\la=[\la]_P+\{\la\}_P
\end{equation}
with $[\la]_P\in X_*(M_{\scnx})$ and
$$\{\la\}_P=\sum_{\al\in\Delta_P}r_\al\al^\vee$$
with $0<r_\la\leq 1$.

For every $\Lambda\in\ago_M^\ast$, let
$$\mathrm{c}_P(\Lambda)=\prod_{\al\in\Delta_P}(\exp(\Lambda(\al^\vee))-1).$$
When $M=G$, this product equals $1$ (by convention every product over the empty set equals $1$). In what follows, we say that $\Lambda\in\ago_M^\ast$ is \emph{generic} if $\mathrm{c}_P(\Lambda)\neq 0$ for every $P\in\pc(M)$, that is to say $\Lambda(\al^\vee)\neq 0$ for all $P\in\pc(M)$ and $\al^\vee\in\Delta_P^\vee$.

\begin{proposition}\label{prop:formule-poids}
For every $g\in G(\AAA)$, one has the following equality
$$\mathrm{w}_M^\xi(g)=\lim_{\Lambda\to 0}\sum_{\mu_0\in X_*(M)/X_*(M_{\scnx})}\ \sum_{P\in\pc(M)}\mathrm{c}_P(\Lambda)^{-1}\exp(-\Lambda(H_P(g)+[\mu_0+\xi_M]_P))$$
where the limit is taken over the generic $\Lambda\in\ago_M^\ast$.
\end{proposition}

\begin{remarque}
In the limit, the inner sum depends on the choice of a system of representatives of $X_*(M)/X_*(M_{\scnx})$. Changing $\mu_0$ into $\mu_0+\mu$ with $\mu\in X_*(M_{\scnx})$ multiplies the inner sum by $\exp(-\Lambda(\mu))$. This therefore does not affect the limit at $\Lambda=0$.
\end{remarque}

\begin{preuve}
It rests entirely on Arthur's methods. Let us give a few indications for the convenience of the reader. Let $(\varpi_\al)_{\al\in\Delta_P}$ be the basis of $\ago_M^\ast$ dual to $\Delta_P^\vee$. For every generic $\Lambda\in\ago_M^\ast$ (in the sense defined above) and every $P\in\pc(M)$, let
$$\Delta_P^\Lambda=\{\al\in\Delta_P\mid\Lambda(\al^\vee)<0\}$$
and let $\varphi_P^\Lambda$ be the characteristic function of those $\la\in\ago_M$ such that $\varpi_\al(\la)>0$ for all $\al\in\Delta_P^\Lambda$ and $\varpi_\al(\la)\leq 0$ for $\al\in\Delta_P-\Delta_P^\Lambda$. Let $g\in G(\AAA)$. By a lemma due to Langlands (cf. \cite[formula (3.8), p.~22]{Art91}), the characteristic function of the convex hull of the points $-H_P(g)$ is equal to the function
$$\mu\mapsto\sum_{P\in\pc(M)}(-1)^{|\Delta_P^\Lambda|}\varphi_P^\Lambda(\mu+H_P(g))$$
of the variable $\mu\in\ago_M$. This formula is valid in our context since the family $(-H_P(g))_{P\in\pc(M)}$ is orthogonal positive in the sense of Arthur (cf. \cite[pp.~19--20]{Art91}). It follows that one has
$$\mathrm{w}_M^\xi(g)=\sum_{\mu\in X_*(M)}\ \sum_{P\in\pc(M)}(-1)^{|\Delta_P^\Lambda|}\varphi_P^\Lambda(\mu+\xi_M+H_P(g)).$$

Following Arthur (cf. \cite[\S 6]{Art91}), we introduce for $P\in\pc(M)$ the Fourier series
$$S_P(\Lambda)=\sum_{\mu\in X_*(M)}\varphi_P^\Lambda(\mu+\xi_M+H_P(g))\exp(\Lambda(\mu)).$$
This series is absolutely convergent on the complement in $\ago_M^\ast$ of the hyperplanes with equation $\Lambda(\al^\vee)=0$ for $\al\in\Delta_P$. Its sum is a continuous function on this open set. Moreover, the alternating sum
$$\sum_{P\in\pc(M)}(-1)^{|\Delta_P^\Lambda|}S_P(\Lambda)$$
has as limit at $\Lambda=0$ (where the limit is taken over the generic $\Lambda$) the weight $\mathrm{w}_M^\xi(g)$.

To conclude, it remains for us to compute the sum of the series $S_P(\Lambda)$. Since $H_P(g)\in X_*(M)$, one has by an obvious change of variable
$$S_P(\Lambda)=\sum_{\mu\in X_*(M)}\varphi_P^\Lambda(\mu+\xi_M)\exp(\Lambda(\mu-H_P(g))).$$
In the above sum, it is more convenient to sum first over $X_*(M_{\scnx})$ and then over a system of representatives of $X_*(M)/X_*(M_{\scnx})$. Let $\mu_0\in X_*(M)/X_*(M_{\scnx})$. The contribution of $\mu_0$ to $S_P(\Lambda)$ is
$$\sum_{\mu\in X_*(M_{\scnx})}\varphi_P^\Lambda(\mu+\mu_0+\xi_M)\exp(\Lambda(\mu+\mu_0-H_P(g))).$$
By a change of variable and the use of the decomposition
$$\mu_0+\xi_M=[\mu_0+\xi_M]_P+\{\mu_0+\xi_M\}_P,$$
one sees that the above sum is equal to
$$\exp(\Lambda(\mu_0-[\mu_0+\xi]_P-H_P(g)))\sum_{\mu\in X_*(M_{\scnx})}\varphi_P^\Lambda(\mu+\{\mu_0+\xi_M\}_P)\exp(\Lambda(\mu)).$$
Since one has $\{\mu_0+\xi_M\}_P=\sum_{\al\in\Delta_P}r_\al\al^\vee$ with $0<r_\al\leq 1$ and since $\Delta_P^\vee$ is a basis of $X_*(M_{\scnx})$, the above sum is equal to
$$\sum_{(m_\al)_{\al\in\Delta_P}}\exp\Bigl(\Lambda\Bigl(\sum_{\al\in\Delta_P}m_\al\al^\vee\Bigr)\Bigr)$$
where one sums over the families of integers $(m_\al)_{\al\in\Delta_P}$ which satisfy $m_\al\geq 0$ for $\al\in\Delta_P^\Lambda$ and $m_\al\leq-1$ for $\al\in\Delta_P-\Delta_P^\Lambda$. These geometric series have as sum
$$\sum_{(m_\al)_{\al\in\Delta_P}}\exp\Bigl(\Lambda\Bigl(\sum_{\al\in\Delta_P}m_\al\al^\vee\Bigr)\Bigr)=(-1)^{|\Delta_P^\Lambda|}\mathrm{c}_P(\Lambda)^{-1}.$$
The announced result follows.
\end{preuve}
\end{paragr}

\begin{paragr}[Arthur's Weights and Weighted Orbital Integrals.] --- \label{S:poids-arthur}
As in \S \ref{S:HN}, we equip $\ago_T$ with a scalar product invariant under $W$. For every $L\in\lc$, we equip $\ago_L$ with the Haar measure which gives covolume $1$ to the lattices generated by orthonormal bases in $\ago_L$. By definition, for every $g\in G(\AAA)$, Arthur's weight
$$\mathrm{v}_M(g)$$
is the volume in $\ago_M$ of the convex hull of the points $-H_P(g)$. We shall recall the analytic expression due to Arthur of this weight. Beforehand, for all $P\in\pc(M)$ and $\Lambda\in\ago_M^\ast$, we introduce the polynomial
$$\mathrm{d}_P(\Lambda)=\vol(\ago_M/X_*(M_{\scnx}))^{-1}\prod_{\al\in\Delta_P}\Lambda(\al^\vee).$$
For $M=G$, one has, by convention, $\vol(\ago_G/X_*(G_{\scnx}))=1$ and the above polynomial is the constant polynomial equal to $1$.

\begin{proposition}\label{prop:poids-arthur}
(Arthur \cite[pp.~219--220]{Art76}). For every $g\in G(\AAA)$, one has the following equality
$$\mathrm{v}_M(g)=\lim_{\Lambda\to 0}\sum_{P\in\pc(M)}\mathrm{d}_P(\Lambda)^{-1}\exp(-\Lambda(H_P(g)))$$
where the limit is taken over the generic $\Lambda\in\ago_M^\ast$.
\end{proposition}

\begin{preuve}
It is analogous to the proof of Proposition \ref{prop:formule-poids}, whose notations we take up again. By the lemma of Langlands already mentioned, the weight $\mathrm{v}_M(g)$ is equal to the limit at $\Lambda=0$ of the Fourier transform
$$\int_{\ago_M}\sum_{P\in\pc(M)}(-1)^{|\Delta_P^\Lambda|}\varphi_P^\Lambda(\mu+H_P(g))\exp(\Lambda(\mu))\,d\mu.$$
Now this can be computed and equals $\sum_{P\in\pc(M)}\mathrm{d}_P(\Lambda)^{-1}\exp(-\Lambda(H_P(g)))$.
\end{preuve}

Let us take up again the notations of \S \ref{S:JMxi}. When one lets $T_Y(\AAA_0)$ act by left translation on $G(\AAA_0)$, the convex hull of the points $(-H_P(g))$ for $P\in\pc(M)$ undergoes a translation by a vector of $H_M(T_Y(\AAA_0))$. The volume $\mathrm{v}_M(g)$ is therefore left invariant under $T_Y(\AAA_0)$. Arthur's weighted orbital integral (cf. \cite[\S 8]{Art78}) is defined by the formula
\begin{equation}
\label{eq:JM}
J_M(Y)=J_M(Y,\mathbf{1}_{D_0})=\int_{T_Y(\AAA_0)\backslash G(\AAA_0)}\mathbf{1}_{D_0}(\Ad(g^{-1})Y)\mathrm{v}_M(g)\,dg.
\end{equation}

These are in fact analogues for Lie algebras of Arthur's weighted orbital integrals. These integrals appear naturally in an analogue for Lie algebras of the Arthur--Selberg trace formula (cf. \cite{Cha02}).
\end{paragr}

\begin{paragr}[The $(G,M)$-Families.] --- \label{S:GM-familles}
In \cite{Art81}, Arthur introduced the notion of $(G,M)$-family. A $(G,M)$-family is a family of functions $(\mathrm{b}_P)_{P\in\pc(M)}$ which are smooth on $\ago_M^\ast$ and which satisfy, for a pair $(P,P')$ of adjacent parabolic subgroups, the ``gluing'' condition
$$\mathrm{b}_P(\Lambda)=\mathrm{b}_{P'}(\Lambda)$$
on the hyperplane $\Lambda(\al^\vee)=0$ where $\al^\vee$ is the unique element of $\Delta_P^\vee\cap(-\Delta_{P'}^\vee)$. One can then define, \emph{a priori} only for generic $\Lambda$, the function
\begin{equation}
\label{eq:bM}
\mathrm{b}_M(\Lambda)=\sum_{P\in\pc(M)}\mathrm{d}_P(\Lambda)^{-1}\mathrm{b}_P(\Lambda)
\end{equation}
and Arthur shows that this function extends to a smooth function on $\ago_M^\ast$ (cf. \cite[Lemma 6.2]{Art81}). One sets
$$\mathrm{b}_M=\mathrm{b}_M(0).$$
The examples of $(G,M)$-families that we shall consider are the following two. The first example, which depends on $g\in G(\AAA)$, is the family $(\mathrm{v}_P(g))_{P\in\pc(M)}$ where, for every $P\in\pc(M)$, one introduces the function of the variable $\Lambda$
\begin{equation}
\label{eq:vP}
\mathrm{v}_P(g,\Lambda)=\exp(-\Lambda(H_P(g))).
\end{equation}
The gluing condition follows from Lemma \ref{lem:pos}. The second example is the family $(\mathrm{w}_P(\mu))_{P\in\pc(M)}$ which depends on $\mu\in\ago_M$ and which is defined by
\begin{equation}
\label{eq:wP}
\mathrm{w}_P(\mu,\Lambda)=\frac{\mathrm{d}_P(\Lambda)}{\mathrm{c}_P(\Lambda)}\exp(-\Lambda([\mu]_P)).
\end{equation}
The reader will verify immediately that the lemma below implies that the preceding family is a $(G,M)$-family.

\begin{lemme}\label{lem:recollement}
Let $\mu\in\ago_M$. Let $P$ and $P'$ be two adjacent parabolic subgroups in $\pc(M)$ and let $\al^\vee$ be the unique element of $\Delta_P^\vee\cap(-\Delta_{P'}^\vee)$. For every $\Lambda\in\ago_M^\ast$ which satisfies $\Lambda(\al^\vee)=0$, one has
$$\Lambda([\mu]_P)=\Lambda([\mu]_{P'}).$$
\end{lemme}

\begin{preuve}
The hyperplane of $\ago_M^\ast$ with equation $\Lambda(\al^\vee)=0$ is none other than the subspace $\ago_Q^\ast$ where $Q$ is the smallest parabolic subgroup of $G$ which contains both $P$ and $P'$. The projection $p_Q$ from $\ago_M$ onto $\ago_Q$ according to the decomposition $\ago_M=\ago_Q\oplus\ago_M^Q$ induces a bijection from $\Delta_P^\vee-\{\al^\vee\}$, respectively $\Delta_{P'}^\vee-\{-\al^\vee\}$, onto $\Delta_Q^\vee$. Let $\mu\in\ago_M$. Let us write this element in the bases $\Delta_P^\vee$ and $\Delta_{P'}^\vee$
$$\mu=\sum_{\be\in\Delta_P}x_\be\be^\vee=\sum_{\gamma\in\Delta_{P'}}y_\gamma\gamma^\vee.$$
For all $\be\in\Delta_P$ and $\gamma\in\Delta_{P'}$ distinct from $\pm\al$ such that $p_Q(\be^\vee)=p_Q(\gamma^\vee)$ one therefore has $x_\be=y_\gamma$. Using the notation $[\cdot]$ for the integer part, one therefore has
$$p_Q([\mu]_P)=\sum_{\be\in\Delta_P-\{\al\}}[x_\be]p_Q(\be^\vee)=\sum_{\gamma\in\Delta_{P'}-\{\al\}}[y_\gamma]p_Q(\gamma^\vee)=p_Q([\mu]_{P'})$$
whence
$$\Lambda([\mu]_P)=\Lambda([\mu]_{P'})$$
for every $\Lambda\in\ago_Q^\ast$.
\end{preuve}
\end{paragr}

\begin{paragr}[Product of Two $(G,M)$-Families.] --- \label{S:produit-GM}
The product of two $(G,M)$-families is obviously a $(G,M)$-family. We denote by $(\mathrm{v}\cdot\mathrm{w})(g,\mu)$ the product of the families $\mathrm{v}(g)$ and $\mathrm{w}(\mu)$ defined in (\ref{eq:vP}) and (\ref{eq:wP}) above. In the following lemma, we reformulate Propositions \ref{prop:formule-poids} and \ref{prop:poids-arthur}.

\begin{lemme}\label{lem:reformulation}
For every $g\in g(\AAA)$, the weights $\mathrm{w}_M^\xi(g)$ and $\mathrm{v}(g)$ are respectively the values at $\Lambda=0$ of the smooth functions
$$\sum_{\mu\in X_*(M)/X_*(M_{\scnx})}(\mathrm{v}\cdot\mathrm{w})_M(g,\mu+\xi_M,\Lambda)$$
and
$$\mathrm{v}_M(g,\Lambda).$$
\end{lemme}

The key to our problem of comparison of weighted orbital integrals is a formula for the product $(\mathrm{v}\cdot\mathrm{w})_M$ due to Arthur. Before being able to state it in Lemma \ref{lem:arthur-bQ}, we shall need a few additional notations.

Let $\mathrm{b}$ be a $(G,M)$-family and $L\in\lc(M)$. For every $Q\in\pc(L)$, the function on $\ago_L^\ast$ obtained by restriction of $\mathrm{b}_P$ does not depend on the choice of the parabolic subgroup $P\in\pc^Q(M)$. The family $(\mathrm{b}_Q)_{Q\in\pc(L)}$ is a $(G,L)$-family. By formula (\ref{eq:bM}) applied to the Levi subgroup $L$, one obtains a smooth function denoted $\mathrm{b}_L(\Lambda)$ on $\ago_L^\ast$ and one sets
$$\mathrm{b}_L=\mathrm{b}_L(0).$$

\begin{lemme}\label{lem:wL}
Let $L\in\lc(M)$. For every $\mu\in\ago_M$ let $(\mathrm{w}_Q(\mu,\Lambda))_{Q\in\pc(L)}$ be the $(G,L)$-family deduced as above from the $(G,M)$-family $(\mathrm{w}_P(\mu,\Lambda))_{P\in\pc(M)}$ defined in (\ref{eq:wP}). Then one has
$$\sum_{\mu\in X_*(M)/X_*(M_{\scnx})}\mathrm{w}_L(\mu+\xi_M)=\frac{\vol(\ago_L/X_*(L))}{\vol(\ago_M/X_*(M))}\cdot\mathrm{w}_L^\xi(1)$$
where:
\begin{itemize}
\item by convention, one sets $\vol(\ago_G/X_*(G))=1$;
\item $1$ is the neutral element of $G(\AAA_0)$;
\item the right-hand side is the one defined in \S \ref{S:poids-wM} when one replaces $M$ by $L$.
\end{itemize}
\end{lemme}

\begin{preuve}
The left-hand side equals by definition
$$\lim_{\Lambda\to 0}\sum_{\mu_0\in X_*(M)/X_*(M_{\scnx})}\ \sum_{Q\in\pc(L)}\mathrm{d}_Q(\Lambda)^{-1}\cdot\frac{\mathrm{d}_P(\Lambda)}{\mathrm{c}_P(\Lambda)}\cdot\exp(-\Lambda([\mu_0+\xi_M]_P))$$
where:
\begin{itemize}
\item the limit is taken over the generic $\Lambda\in\ago_L^\ast$;
\item for each $Q\in\pc(L)$, one chooses $P\in\pc^Q(M)$;
\item the fraction $\mathrm{d}_P(\Lambda)/\mathrm{c}_P(\Lambda)$ is well defined at the generic points of $\ago_M^\ast$ and extends to a smooth function on $\ago_M^\ast$.
\end{itemize}
For such $\Lambda$, $Q$ and $P$, one verifies the formulas
$$\frac{\mathrm{d}_P(\Lambda)}{\mathrm{c}_P(\Lambda)}=\frac{\vol(\ago_L/X_*(L_{\scnx}))}{\vol(\ago_M/X_*(M_{\scnx}))}\cdot\frac{\mathrm{d}_Q(\Lambda)}{\mathrm{c}_Q(\Lambda)}$$
and
$$[\mu_0+\xi_M]_P=[\mu_0'+\xi_L]_Q$$
where $\mu_0'$ is the projection of $\mu_0$ onto $\ago_L$. Lemma \ref{lem:diagramme} implies that the projection from $\ago_M$ onto $\ago_L$ induces a surjection from $X_*(M)/X_*(M_{\scnx})$ onto $X_*(L)/X_*(L_{\scnx})$. The order of its kernel combines with the ratio of the covolumes of the lattices $X_*(L_{\scnx})$ and $X_*(M_{\scnx})$ to give the factor
$$\frac{\vol(\ago_L/X_*(L))}{\vol(\ago_M/X_*(M))}.$$
One concludes the proof by applying Proposition \ref{prop:formule-poids} to the Levi subgroup $L$ and to $g=1$.
\end{preuve}

Let $Q$ be a parabolic subgroup with Levi $L$. Arthur also defines an $(L,M)$-family $\mathrm{b}^Q$ in the following way: for every $P\in\pc^L(M)$, $\mathrm{b}_P^Q(\Lambda)$ is the smooth function on $\ago_M^\ast$ defined by
$$\mathrm{b}_P^Q(\Lambda)=\mathrm{b}_{PN_Q}(\Lambda),$$
where the group $PN_Q$ is the element of $\pc(M)$ generated by $P$ and $N_Q$. One can then form a smooth function
$$\mathrm{b}_M^Q(\Lambda)=\sum_{P\in\pc^L(M)}\mathrm{d}_P^Q(\Lambda)^{-1}\mathrm{b}_P^Q(\Lambda)$$
using the coefficient
$$\mathrm{d}_P^Q(\Lambda)=\vol(\ago_M^Q/X_*(M'))^{-1}\prod_{\al\in\Delta_P^Q}\Lambda(\al^\vee)$$
where $\Delta_P^Q\subset\Delta_P$ is the subset of the roots in $L$ and $M'$ is the inverse image of $M$ in the simply connected covering of the derived group of $L$. The $\ZZ$-module $X_*(M')$ is then the lattice of $\ago_M^Q$ generated by $\Delta_P^Q$.

Let us note the following lemma, which is due to Arthur and whose proof we leave to the reader.

\begin{lemme}\label{lem:vQ}
Let $L\in\lc(M)$ and $Q\in\pc(L)$. Let $g\in G(\AAA_0)$.

Let $\mathrm{v}^Q(g)$ be the $(L,M)$-family deduced from $\mathrm{v}(g)$.

For every $l\in L(\AAA_0)$, let $\mathrm{v}^L(l)$ be the $(L,M)$-family defined by $\mathrm{v}_P^L(l,\Lambda)=\exp(-\Lambda(H_P(l)))$.

One has the equality of $(L,M)$-families
$$\mathrm{v}^Q(g)=\mathrm{v}^L(l_Q(g))$$
where $l_Q(g)\in L(\AAA_0)$ is given by the Iwasawa decomposition
$$g\in l_Q(g)N_Q(\AAA_0)K.$$
In particular, one has the equality
\begin{equation}
\label{eq:vQM}
\mathrm{v}_M^Q(g)=\mathrm{v}_M^L(l_Q(g)).
\end{equation}
\end{lemme}

Arthur also deduces from a $(G,M)$-family $\mathrm{b}$ a smooth function on $\ago_Q^\ast$ denoted $\mathrm{b}_Q'(\Lambda)$ (cf. \cite[\S 6 (6.3) and Lemma 6.1]{Art81}). We shall not recall its definition; only the following two properties will be useful to us.

\begin{lemme}\label{lem:arthur-bQ}
(Cf. \cite[Lemma 6.3 and Corollary 6.4]{Art81}). Let $\mathrm{b}$ be a $(G,M)$-family.
\begin{enumerate}
\item For every $L\in\lc(M)$ and $\Lambda\in\ago_L^\ast$, one has
$$\mathrm{b}_L(\Lambda)=\sum_{Q\in\pc(L)}\mathrm{b}_Q'(\Lambda).$$
\item Let $\mathrm{e}$ be a $(G,M)$-family and $\mathrm{e}\cdot\mathrm{b}$ the product of $\mathrm{e}$ and $\mathrm{b}$. For every $\Lambda\in\ago_M^\ast$, one has the equality
$$(\mathrm{e}\cdot\mathrm{b})_M(\Lambda)=\sum_{Q\in\fc(M)}\mathrm{e}_M^Q(\Lambda)\mathrm{b}_Q'(\Lambda).$$
\end{enumerate}
\end{lemme}
\end{paragr}

\begin{paragr}[Comparison of Weighted Orbital Integrals.] --- \label{S:comparaison}
Let $L\in\lc(M)$ and $Q\in\pc(L)$. We take up again the notations of \S \ref{S:JMxi}. The ``weights'' $\mathrm{v}_M^Q$ and $\mathrm{v}_M^L$ deduced from the eponymous $(L,M)$-families make it possible to define the weighted orbital integrals below
$$J_M^Q(Y,\mathbf{1}_{D_0})=\int_{T_Y(\AAA_0)\backslash G(\AAA_0)}\mathbf{1}_{D_0}(\Ad(g^{-1})Y)\mathrm{v}_M^Q(g)\,dg$$
and
$$J_M^L(Y,\mathbf{1}_{D_0})=\int_{T_Y(\AAA_0)\backslash L(\AAA_0)}\mathbf{1}_{D_0}(\Ad(l^{-1})Y)\,\mathrm{v}_M^L(l)\,dl.$$

For $Q=G$ or $L=G$, one recovers the integral $J_M$ defined in (\ref{eq:JM}). The two integrals are related by the descent formula stated below.

\begin{lemme}\label{lem:descente}
With the notations below, one has
$$J_M^Q(Y,\mathbf{1}_{D_0})=q^{\dim(\ngo_Q)\deg(D_0)}J_M^L(Y,\mathbf{1}_{D_0}).$$
\end{lemme}

\begin{preuve}
It rests on formula (\ref{eq:vQM}) of Lemma \ref{lem:vQ}, on the Iwasawa decomposition $G(\AAA_0)=L(\AAA_0)N_Q(\AAA_0)K$, which is compatible with our choices of measures, as well as, for $l\in L(\AAA_0)$, on the change of variables $n\mapsto\Ad(ln)^{-1}Y-Y$ which induces an isomorphism between the measured spaces $N_Q(\AAA_0)$ and $\ngo_Q(\AAA_0)$, the latter being equipped with the Haar measure which gives volume $1$ to the product $\prod_{v\in V_0}\ngo_Q[[z_v]](\Fq)$. The factor in $q$ comes from the equality
$$\int_{\ngo_Q(\AAA_0)}\mathbf{1}_{D_0}(Z)\,dZ=q^{\dim(\ngo_Q)\deg(D_0)}.$$
\end{preuve}

One can then state the following theorem.

\begin{theoreme}\label{thm:comparaison}
With the notations above, one has
$$J_M^\xi(Y)=\sum_{L\in\lc(M)}q^{1/2(\dim(G)-\dim(L))\deg(D_0)}\cdot\frac{\vol(\ago_L/X_*(L))}{\vol(\ago_M/X_*(M))}\cdot J_M^L(Y)\cdot\mathrm{w}_L^\xi(1)$$
where $1\in G(\AAA_0)$ is the neutral element.
\end{theoreme}

\begin{preuve}
Let $g\in G(\AAA_0)$ and $\mu\in\ago_M$. By assertion (2) of Lemma \ref{lem:arthur-bQ}, one has
$$(\mathrm{v}\cdot\mathrm{w})_M(g,\mu,\Lambda)=\sum_{Q\in\fc(M)}\mathrm{v}_M^Q(g,\Lambda)\mathrm{w}_Q'(\mu,\Lambda).$$
Using this formula and Lemma \ref{lem:reformulation}, one obtains
$$\mathrm{w}_M^\xi(g)=\sum_{Q\in\fc(M)}\mathrm{v}_M^Q(g)\cdot\Bigl(\sum_{\mu\in X_*(M)/X_*(M_{\scnx})}\mathrm{w}_Q'(\mu+\xi_M)\Bigr)$$
whence one deduces the relation
\begin{equation}
\label{eq:JMxi1}
J_M^\xi(Y)=\sum_{Q\in\lc(M)}J_M^Q(Y)\cdot\Bigl(\sum_{\mu\in X_*(M)/X_*(M_{\scnx})}\mathrm{w}_Q'(\mu+\xi_M)\Bigr).
\end{equation}
Using the descent formula of Lemma \ref{lem:descente}, line (\ref{eq:JMxi1}) above can be rewritten
\begin{equation}
\label{eq:JMxi2}
J_M^\xi(Y)=\sum_{L\in\lc(M)}q^{1/2(\dim(G)-\dim(L))\deg(D_0)}J_M^L(Y)\cdot\Bigl(\sum_{\mu\in X_*(M)/X_*(M_{\scnx})}\ \sum_{Q\in\pc(L)}\mathrm{w}_Q'(\mu+\xi_M)\Bigr).
\end{equation}
By assertion (1) of Lemma \ref{lem:arthur-bQ}, one has
$$\sum_{Q\in\pc(L)}\mathrm{w}_Q'(\mu+\xi_M)=\mathrm{w}_L(\mu+\xi_M)$$
and by Lemma \ref{lem:wL}, one has
$$\sum_{\mu\in X_*(M)/X_*(M_{\scnx})}\mathrm{w}_L(\mu+\xi_M)=\frac{\vol(\ago_L/X_*(L))}{\vol(\ago_M/X_*(M))}\cdot\mathrm{w}_L^\xi(1).$$
This concludes the proof.
\end{preuve}

\begin{corollaire}\label{cor:comparaison}
With the notations above, one has for every $\xi\in\ago_T$ in general position the equality
$$J_M^\xi(Y)=\vol(\ago_M/X_*(M))^{-1}\cdot J_M^G(Y).$$
\end{corollaire}

\begin{preuve}
Let $L\in\lc(M)$. By definition (cf. \S \ref{S:poids-wM}), the weight $\mathrm{w}_L^\xi(1)$ equals $1$ if $0$ belongs to $\xi_L+X_*(L)$ and $0$ otherwise. Now since $\xi$ is in general position, the first case is only possible if $L=G$ (cf. Definition \ref{def:posgenerale}). Whence the corollary.
\end{preuve}
\end{paragr}

\begin{paragr}[Proof of Theorem \ref{thm:comptage}.] --- \label{S:dem-thm-comptage}
Theorem \ref{thm:comptage} is simply a reformulation of Proposition \ref{prop:premier-comptage}, for $\xi$ in general position, by means of Corollary \ref{cor:comparaison}.
\end{paragr}

% ---------- sec12 ----------
\section{The Example of a Rank One Semisimple Group}\label{sec:rang1}

\begin{paragr}[Introduction.] --- \label{S:rang1-intro}
The aim of this section is to describe, in a simple case, a non-elliptic Hitchin fiber. The situation will be the following. The base field $k$ is an algebraic closure of a finite field of characteristic $\neq 2$. The group $G$ is the group $\operatorname{PGL}(2)$ or $\SL(2)$ over $k$. Let $T\subset G$ be the standard maximal torus, with Lie algebra $\tgo\subset \ggo$. Let $\ggo$ be the Lie algebra common to $\operatorname{PGL}(2)$ and $\SL(2)$: it is the Lie algebra of square matrices of size $2$ and trace zero. The other notations are those used in the article (in particular \S\S\,\ref{sec:morcar} and \ref{sec:schema-car}).

The algebra $k[\ggo]^G$ is the subalgebra of $k[\ggo]$ generated by the determinant. One therefore identifies the space $\car$ with the affine line $\AAA^1_k$, and the characteristic morphism $\chi:\ggo\to \car$ is identified with the morphism $\ggo\to \AAA^1$ given by the determinant. This morphism is $\Gmk$-equivariant when $\Gmk$ acts naturally on $\ggo$ and through $z\mapsto z^2$ on $\AAA^1_k$. The open subset $\car^{\reg}$ is the open subset formed by the affine line with the origin removed.

The curve $C$ is the standard projective line $\mathbb{P}^1_k$, which we equip with homogeneous coordinates $(X:Y)$. Let $0$ and $\infty$ be the points of respective coordinates $(0:1)$ and $(1:0)$. We thus have
$$\mathbb{P}^1_k-\{\infty\}=\Spec(k[x])$$
and
$$\mathbb{P}^1_k-\{0\}=\Spec(k[y])$$
with $x=X/Y$ and $y=Y/X$. Let $D$ be the divisor
$$D=2[0].$$

One obtains the ruled surface $\car_D\to C$ by gluing $\Spec(k[x,u])\to \Spec(k[x])$ and $\Spec(k[y,v])\to \Spec(k[y])$ along the open subsets $x\neq 0$ and $y\neq 0$ by the change of coordinates
$$y=1/x \quad\text{and}\quad v=u/x^4.$$
In an analogous way, one obtains the ruled surface $\tgo_D\to C$ by gluing the same open subsets as previously but according to the change of coordinates
$$y=1/x \quad\text{and}\quad v=u/x^2.$$
\end{paragr}

\begin{paragr}[The Base of the Hitchin Fibration.] --- \label{S:rang1-base}
By the above description of the ruled surface $\car_D$, one has the identification
$$H^0(C,\car_D)=\{a\in k[y]\mid \deg(a)\leq 4\}$$
so that the base of the Hitchin fibration is the $5$-dimensional scheme
$$\Ac=\{(a,t)\in k[y]\times k^\times \mid \deg(a)\leq 4 \text{ and } t^2=a(0)\}.$$
Likewise, one has
$$\Ac_T=\{(a,t)\in k[y]\times k^\times \mid \deg(a)\leq 2 \text{ and } t=a(0)\}.$$
The closed immersion $\Ac_T\hookrightarrow \Ac$ is given by
$$(a,t)\mapsto (a^2,t).$$

The elliptic open subset
$$\Ac^{\el}=\Ac-\Ac_T$$
is formed of the pairs $(a,t)$ such that $a$ is not a square in $k[y]$.
\end{paragr}

\begin{paragr}[The Hitchin Fibration.] --- \label{S:rang1-fibration}
The Hitchin space $\mc_{\SL(2)}$ (cf. \S\,\ref{sec:fib-Hitchin}) is isomorphic to the stack of triples $(\vc,\theta,t)$ where $\vc$ is a rank $2$ vector bundle equipped with a trivialization of its determinant,
$$\theta:\vc\to \vc(D)$$
is a homomorphism of $\oc_C$-modules of trace zero such that $\theta_\infty:\vc_\infty\to \vc_\infty$ is regular, semisimple and of determinant equal to $t^2$. The Hitchin space $\mc_{\operatorname{PGL}(2)}$ is isomorphic to the stack of triples $(\vc,\theta,t)$, where $\vc$ is a rank $2$ vector bundle and $(\theta,t)$ is as above, taken ``up to tensorization'' by a line bundle on $C$. In both cases, the Hitchin morphism sends a triple $(\vc,\theta,t)$ to $(\det(\theta),t)$.
\end{paragr}

\begin{paragr}[Connected Components.] --- \label{S:rang1-composantes}
It is well known that the total space $\mc_{\SL(2)}$ of the Hitchin fibration for $\SL(2)$ is connected. The same does not hold for $\mc_{\operatorname{PGL}(2)}$, as we shall see. The exact sequence
$$1\to \mu_2\to \SL(2)\to \operatorname{PGL}(2)\to 1$$
induces an exact sequence of groups
$$0\to H^0(C,\mu_2)\to H^0(C,\SL(2))\to H^0(C,\operatorname{PGL}(2))\to 1$$
which is none other than the initial exact sequence, and an exact sequence of pointed sets
\begin{equation}
  \label{eq:rang1-suite}
  0\to H^1(C,\SL(2))\to H^1(C,\operatorname{PGL}(2))\to H^2(C,\mu_2)\to 1.
\end{equation}
Note that one has $H^1(C,\mu_2)=(1)$ since $C=\mathbb{P}^1_k$ and
$$H^2(C,\mu_2)\simeq \ZZ/2\ZZ.$$
Let us verify the surjectivity on the right in the sequence (\ref{eq:rang1-suite}). We start from the sequence
$$1\to \mu_2\to T_{\SL(2)}\to T_{\operatorname{PGL}(2)}\to 1$$
where $T_{\SL(2)}$ is a maximal subtorus of $\SL(2)$ and $T_{\operatorname{PGL}(2)}$ is its image in $\operatorname{PGL}(2)$. One deduces the sequence
$$0\to H^1(C,T_{\operatorname{PGL}(2)})\to H^2(C,\mu_2)\to H^2(C,T_{\SL(2)}).$$
But $H^2(C,T_{\SL(2)})=(1)$ since $T_{\SL(2)}\simeq \Gmk$ and $k$ is algebraically closed, whence the surjectivity of the arrow $H^1(C,T_{\operatorname{PGL}(2)})\to H^2(C,\mu_2)$ and \emph{a fortiori} of $H^1(C,\operatorname{PGL}(2))\to H^2(C,\mu_2)$.

The sequence (\ref{eq:rang1-suite}) implies that the space $\mc_{\operatorname{PGL}(2)}$ possesses two connected components $\mc^i_{\operatorname{PGL}(2)}$ indexed by $i\in\ZZ/2\ZZ$. The natural morphism
$$\mc_{\SL(2)}\to \mc^0_{\operatorname{PGL}(2)}$$
is a $\mu_2$-gerbe. Indeed, by the sequence (\ref{eq:rang1-suite}), every $\operatorname{PGL}(2)$-torsor whose class has trivial image in $H^2(C,\mu_2)$ lifts to an $\SL(2)$-torsor, and this lifting is unique up to an element of
$$\mu_2=\Ker(H^0(C,\SL(2))\to H^0(C,\operatorname{PGL}(2))).$$
\end{paragr}

\begin{paragr}[A Hyperbolic Fiber.] --- \label{S:rang1-hyperbolique}
From now on, \emph{we restrict ourselves} to the case $G=\operatorname{PGL}(2)$ and we simply write $\mc=\mc_G$. We are going to study the Hitchin fiber $\mc_a$ above the point $a=(y^2-1,-1)\in \Ac_T$ (which is also identified with the point $((y^2-1)^2,-1)\in \Ac_G$). Let
$$U=\mathbb{P}^1_k-\{-1,1\}=\Spec\left(k\left[\left(\frac{y+1}{y-1}\right)^{\pm 1}\right]\right).$$
Let $X\in\tgo(k[y])$ be defined by
$$X=\begin{pmatrix} y^2-1 & 0\\ 0 & 1-y^2\end{pmatrix}.$$
This element is $G$-regular on $U$ with centralizer the constant torus scheme $T\times_k U$. Indeed, its trace is always zero and its determinant does not vanish on $U$. Let $(\ec,\theta,t)\in\mc_a$. It follows from the considerations of \S\,\ref{S:reductouvert} that $(\ec,\theta,t)$ possesses on $U$ a canonical reduction to $T$, say $(\ec_T,\theta,t)$. We identify $T$ once and for all with $\Gmk$ via the morphism
$$z\mapsto \begin{pmatrix} z & 0\\ 0 & 1\end{pmatrix}_{\!\ast}$$
where $\ast$ indicates that one takes the image in $\operatorname{PGL}(2)$. Since every $\Gmk$-torsor is trivial on $U$, one may choose a trivialization of $\ec_T$. Let $\xc_a$ be the scheme formed of the triples
$$(\ec,\theta,\iota)$$
such that $(\ec,\theta,t)$ is a Hitchin triple above $a$ and $\iota$ is a trivialization of $\ec$ on $U$ which sends $\theta$ to $X$. The group
$$H^0(U,T)\simeq \Gmk(U)=k^\times\left(\frac{y+1}{y-1}\right)^{\ZZ}\simeq k^\times\times \ZZ$$
acts on the left on $\xc$ by the obvious action on the factor $\iota$. The fiber $\mc_a$ is therefore isomorphic to the quotient stack $[\Gmk\times\ZZ\back \xc_a]$.
\end{paragr}

\begin{paragr}[Affine Springer Fibers.] --- \label{S:rang1-springer}
Let $z$ be an indeterminate over $k$. Let $\oc=k[[z]]$ and $F=k((z))$ its field of fractions. Let $\pgo$ be the ideal of $\oc$ generated by $z$. Let
$$\val:F^\times\to \ZZ$$
be the discrete valuation normalized by $\val(z)=1$. Let
$$Y=\begin{pmatrix} z & 0\\ 0 & -z\end{pmatrix}\in\tgo(\oc)$$
and $M$ the affine Springer fiber (cf. \cite{KL88} and \cite{GKM04}) defined by
$$M=\{g\in G(F)/G(\oc)\mid \Ad(g^{-1})Y\in\ggo(\oc)\}.$$
This is a $k$-scheme locally of finite type which possesses two connected components corresponding to the fibers of the map
$$M\to \ZZ/2\ZZ$$
which associates to $g$ the class $\val(\det(\tilde{g}))\bmod 2\ZZ$, where $\tilde{g}\in \operatorname{GL}(2,F)$ is an arbitrary lift of $g$. Using the bijection between $\operatorname{GL}(2,F)/\operatorname{GL}(2,\oc)$ and the set of $\oc$-lattices of $F^2=F\oplus F$, one sees that $M$ is identified with the scheme of classes of lattices of $F^2$ stable under the endomorphism $Y$, for the equivalence relation on lattices defined by the homotheties with ratio in $z^{\ZZ}$. One obtains the component of index $i\in\ZZ/2\ZZ$ by considering only the lattices of index in $i+2\ZZ$ (where the index of the standard lattice $\oc\oplus\oc$ is $0$ by definition). The matrix
\begin{equation}
  \label{eq:rang1-matrice}
  \begin{pmatrix} z & 0\\ 0 & 1\end{pmatrix}_{\!\ast}
\end{equation}
acts on $M$ (according to the points of view: either by left translation on $g$, or by its natural action on the classes of lattices) and permutes the connected components of $M$. It therefore suffices to describe the component $M^0$ of the lattices of index $0\bmod 2\ZZ$. Let $n\in\ZZ$ and let $M^0_n$ be the set of classes of lattices of $F^2$, of index $0\bmod 2\ZZ$, which admit a representative $R$ (necessarily unique) satisfying the inclusions
\begin{equation}
  \label{eq:rang1-inclusions}
  \pgo^{1-n}\oplus \pgo^n\subsetneq R\subsetneq \pgo^{-n}\oplus \pgo^{n-1}.
\end{equation}
Such lattices are automatically stable under $Y$, in other words $M^0_n\subset M^0$. A lattice $R$ satisfying the above inclusions defines a line of the $2$-dimensional quotient $k$-vector space $\pgo^{1-n}\oplus\pgo^n\back \pgo^{-n}\oplus\pgo^{n-1}$. One then obtains a bijection
$$M^0_n\to \mathbb{P}(\pgo^{1-n}\oplus\pgo^n\back \pgo^{-n}\oplus\pgo^{n-1})\simeq \mathbb{P}^1_k$$
from $M^0_n$ onto the $k$-points of the projective line $\mathbb{P}^1_k$, which sends the class of the lattice $\pgo^{-n}\oplus\pgo^n$ to the point $(1:0)$ (in homogeneous coordinates) and that of the lattice $\pgo^{1-n}\oplus\pgo^{n-1}$ to the point $(0:1)$. These two classes are the only ones in $M^0_n$ that are classes of decomposed lattices. In this way one identifies $M^0$ with the chain of projective lines $M^0_n$, $n\in\ZZ$, the lines $M^0_n$ and $M^0_{n+1}$ being glued at the point corresponding to the class of the lattice $\pgo^{-n}\oplus\pgo^n$.

The left action of $T(F)$ on $G(F)$ preserves $M$. One deduces from it an action of $F^\times$ on $M$ via the isomorphism $\Gmk\simeq T$. The subgroup of the elements of $F^\times$ of valuation $2$ stabilizes $M^0$. This group is written as a direct product
$$z^{2\ZZ}\times k^\times\times (1+\pgo).$$
The component $(1+\pgo)$ acts trivially on $M^0$. The action of $z^{2k}$ induces a translation in the chain of projective lines which sends $M^0_n$ to $M^0_{n-k}$. The action of $k^\times$ stabilizes each irreducible component $M^0_n$. In homogeneous coordinates, one has $\la\cdot(x:y)=(\la x:y)$. This action therefore possesses on each $M^0_n$ two fixed points, which are the two classes of decomposed lattices, and the action of $k^\times$ on $M^0_n$ with these two points removed is simply transitive.

Let us describe in terms of lattices the functions $H_B$ and $H_{\overline{B}}$ (cf. Definition \ref{def:HP}) where $B$ is the ``upper'' Borel subgroup and $\overline{B}$ is its opposite. The space $\ago_T$ is of dimension $1$. Let $\al$ be the root of $T$ in $B$: it induces on $\Gmk\simeq T$ the identity isomorphism $\Gmk\to\Gmk$. This root (viewed as an element of $\ago_T^\ast$) induces an isomorphism $\ago_T\simeq \RR$. We set, for every $g\in G(F)/G(\oc)$,
$$h_B(g)=\al(H_B(g))\quad\text{and}\quad h_{\overline{B}}(g)=\al(H_{\overline{B}}(g)).$$
For a lattice $R$ of $F^2$, let $R_1=R\cap (F\oplus (0))$ and let $p_1(R)$ be the projection of $R$ onto $F\oplus (0)$. Using $0\oplus F$, one defines $R_2$ and $p_2(R)$. For every fractional ideal $\pgo^n$ of $F$, we set $\val(\pgo^n)=n$. This definition thus applies to $R_1$, $p_1(R)$, etc. We then set
$$h_B(R)=\val(p_2(R))-\val(R_1)\quad\text{and}\quad h_{\overline{B}}(R)=\val(R_2)-\val(p_1(R)).$$
This definition depends only on the class of the lattice $R$. We leave to the reader the task of verifying that if the class of $R$ corresponds to $g$ then one has $h_B(R)=h_B(g)$ and $h_{\overline{B}}(R)=h_{\overline{B}}(g)$ (cf. \cite[\S 9.1]{CL08}, where the function $H_B$ of that article is the \emph{opposite} of our function $H_B$).

Note that one has
$$h_{\overline{B}}(R)-h_B(R)=\val(R_1)-\val(p_1(R))+\val(R_2)-\val(p_2(R))\geq 0$$
(as it should be, cf. Lemma \ref{lem:pos}) and that this difference is zero only if $R$ is a decomposed lattice. If the class $R$ belongs to $M$, then by the inclusions (\ref{eq:rang1-inclusions}) one sees that one has
$$0\leq \val(R_1)-\val(p_1(R))\leq 1.$$
Therefore for a class of lattices $R$ in $M$ that is not decomposed, one has
$$h_{\overline{B}}(R)-h_B(R)=2.$$
Moreover, for a class of lattices $R$ in $M^0_n$, one has
\begin{equation}
  \label{eq:rang1-hB}
  h_B(R)=\begin{cases} 2n & \text{if } R=\pgo^{-n}\oplus\pgo^n;\\ 2(n-1) & \text{otherwise};\end{cases}
\end{equation}
and
\begin{equation}
  \label{eq:rang1-hBbar}
  h_{\overline{B}}(R)=\begin{cases} 2(n-1) & \text{if } R=\pgo^{1-n}\oplus\pgo^{n-1};\\ 2n & \text{otherwise.}\end{cases}
\end{equation}
\end{paragr}

\begin{paragr}[A Non-Separated Fiber.] --- \label{S:rang1-nonsep}
We are going to see that the fiber $\mc_a$ of \S\,\ref{S:rang1-hyperbolique} is neither separated nor of finite type. We have seen that $\mc_a$ is isomorphic to a certain quotient of $\xc_a$. It follows from the considerations of Section \ref{sec:description} that $\xc_a$ is identified with a product of two Springer fibers. More precisely, the formal gluing of Beauville--Laszlo identifies $\xc_a$ with the product of
$$M_+=\{g\in G(k((y-1)))/G(k[[y-1]])\mid \Ad(g^{-1})X\in\ggo(k[[y-1]])\}$$
with
$$M_-=\{g\in G(k((y+1)))/G(k[[y+1]])\mid \Ad(g^{-1})X\in\ggo(k[[y+1]])\}$$
where
$$X=\begin{pmatrix} y^2-1 & 0\\ 0 & 1-y^2\end{pmatrix}$$
and the embeddings $k[y]\hookrightarrow k[[y-1]]$ and $k[y]\hookrightarrow k[[y+1]]$ are the obvious embeddings. Since $y+1$ is a unit in $k[[y-1]]$, one sees clearly that $M_+$ is the Springer fiber $M$ defined in \S\,\ref{S:rang1-springer}. Likewise $M_-=M$.

Let us now interpret the action of $\Gmk\times\ZZ$ on $\xc_a$. This corresponds to the diagonal action of $k^\times\times ((y+1)/(y-1))^n$ on $G(k((y-1)))\times G(k((y+1)))$. In terms of $M_+\times_k M_-$, one sees that $\Gmk$ acts diagonally on $M_+\times_k M_-$ and that $\ZZ$ acts on $M_+\times_k M_-$ through $n\mapsto (-n,n)$, the actions on each factor being those described in \S\,\ref{S:rang1-springer}.

We saw in \S\,\ref{S:rang1-composantes} that $\mc$ possesses two connected components $\mc^i$ indexed by $i\in\ZZ/2\ZZ$. The fiber $\mc^i_a=\mc_a\cap\mc^i$ is therefore the quotient by $\Gmk\times\ZZ$ of the product of the components
$$M_+^{i_+}\times_k M_-^{i_-}$$
for $i_++i_-=i$, where $M_\pm^{i_\pm}$ is the component of $M_\pm$ of index $i_\pm$. In fact, the action of the matrix (\ref{eq:rang1-matrice}) on one of the factors exchanges the connected components, which are therefore isomorphic. To simplify the discussion, in what follows we consider only the component $\mc^0_a$. The action of $1\in\ZZ$, which is identified with $(y+1)/(y-1)$, permutes $M^0_+\times_k M^0_-$ and $M^1_+\times_k M^1_-$. It follows that $\mc^0_a$ is identified with the quotient of $M^0_+\times_k M^0_-$ by the action of $\Gmk\times 2\ZZ$. Figure \ref{fig:1} below represents $M^0_+\times_k M^0_-$, that is to say the product of two chains of $\mathbb{P}^1_k$. The vertices of the (infinite) grid correspond to the pairs of decomposed lattices, the vertical, respectively horizontal, segments to the pairs of lattices whose first, respectively second, component is a decomposed lattice. The points in the interior of the squares are the pairs of classes of lattices which are not decomposed.

The action of $\Gmk$ preserves each square (which is a product $M^0_{+,n_+}\times_k M^0_{-,n_-}$) and the action of $2\in 2\ZZ$ translates a square by $1$ to the right and by $1$ downwards. One thus sees that the quotient of $M^0_+\times_k M^0_-$ by $2\ZZ\times \Gmk$ is not of finite type.

\begin{figure}[ht]
\centering
\setlength{\unitlength}{1pt}
\begin{picture}(140,140)(-10,-10)
  \put(0,-10){\line(0,1){140}}
  \put(30,-10){\line(0,1){140}}
  \put(60,-10){\line(0,1){140}}
  \put(90,-10){\line(0,1){140}}
  \put(120,-10){\line(0,1){140}}
  \put(-10,0){\line(1,0){140}}
  \put(-10,30){\line(1,0){140}}
  \put(-10,60){\line(1,0){140}}
  \put(-10,90){\line(1,0){140}}
  \put(-10,120){\line(1,0){140}}
\end{picture}
\caption{The product of two chains of $\mathbb{P}^1_k$.}
\label{fig:1}
\end{figure}

In Figure \ref{fig:2} below, the vertices of the square (symbolized by the small circles) are fixed points of the action of $\Gmk$. The residual gerbe of such points is $B(\Gm)$. We have represented in thin lines the orbits of $\Gmk$ of dimension $1$. The quotient of the square with the vertices removed is therefore a non-separated projective line, symbolized by a bold line with its two origins and its two points at infinity, which are the quotients by $\Gm$ of the sides of the square. This line is isomorphic to the one obtained by gluing two standard projective lines along the open subset complementary to the origin and to the point at infinity.

\begin{figure}[ht]
\centering
\setlength{\unitlength}{1pt}
\begin{picture}(176,176)(0,0)
  \linethickness{0.4pt}
  \put(5.0,0){\line(1,0){160.0}}
  \put(5.0,170.0){\line(1,0){160.0}}
  \put(0,5.0){\line(0,1){160.0}}
  \put(170.0,5.0){\line(0,1){160.0}}
  \qbezier(1.2,4.8)(8.1,31.7)(17.0,52.7)
  \qbezier(17.0,52.7)(25.5,73.1)(36.4,89.2)
  \qbezier(36.4,89.2)(46.2,103.7)(58.3,115.5)
  \qbezier(58.3,115.5)(70.3,127.0)(85.0,136.4)
  \qbezier(85.0,136.4)(97.2,144.1)(111.7,150.6)
  \qbezier(111.7,150.6)(122.0,155.2)(133.6,159.3)
  \qbezier(133.6,159.3)(142.8,162.5)(153.0,165.5)
  \qbezier(153.0,165.5)(158.9,167.2)(165.2,168.8)
  \qbezier(2.4,4.4)(9.4,16.9)(17.0,28.6)
  \qbezier(17.0,28.6)(26.3,43.1)(36.4,56.4)
  \qbezier(36.4,56.4)(46.9,70.2)(58.3,82.9)
  \qbezier(58.3,82.9)(71.0,97.0)(85.0,109.8)
  \qbezier(85.0,109.8)(97.8,121.4)(111.7,132.1)
  \qbezier(111.7,132.1)(122.3,140.3)(133.6,147.9)
  \qbezier(133.6,147.9)(143.0,154.2)(153.0,160.3)
  \qbezier(153.0,160.3)(159.2,164.0)(165.6,167.6)
  \qbezier(3.5,3.5)(10.2,10.2)(17.0,17.0)
  \qbezier(17.0,17.0)(26.7,26.7)(36.4,36.4)
  \qbezier(36.4,36.4)(47.4,47.4)(58.3,58.3)
  \qbezier(58.3,58.3)(71.7,71.7)(85.0,85.0)
  \qbezier(85.0,85.0)(98.3,98.3)(111.7,111.7)
  \qbezier(111.7,111.7)(122.6,122.6)(133.6,133.6)
  \qbezier(133.6,133.6)(143.3,143.3)(153.0,153.0)
  \qbezier(153.0,153.0)(159.8,159.8)(166.5,166.5)
  \qbezier(4.4,2.4)(10.8,6.0)(17.0,9.7)
  \qbezier(17.0,9.7)(27.0,15.8)(36.4,22.1)
  \qbezier(36.4,22.1)(47.7,29.7)(58.3,37.9)
  \qbezier(58.3,37.9)(72.2,48.6)(85.0,60.2)
  \qbezier(85.0,60.2)(99.0,73.0)(111.7,87.1)
  \qbezier(111.7,87.1)(123.1,99.8)(133.6,113.6)
  \qbezier(133.6,113.6)(143.7,126.9)(153.0,141.4)
  \qbezier(153.0,141.4)(160.6,153.1)(167.6,165.6)
  \qbezier(4.8,1.2)(11.1,2.8)(17.0,4.5)
  \qbezier(17.0,4.5)(27.2,7.5)(36.4,10.7)
  \qbezier(36.4,10.7)(48.0,14.8)(58.3,19.4)
  \qbezier(58.3,19.4)(72.8,25.9)(85.0,33.6)
  \qbezier(85.0,33.6)(99.7,43.0)(111.7,54.5)
  \qbezier(111.7,54.5)(123.8,66.3)(133.6,80.8)
  \qbezier(133.6,80.8)(144.5,96.9)(153.0,117.3)
  \qbezier(153.0,117.3)(161.9,138.3)(168.8,165.2)
  \linethickness{1.4pt}
  \put(21.1,148.2){\line(0,1){21.8}}
  \put(0.0,148.2){\line(1,0){21.1}}
  \thicklines
  \put(20.63,147.79){\line(1,-1){127.8}}
  \put(21.08,148.24){\line(1,-1){127.8}}
  \put(21.53,148.69){\line(1,-1){127.8}}
  \linethickness{1.4pt}
  \put(148.9,21.8){\line(1,0){21.1}}
  \put(148.9,21.8){\line(0,-1){21.8}}
  \put(21.1,170.0){\circle*{5}}
  \put(0.0,148.2){\circle*{5}}
  \put(170.0,21.8){\circle*{5}}
  \put(148.9,0.0){\circle*{5}}
  \linethickness{0.4pt}
  \put(0.0,0.0){\circle{10.0}}
  \put(170.0,0.0){\circle{10.0}}
  \put(0.0,170.0){\circle{10.0}}
  \put(170.0,170.0){\circle{10.0}}
\end{picture}
\caption{In thin lines, some orbits of $\Gmk$ on $\mathbb{P}^1_k\times \mathbb{P}^1_k$ with the vertices removed, and in bold the non-separated quotient.}
\label{fig:2}
\end{figure}

In particular, $\mc^0_a$, with the points corresponding to the vertices of the grid removed, is an infinite chain of two-sided N\'eron polygons (cf. Figure \ref{fig:3} above).

\begin{figure}[ht]
\centering
\setlength{\unitlength}{1pt}
\begin{picture}(200,120)(0,-4)
  \linethickness{1.2pt}
  \qbezier(0,115)(100,-115)(200,115)
  \qbezier(0,0)(100,230)(200,0)
  \put(29.3,57.5){\circle*{6}}
  \put(170.7,57.5){\circle*{6}}
  \linethickness{0.4pt}
\end{picture}
\caption{A two-sided N\'eron polygon.}
\label{fig:3}
\end{figure}

More precisely, each N\'eron polygon of this chain is indexed by $n\in\ZZ$. Such a polygon is the union of two projective lines $P_n$ and $P'_n$ over $k$. The open subsets $(P_n\cup P'_n)-P_n$ and $(P_n\cup P'_n)-P'_n$ are isomorphic to $\Gmk$. In this chain, one glues the polygon $P_n\cup P'_n$ to the polygon $P_{n+1}\cup P'_{n+1}$ by identifying $(P_n\cup P'_n)-P_n$ with $(P_{n+1}\cup P'_{n+1})-P'_{n+1}$. In Figure \ref{fig:4} below, we have represented a piece of this chain of polygons. In order to suggest that this chain is infinite, and so that the reader may recognize in Figure \ref{fig:4} the non-separated line represented in bold in Figure \ref{fig:2}, the extremal polygons are only partially drawn (they have been amputated of a $\Gmk$). The reader who would like to recover Figure \ref{fig:4} will note that the line of Figure \ref{fig:2} is every other time ``reversed''.

\begin{figure}[ht]
\centering
\setlength{\unitlength}{1pt}
\begin{picture}(192,192)(0,0)
  \thicklines
  \put(84.0,156.0){\line(0,1){24.0}}
  \put(60.0,156.0){\line(1,0){24.0}}
  \put(60.0,132.0){\line(0,1){24.0}}
  \put(36.0,132.0){\line(1,0){24.0}}
  \put(36.0,108.0){\line(0,1){24.0}}
  \put(12.0,108.0){\line(1,0){24.0}}
  \put(12.0,96.0){\line(0,1){12.0}}
  \put(84.0,180.0){\line(1,0){12.0}}
  \put(156.0,84.0){\line(1,0){24.0}}
  \put(156.0,60.0){\line(0,1){24.0}}
  \put(132.0,60.0){\line(1,0){24.0}}
  \put(132.0,36.0){\line(0,1){24.0}}
  \put(108.0,36.0){\line(1,0){24.0}}
  \put(108.0,12.0){\line(0,1){24.0}}
  \put(180.0,84.0){\line(0,1){12.0}}
  \put(96.0,12.0){\line(1,0){12.0}}
  \put(84.0,156.0){\line(1,-1){72.0}}
  \put(60.0,132.0){\line(1,-1){72.0}}
  \put(36.0,108.0){\line(1,-1){72.0}}
  \put(84.0,180.0){\circle*{6}}
  \put(60.0,156.0){\circle*{6}}
  \put(36.0,132.0){\circle*{6}}
  \put(12.0,108.0){\circle*{6}}
  \put(180.0,84.0){\circle*{6}}
  \put(156.0,60.0){\circle*{6}}
  \put(132.0,36.0){\circle*{6}}
  \put(108.0,12.0){\circle*{6}}
  \thinlines
\end{picture}
\caption{A chain of two-sided N\'eron polygons.}
\label{fig:4}
\end{figure}
\end{paragr}

\begin{paragr}[A $\xi$-Stable Hyperbolic Fiber.] --- \label{S:rang1-xistable}
As we are going to see, the situation improves radically when one takes the $\xi$-stable fiber $\mc^\xi_a$ for a $\xi$ in general position (cf. Section \ref{sec:xi-stab}). Let $\xi\in\ago_T\simeq \RR$. Let $\mc^{\xi,0}_a$ be the component of index $0$ of $\mc^\xi_a$: it is the quotient of $(M^0_+\times_k M^0_-)^\xi$ by the action of $\Gmk\times 2\ZZ$, where $(M^0_+\times_k M^0_-)^\xi\subset M^0_+\times_k M^0_-$ is the open subset formed of the (classes of) lattices $(R_+,R_-)$ which satisfy
\begin{equation}
  \label{eq:rang1-xistab}
  -h_{\overline{B}}(R_+)-h_{\overline{B}}(R_-)<\xi<-h_B(R_+)-h_B(R_-).
\end{equation}
It follows from (\ref{eq:rang1-hB}) and (\ref{eq:rang1-hBbar}) that if $\xi\notin 2\ZZ$, the above inequalities are always strict. In what follows, in order to lighten the discussions we take $\xi=1$. Let us remark that condition (\ref{eq:rang1-xistab}) excludes the pairs of decomposed lattices. Let $R_+$ be a lattice in the class of $\pgo^{-n}\oplus\pgo^n$ and $R_-$ such that $(R_+,R_-)\in (M^0_+\times_k M^0_-)^\xi$. By (\ref{eq:rang1-hB}) and (\ref{eq:rang1-hBbar}), the inequality (\ref{eq:rang1-xistab}) becomes
$$-2n-h_{\overline{B}}(R_-)<1<-2n-h_B(R_-)$$
so necessarily $h_B(R_-)=-2n-2$ and $h_{\overline{B}}(R_-)=-2n$, so $R_-$ is a non-decomposed lattice of $M^0_{-,-n}$. Likewise, if $R_-=\pgo^{-n}\oplus\pgo^n$, every lattice $R_+$ such that $(R_+,R_-)\in (M^0_+\times_k M^0_-)^\xi$ is a non-decomposed lattice of $M^0_{+,-n}$.

Let $(R_+,R_-)$ be a pair of non-decomposed lattices with $R_+\in M^0_{+,n}$. By (\ref{eq:rang1-hB}) and (\ref{eq:rang1-hBbar}), the inequality (\ref{eq:rang1-xistab}) becomes
$$-2n-h_{\overline{B}}(R_-)<1<-2n+2-h_B(R_-).$$
There are therefore two possibilities for $h_B(R_-)$, namely $-2n$ and $-2n-2$, so $R_-\in M^0_{-,-n}\cup M^0_{-,-n-1}$.

One may therefore represent $(M^0_+\times_k M^0_-)^\xi$ by Figure \ref{fig:5}.

\begin{figure}[ht]
\centering
\setlength{\unitlength}{1pt}
\begin{picture}(164,122)(-19,-19)
  \linethickness{0.4pt}
  \put(-15.1,84.0){\line(1,0){11.1}}
  \put(42.0,88.0){\line(0,1){11.1}}
  \put(0.0,46.0){\line(0,1){34.0}}
  \put(4.0,42.0){\line(1,0){34.0}}
  \put(46.0,84.0){\line(1,0){34.0}}
  \put(84.0,46.0){\line(0,1){34.0}}
  \put(42.0,4.0){\line(0,1){34.0}}
  \put(46.0,0.0){\line(1,0){34.0}}
  \put(88.0,42.0){\line(1,0){34.0}}
  \put(126.0,4.0){\line(0,1){34.0}}
  \put(84.0,-15.1){\line(0,1){11.1}}
  \put(130.0,0.0){\line(1,0){11.1}}
  \linethickness{1.6pt}
  \put(0.0,88.0){\line(0,1){11.1}}
  \put(4.0,84.0){\line(1,0){34.0}}
  \put(42.0,46.0){\line(0,1){34.0}}
  \put(46.0,42.0){\line(1,0){34.0}}
  \put(84.0,4.0){\line(0,1){34.0}}
  \put(88.0,0.0){\line(1,0){34.0}}
  \put(126.0,-15.1){\line(0,1){11.1}}
  \linethickness{0.4pt}
  \put(0.0,84.0){\circle{8.0}}
  \put(42.0,84.0){\circle{8.0}}
  \put(84.0,84.0){\circle{8.0}}
  \put(0.0,42.0){\circle{8.0}}
  \put(42.0,42.0){\circle{8.0}}
  \put(84.0,42.0){\circle{8.0}}
  \put(126.0,42.0){\circle{8.0}}
  \put(42.0,0.0){\circle{8.0}}
  \put(84.0,0.0){\circle{8.0}}
  \put(126.0,0.0){\circle{8.0}}
\end{picture}
\caption{A $\xi$-stable open subset of the product $\mathbb{P}^1_k\times \mathbb{P}^1_k$.}
\label{fig:5}
\end{figure}

In Figure \ref{fig:5} above, $(M^0_+\times_k M^0_-)^\xi$ is the open subset which is the union of the interiors of the squares that are drawn (namely the products $M^0_{+,n}\times_k M^0_{-,-n}$ and $M^0_{+,n}\times_k M^0_{-,-n-1}$ for $n\in\ZZ$) and of the open segments in bold. The vertices, that is to say the pairs of decomposed lattices, are symbolized by small circles and are removed. The group $\Gmk\times 2\ZZ$ preserves, as it should, this open subset. The quotient of $(M^0_+\times_k M^0_-)^\xi$ by the factor $\Gmk$ is a chain of projective lines (from now on non-separated) on which the factor $2\in 2\ZZ$ acts by translation by $2$. The quotient $\mc^{0,\xi}_a$ is therefore a two-sided N\'eron polygon (cf. Figure \ref{fig:3}). The $\xi$-stability, for $\xi$ in general position, therefore keeps only a single polygon in the chain of N\'eron polygons of Figure \ref{fig:4}.
\end{paragr}

\section*{Acknowledgements}
Part of this article was written during a stay of the first author at the
\emph{Institute for Advanced Study} in Princeton in the fall of 2008. He wishes to thank that institute for its hospitality as well as the \emph{National Science Foundation} for the support (agreement No.\ DMS-0635607) which made this stay possible. The authors also thank the referees of this article for their careful reading.

% ---------- biblio ----------
\begin{thebibliography}{XXXXXX}

\bibitem[Art76]{Art76} J. Arthur, \emph{The characters of discrete series as orbital integrals}, Invent. Math. \textbf{32} (1976), 205--261. \doi{10.1007/BF01425569}

\bibitem[Art78]{Art78} J. Arthur, \emph{A trace formula for reductive groups. I. Terms associated to classes in $G(\mathbb{Q})$}, Duke Math. J. \textbf{45} (1978), 911--952. \doi{10.1215/S0012-7094-78-04542-8}

\bibitem[Art81]{Art81} J. Arthur, \emph{The trace formula in invariant form}, Ann. of Math. (2) \textbf{114} (1981), 1--74. \doi{10.2307/1971376}

\bibitem[Art91]{Art91} J. Arthur, \emph{A local trace formula}, Publ. Math. Inst. Hautes \'Etudes Sci. \textbf{73} (1991), 5--96. \doi{10.1007/BF02699256}

\bibitem[BL95]{BL95} A. Beauville and Y. Laszlo, \emph{Un lemme de descente}, C. R. Acad. Sci. Paris S\'er. I Math. \textbf{320} (1995), 335--340. \urlref{https://math.univ-cotedazur.fr/u/beauvill/pubs/descente.pdf}

\bibitem[BNR89]{BNR89} A. Beauville, M. S. Narasimhan and S. Ramanan, \emph{Spectral curves and the generalised theta divisor}, J. Reine Angew. Math. \textbf{398} (1989), 169--179. \doi{10.1515/crll.1989.398.169}

\bibitem[Beh]{Beh} K. Behrend, \emph{The Lefschetz trace formula for the moduli stack of principal bundles}, PhD thesis, University of California, Berkeley (1991). \urlref{https://personal.math.ubc.ca/~behrend/thesis.pdf}

\bibitem[Beh95]{Beh95} K. Behrend, \emph{Semi-stability of reductive group schemes over curves}, Math. Ann. \textbf{301} (1995), 281--305. \doi{10.1007/BF01446630}

\bibitem[BR94]{BR94} I. Biswas and S. Ramanan, \emph{An infinitesimal study of the moduli of Hitchin pairs}, J. London Math. Soc. (2) \textbf{49} (1994), 219--231. \doi{10.1112/jlms/49.2.219}

\bibitem[BY96]{BY96} H. U. Boden and K. Yokogawa, \emph{Moduli spaces of parabolic Higgs bundles and parabolic $K(D)$ pairs over smooth curves. I}, Internat. J. Math. \textbf{7} (1996), 573--598. \doi{10.1142/S0129167X96000311}

\bibitem[Bou68]{Bou68} N. Bourbaki, \emph{\'El\'ements de math\'ematique. Fasc. XXXIV. Groupes et alg\`ebres de Lie. Chapitre IV: Groupes de Coxeter et syst\`emes de Tits. Chapitre V: Groupes engendr\'es par des r\'eflexions. Chapitre VI: syst\`emes de racines}, Actualit\'es Scientifiques et Industrielles, No. 1337 (Hermann, Paris, 1968); reprint \doi{10.1007/978-3-540-34491-9}

\bibitem[Cha02]{Cha02} P.-H. Chaudouard, \emph{La formule des traces pour les alg\`ebres de Lie}, Math. Ann. \textbf{322} (2002), 347--382. \doi{10.1007/s002080100274}

\bibitem[CL08]{CL08} P.-H. Chaudouard and G. Laumon, \emph{Sur l'homologie des fibres de Springer affines tronqu\'ees}, Duke Math. J. \textbf{145} (2008), 443--535. \doi{10.1215/00127094-2008-057}

\bibitem[CS78]{CS78} J.-L. Colliot-Th\'el\`ene and J.-J. Sansuc, \emph{Cohomologie des groupes de type multiplicatif sur les sch\'emas r\'eguliers}, C. R. Acad. Sci. Paris S\'er. A--B \textbf{287} (1978), A449--A452. \urlref{https://www.imo.universite-paris-saclay.fr/~jean-louis.colliot-thelene/78multiplicatif.pdf}

\bibitem[CS79]{CS79} J.-L. Colliot-Th\'el\`ene and J.-J. Sansuc, \emph{Fibr\'es quadratiques et composantes connexes r\'eelles}, Math. Ann. \textbf{244} (1979), 105--134. \doi{10.1007/BF01420486}

\bibitem[DH98]{DH98} I. V. Dolgachev and Y. Hu, \emph{Variation of geometric invariant theory quotients}, Publ. Math. Inst. Hautes \'Etudes Sci. \textbf{87} (1998), 5--51 (with an appendix by Nicolas Ressayre). \doi{10.1007/BF02698859}

\bibitem[DS95]{DS95} V. G. Drinfel'd and C. Simpson, \emph{$B$-structures on $G$-bundles and local triviality}, Math. Res. Lett. \textbf{2} (1995), 823--829. \doi{10.4310/MRL.1995.v2.n6.a13}

\bibitem[Est01]{Est01} E. Esteves, \emph{Compactifying the relative Jacobian over families of reduced curves}, Trans. Amer. Math. Soc. \textbf{353} (2001), 3045--3095. \doi{10.1090/S0002-9947-01-02746-5}

\bibitem[Fal93]{Fal93} G. Faltings, \emph{Stable $G$-bundles and projective connections}, J. Algebraic Geom. \textbf{2} (1993), 507--568.

\bibitem[GKM04]{GKM04} M. Goresky, R. Kottwitz and R. Macpherson, \emph{Homology of affine Springer fibers in the unramified case}, Duke Math. J. \textbf{121} (2004), 509--561. \doi{10.1215/S0012-7094-04-12135-9}

\bibitem[Gro61]{Gro61} A. Grothendieck, \emph{\'El\'ements de g\'eom\'etrie alg\'ebrique. III. \'Etude cohomologique des faisceaux coh\'erents. I}, Inst. Hautes \'Etudes Sci. Publ. Math. \textbf{11} (1961), 167. \doi{10.1007/BF02684273}

\bibitem[Gro67]{Gro67} A. Grothendieck, \emph{\'El\'ements de g\'eom\'etrie alg\'ebrique. IV. \'Etude locale des sch\'emas et des morphismes de sch\'emas IV}, Inst. Hautes \'Etudes Sci. Publ. Math. \textbf{32} (1967), 361. \doi{10.1007/BF02732123}

\bibitem[Har69]{Har69} G. Harder, \emph{Minkowskische Reduktionstheorie \"uber Funktionenk\"orpern}, Invent. Math. \textbf{7} (1969), 33--54. \doi{10.1007/BF01418773}

\bibitem[Har74]{Har74} G. Harder, \emph{Chevalley groups over function fields and automorphic forms}, Ann. of Math. (2) \textbf{100} (1974), 249--306. \doi{10.2307/1971073}

\bibitem[Hei08]{Hei08} J. Heinloth, \emph{Semistable reduction for $G$-bundles on curves}, J. Algebraic Geom. \textbf{17} (2008), 167--183. \doi{10.1090/S1056-3911-07-00476-6}

\bibitem[HS]{HS} J. Heinloth and A. H. W. Schmitt, \emph{The cohomology rings of moduli stacks of principal bundles over curves}, Doc. Math. \textbf{15} (2010), 423--488. \doi{10.4171/DM/302}

\bibitem[KL88]{KL88} D. Kazhdan and G. Lusztig, \emph{Fixed point varieties on affine flag manifolds}, Israel J. Math. \textbf{62} (1988), 129--168. \doi{10.1007/BF02787119}

\bibitem[Kos63]{Kos63} B. Kostant, \emph{Lie group representations on polynomial rings}, Amer. J. Math. \textbf{85} (1963), 327--404. \doi{10.2307/2373130}

\bibitem[Kot99]{Kot99} R. E. Kottwitz, \emph{Transfer factors for Lie algebras}, Represent. Theory \textbf{3} (1999), 127--138 (electronic). \doi{10.1090/S1088-4165-99-00068-0}

\bibitem[Lan75]{Lan75} S. G. Langton, \emph{Valuative criteria for families of vector bundles on algebraic varieties}, Ann. of Math. (2) \textbf{101} (1975), 88--110. \doi{10.2307/1970987}

\bibitem[Ngo08]{Ngo08} B. C. Ng\^o, \emph{Le lemme fondamental pour les alg\`ebres de Lie}, Preprint (2008), \urlref{https://arxiv.org/abs/0801.0446}; published in Publ. Math. Inst. Hautes \'Etudes Sci. \textbf{111} (2010), 1--169, \doi{10.1007/s10240-010-0026-7}

\bibitem[Ngo06]{Ngo06} B. C. Ng\^o, \emph{Fibration de Hitchin et endoscopie}, Invent. Math. \textbf{164} (2006), 399--453. \doi{10.1007/s00222-005-0483-7}

\bibitem[Nit91]{Nit91} N. Nitsure, \emph{Moduli space of semistable pairs on a curve}, Proc. London Math. Soc. (3) \textbf{62} (1991), 275--300. \doi{10.1112/plms/s3-62.2.275}

\bibitem[Ser94]{Ser94} J.-P. Serre, \emph{Cohomologie galoisienne}, Lecture Notes in Mathematics, vol. 5, fifth edition (Springer, Berlin, 1994). \doi{10.1007/BFb0108758}

\bibitem[Spr79]{Spr79} T. A. Springer, \emph{Reductive groups}, in \emph{Automorphic forms, representations and $L$-functions} (Proc. Sympos. Pure Math., Oregon State Univ., Corvallis, Ore., 1977), Part 1, Proceedings of Symposia in Pure Mathematics, vol. XXXIII (American Mathematical Society, Providence, RI, 1979), 3--27. \doi{10.1090/pspum/033.1}

\bibitem[Tha96]{Tha96} M. Thaddeus, \emph{Geometric invariant theory and flips}, J. Amer. Math. Soc. \textbf{9} (1996), 691--723. \doi{10.1090/S0894-0347-96-00204-4}

\end{thebibliography}

\bigskip
\begin{flushleft}
Pierre-Henri Chaudouard \ \ \href{mailto:Pierre-Henri.Chaudouard@math.u-psud.fr}{\texttt{Pierre-Henri.Chaudouard@math.u-psud.fr}}\\
CNRS et Universit\'{e} Paris-Sud, UMR 8628, Math\'{e}matique, B\^{a}timent 425,\\
F-91405 Orsay Cedex, France
\bigskip

G\'{e}rard Laumon \ \ \href{mailto:Gerard.Laumon@math.u-psud.fr}{\texttt{Gerard.Laumon@math.u-psud.fr}}\\
CNRS et Universit\'{e} Paris-Sud, UMR 8628, Math\'{e}matique, B\^{a}timent 425,\\
F-91405 Orsay Cedex, France
\end{flushleft}
\end{document}
